\documentclass[11pt]{book}
\usepackage[T1]{fontenc}
\usepackage{lmodern}
\usepackage[margin=0.8in]{geometry}
\usepackage{microtype}
\usepackage{amsmath,amssymb,mathtools}
\usepackage{booktabs,longtable,tabularx,array}
\usepackage{enumitem}
\usepackage[most]{tcolorbox}
\usepackage{xcolor}
\usepackage{hyperref}
\usepackage[nameinlink,noabbrev]{cleveref}
\usepackage{fancyhdr}
\usepackage{titlesec}
\usepackage{needspace}
\definecolor{navy}{HTML}{17365D}
\definecolor{blue}{HTML}{1F4E79}
\definecolor{teal}{HTML}{0F6B78}
\definecolor{lightblue}{HTML}{EAF2F8}
\definecolor{lightgray}{HTML}{F5F7F8}
\definecolor{green}{HTML}{2E7D32}
\definecolor{orange}{HTML}{A85600}
\hypersetup{colorlinks=true,linkcolor=navy,urlcolor=blue,citecolor=teal}
\pagestyle{fancy}\fancyhf{}\lhead{Calculus Applications Companion Workbook}\rhead{}\cfoot{\thepage}\setlength{\headheight}{16pt}
\titleformat{\chapter}{\LARGE\bfseries\color{navy}\raggedright}{\thechapter}{0.65em}{}
\titleformat{\section}{\Large\bfseries\color{blue}}{\thesection}{0.6em}{}
\titleformat{\subsection}{\large\bfseries\color{teal}}{\thesubsection}{0.5em}{}
\setlist[itemize]{leftmargin=1.4em,itemsep=0.2em,topsep=0.3em}
\setlist[enumerate]{leftmargin=1.6em,itemsep=0.25em,topsep=0.3em}
\setlength{\emergencystretch}{3em}
\raggedbottom
\newtcolorbox{strategybox}[1][]{breakable,colback=lightblue,colframe=blue,title=Problem-Solving Strategy,#1}
\newtcolorbox{problembox}[1][]{breakable,colback=white,colframe=navy,title=Problem,#1}
\newtcolorbox{promptbox}[1][]{breakable,colback=lightgray,colframe=teal,title=Before You Calculate,#1}
\newtcolorbox{errorbox}[1][]{breakable,colback=orange!6,colframe=orange,title=Common Modeling Errors,#1}
\newtcolorbox{solutionbox}[1][]{breakable,colback=green!4,colframe=green,title=Fully Worked Solution,#1}
\newtcolorbox{reflectionbox}[1][]{breakable,colback=teal!5,colframe=teal,title=Reflection,#1}
\newcommand{\pid}[1]{\texttt{#1}}
\newcommand{\workspace}{\par\vspace{0.7in}\hrulefill\par\vspace{0.55in}\hrulefill\par}
\begin{document}


\frontmatter
\begin{titlepage}
\centering
\vspace*{1.1in}
{\Huge\bfseries\color{navy} Calculus Applications Companion Workbook\par}
\vspace{0.25in}
{\Large\color{blue} Original Word Problems and Mathematical Modeling Practice\par}
\vspace{0.45in}
{\large Student Workbook: from verbal situation to model, calculus, validation, and interpretation\par}
\vfill
\begin{minipage}{0.82\textwidth}\small
This workbook is independently authored. Application taxonomy and calculus prerequisites are
cross-referenced to the supplied calculus source; the problem-solving pedagogy is an original synthesis
of \emph{How to Solve Word Problems in Calculus: A Solved Problem Approach}; optional computational
extensions are cross-referenced to \emph{Numerical Recipes: The Art of Scientific Computing}, Third Edition.
No source exercise is reproduced.
\end{minipage}
\vfill
{\large September 2026\par}
\end{titlepage}
\tableofcontents
\mainmatter
\chapter{How to Solve Calculus Word Problems}
\begin{strategybox}
A calculus word problem has two major phases: \textbf{build the mathematical model} and then
\textbf{solve the model with the appropriate mathematical tools}. The discipline in this workbook is to
make the modeling phase visible rather than jumping directly to symbolic manipulation.
\end{strategybox}
\section{Nine-step workflow}
\begin{enumerate}
\item \textbf{Read for meaning.} Identify the situation, given data, unknown quantity, units, and restrictions.
\item \textbf{Draw or organize the situation.} Use a diagram, coordinate system, table, timeline, force diagram, or compartment sketch when it clarifies relationships.
\item \textbf{Define variables.} State every variable and its units. Distinguish independent variables, dependent variables, parameters, and constants.
\item \textbf{Translate words into mathematical relationships.} Write equations, inequalities, physical laws, constraints, initial conditions, or boundary conditions before doing calculus.
\item \textbf{Reduce the model.} Use constraints to eliminate unnecessary variables when possible.
\item \textbf{Identify the calculus technique.} Explain why a derivative, integral, ODE, optimization method, vector operation, or numerical method is appropriate.
\item \textbf{Execute the mathematics.} Carry out the calculus and algebra, keeping exact quantities as long as useful.
\item \textbf{Validate.} Check units, domain, sign, magnitude, feasibility, endpoints, initial conditions, and residuals.
\item \textbf{Interpret.} End with a sentence that answers the original applied question in context.
\end{enumerate}
\section{How to use the workbook}
First attempt each problem without looking at a worked solution. If stuck, identify the exact modeling step that failed, review only that step, then restart the problem. Mixed reviews intentionally omit the technique label so that method selection becomes part of the task.
\section{Difficulty levels}
\begin{description}
\item[Foundation (F)] The model is nearly explicit; emphasis is on recognizing the calculus idea.
\item[Standard (S)] The learner must build the model and carry out several linked steps.
\item[Advanced (A)] The problem combines nontrivial modeling, several mathematical steps, or a computational extension.
\end{description}

\part{Modeling Foundations}
\chapter{Population, cell-phone subscribers, GDP, apartment rental and other baseline trend models}\label{app:A001}
\textbf{Application ID:} \pid{A001}\quad \textbf{Source locus:} P.1–P.3\quad \textbf{Source pages:} 9, 12, 16, 18
\section*{Application Context}
This family focuses on translating observations and verbal relationships into functions that can be evaluated, compared, and differentiated. The canonical application family is \emph{Population, cell-phone subscribers, GDP, apartment rental and other baseline trend models}.
\section*{Calculus Toolkit}
\begin{tabularx}{\textwidth}{@{}lX@{}}\toprule
Primary technique & Functions, graphs, linear models and rates \\
Prerequisites & algebra, functions, units, graphs \\
Numerical extension & Ch. 15 (modeling of data); Ch. 3 when interpolation is used \\
Analytical/numerical class & NUMERICAL-OPTIONAL \\
\bottomrule\end{tabularx}
\section*{Modeling Questions to Ask}\begin{itemize}
\item Which quantity is the independent variable?
\item What output is being modeled?
\item What units belong to the model parameters?
\item Over what domain is the model credible?
\end{itemize}
\begin{errorbox}Do not start differentiating before the variables, units, and functional relationship are explicit. Check the domain before extrapolating a fitted model.\end{errorbox}
\section*{Problems}
\needspace{8\baselineskip}
\subsection*{\pid{APP-001-F-01} — Foundation}
\begin{problembox}
A monitoring team studying population, cell-phone subscribers, gdp, apartment rental and other baseline trend models records a value of 22 units at time 0 and 38 units at time 4. Assume a linear model over the next few time units. Construct the model and predict the value at time 7.
\end{problembox}
\begin{promptbox}\begin{enumerate}
\item Which quantity is the independent variable?
\item What output is being modeled?
\item What units belong to the model parameters?
\item Over what domain is the model credible?
\end{enumerate}\end{promptbox}
\textbf{Workspace}\workspace
\needspace{8\baselineskip}
\subsection*{\pid{APP-001-S-01} — Standard}
\begin{problembox}
For a short planning horizon, a quantity associated with population, cell-phone subscribers, gdp, apartment rental and other baseline trend models is modeled by $Q(t)=53e^{0.045t}$, where $t$ is measured in years. Find the instantaneous growth rate at $t=5$ and interpret it.
\end{problembox}
\begin{promptbox}\begin{enumerate}
\item Which quantity is the independent variable?
\item What output is being modeled?
\item What units belong to the model parameters?
\item Over what domain is the model credible?
\end{enumerate}\end{promptbox}
\textbf{Workspace}\workspace
\needspace{8\baselineskip}
\subsection*{\pid{APP-001-A-01} — Advanced}
\begin{problembox}
Measurements related to population, cell-phone subscribers, gdp, apartment rental and other baseline trend models are 32 units initially and 47.4 units after 6 time units. Assume $Q(t)=Q_0e^{rt}$. Estimate $r$, then determine when the model first reaches 57.6 units.
\end{problembox}
\begin{promptbox}\begin{enumerate}
\item Which quantity is the independent variable?
\item What output is being modeled?
\item What units belong to the model parameters?
\item Over what domain is the model credible?
\end{enumerate}\end{promptbox}
\textbf{Workspace}\workspace
\begin{reflectionbox}\begin{enumerate}
\item Which sentence or datum in these problems supplied the main mathematical constraint?
\item Which assumptions were essential to the model, and how would the model change if one were removed?
\item When would an analytical solution become inconvenient enough to justify a numerical method?
\end{enumerate}\end{reflectionbox}
\chapter{Carbon-dioxide concentration and environmental trend models}\label{app:A002}
\textbf{Application ID:} \pid{A002}\quad \textbf{Source locus:} P.1\quad \textbf{Source pages:} 7
\section*{Application Context}
This family focuses on translating observations and verbal relationships into functions that can be evaluated, compared, and differentiated. The canonical application family is \emph{Carbon-dioxide concentration and environmental trend models}.
\section*{Calculus Toolkit}
\begin{tabularx}{\textwidth}{@{}lX@{}}\toprule
Primary technique & Function modeling from data \\
Prerequisites & algebra, functions, units, graphs \\
Numerical extension & Ch. 15 (modeling of data); Ch. 3 when interpolation is used \\
Analytical/numerical class & NUMERICAL-NATURAL \\
\bottomrule\end{tabularx}
\section*{Modeling Questions to Ask}\begin{itemize}
\item Which quantity is the independent variable?
\item What output is being modeled?
\item What units belong to the model parameters?
\item Over what domain is the model credible?
\end{itemize}
\begin{errorbox}Do not start differentiating before the variables, units, and functional relationship are explicit. Check the domain before extrapolating a fitted model.\end{errorbox}
\section*{Problems}
\needspace{8\baselineskip}
\subsection*{\pid{APP-002-F-01} — Foundation}
\begin{problembox}
A monitoring team studying carbon-dioxide concentration and environmental trend models records a value of 24 units at time 0 and 44 units at time 4. Assume a linear model over the next few time units. Construct the model and predict the value at time 7.
\end{problembox}
\begin{promptbox}\begin{enumerate}
\item Which quantity is the independent variable?
\item What output is being modeled?
\item What units belong to the model parameters?
\item Over what domain is the model credible?
\end{enumerate}\end{promptbox}
\textbf{Workspace}\workspace
\needspace{8\baselineskip}
\subsection*{\pid{APP-002-S-01} — Standard}
\begin{problembox}
For a short planning horizon, a quantity associated with carbon-dioxide concentration and environmental trend models is modeled by $Q(t)=56e^{0.050t}$, where $t$ is measured in years. Find the instantaneous growth rate at $t=5$ and interpret it.
\end{problembox}
\begin{promptbox}\begin{enumerate}
\item Which quantity is the independent variable?
\item What output is being modeled?
\item What units belong to the model parameters?
\item Over what domain is the model credible?
\end{enumerate}\end{promptbox}
\textbf{Workspace}\workspace
\needspace{8\baselineskip}
\subsection*{\pid{APP-002-A-01} — Advanced}
\begin{problembox}
Measurements related to carbon-dioxide concentration and environmental trend models are 34 units initially and 51.3 units after 6 time units. Assume $Q(t)=Q_0e^{rt}$. Estimate $r$, then determine when the model first reaches 61.2 units.
\end{problembox}
\begin{promptbox}\begin{enumerate}
\item Which quantity is the independent variable?
\item What output is being modeled?
\item What units belong to the model parameters?
\item Over what domain is the model credible?
\end{enumerate}\end{promptbox}
\textbf{Workspace}\workspace
\begin{reflectionbox}\begin{enumerate}
\item Which sentence or datum in these problems supplied the main mathematical constraint?
\item Which assumptions were essential to the model, and how would the model change if one were removed?
\item When would an analytical solution become inconvenient enough to justify a numerical method?
\end{enumerate}\end{reflectionbox}
\chapter{Biodiesel production, sales, salary and simple production models}\label{app:A003}
\textbf{Application ID:} \pid{A003}\quad \textbf{Source locus:} P.1–P.3; 9\quad \textbf{Source pages:} 17, 40, 607
\section*{Application Context}
This family focuses on translating observations and verbal relationships into functions that can be evaluated, compared, and differentiated. The canonical application family is \emph{Biodiesel production, sales, salary and simple production models}.
\section*{Calculus Toolkit}
\begin{tabularx}{\textwidth}{@{}lX@{}}\toprule
Primary technique & Function modeling; sequences \\
Prerequisites & algebra, functions, units, graphs \\
Numerical extension & Ch. 15 (modeling of data); Ch. 3 when interpolation is used \\
Analytical/numerical class & NUMERICAL-OPTIONAL \\
\bottomrule\end{tabularx}
\section*{Modeling Questions to Ask}\begin{itemize}
\item Which quantity is the independent variable?
\item What output is being modeled?
\item What units belong to the model parameters?
\item Over what domain is the model credible?
\end{itemize}
\begin{errorbox}Do not start differentiating before the variables, units, and functional relationship are explicit. Check the domain before extrapolating a fitted model.\end{errorbox}
\section*{Problems}
\needspace{8\baselineskip}
\subsection*{\pid{APP-003-F-01} — Foundation}
\begin{problembox}
A monitoring team studying biodiesel production, sales, salary and simple production models records a value of 26 units at time 0 and 50 units at time 4. Assume a linear model over the next few time units. Construct the model and predict the value at time 7.
\end{problembox}
\begin{promptbox}\begin{enumerate}
\item Which quantity is the independent variable?
\item What output is being modeled?
\item What units belong to the model parameters?
\item Over what domain is the model credible?
\end{enumerate}\end{promptbox}
\textbf{Workspace}\workspace
\needspace{8\baselineskip}
\subsection*{\pid{APP-003-S-01} — Standard}
\begin{problembox}
For a short planning horizon, a quantity associated with biodiesel production, sales, salary and simple production models is modeled by $Q(t)=59e^{0.055t}$, where $t$ is measured in years. Find the instantaneous growth rate at $t=5$ and interpret it.
\end{problembox}
\begin{promptbox}\begin{enumerate}
\item Which quantity is the independent variable?
\item What output is being modeled?
\item What units belong to the model parameters?
\item Over what domain is the model credible?
\end{enumerate}\end{promptbox}
\textbf{Workspace}\workspace
\needspace{8\baselineskip}
\subsection*{\pid{APP-003-A-01} — Advanced}
\begin{problembox}
Measurements related to biodiesel production, sales, salary and simple production models are 36 units initially and 52.2 units after 6 time units. Assume $Q(t)=Q_0e^{rt}$. Estimate $r$, then determine when the model first reaches 64.8 units.
\end{problembox}
\begin{promptbox}\begin{enumerate}
\item Which quantity is the independent variable?
\item What output is being modeled?
\item What units belong to the model parameters?
\item Over what domain is the model credible?
\end{enumerate}\end{promptbox}
\textbf{Workspace}\workspace
\begin{reflectionbox}\begin{enumerate}
\item Which sentence or datum in these problems supplied the main mathematical constraint?
\item Which assumptions were essential to the model, and how would the model change if one were removed?
\item When would an analytical solution become inconvenient enough to justify a numerical method?
\end{enumerate}\end{reflectionbox}
\chapter{Ferris-wheel height, angular position and periodic geometry}\label{app:A004}
\textbf{Application ID:} \pid{A004}\quad \textbf{Source locus:} P.4; 12\quad \textbf{Source pages:} 40, 870
\section*{Application Context}
This family focuses on translating observations and verbal relationships into functions that can be evaluated, compared, and differentiated. The canonical application family is \emph{Ferris-wheel height, angular position and periodic geometry}.
\section*{Calculus Toolkit}
\begin{tabularx}{\textwidth}{@{}lX@{}}\toprule
Primary technique & Trigonometric functions and parametric motion \\
Prerequisites & algebra, functions, units, graphs \\
Numerical extension & Ch. 15 (modeling of data); Ch. 3 when interpolation is used \\
Analytical/numerical class & ANALYTICAL \\
\bottomrule\end{tabularx}
\section*{Modeling Questions to Ask}\begin{itemize}
\item Which quantity is the independent variable?
\item What output is being modeled?
\item What units belong to the model parameters?
\item Over what domain is the model credible?
\end{itemize}
\begin{errorbox}Do not start differentiating before the variables, units, and functional relationship are explicit. Check the domain before extrapolating a fitted model.\end{errorbox}
\section*{Problems}
\needspace{8\baselineskip}
\subsection*{\pid{APP-004-F-01} — Foundation}
\begin{problembox}
A monitoring team studying ferris-wheel height, angular position and periodic geometry records a value of 28 units at time 0 and 40 units at time 4. Assume a linear model over the next few time units. Construct the model and predict the value at time 7.
\end{problembox}
\begin{promptbox}\begin{enumerate}
\item Which quantity is the independent variable?
\item What output is being modeled?
\item What units belong to the model parameters?
\item Over what domain is the model credible?
\end{enumerate}\end{promptbox}
\textbf{Workspace}\workspace
\needspace{8\baselineskip}
\subsection*{\pid{APP-004-S-01} — Standard}
\begin{problembox}
For a short planning horizon, a quantity associated with ferris-wheel height, angular position and periodic geometry is modeled by $Q(t)=62e^{0.040t}$, where $t$ is measured in years. Find the instantaneous growth rate at $t=5$ and interpret it.
\end{problembox}
\begin{promptbox}\begin{enumerate}
\item Which quantity is the independent variable?
\item What output is being modeled?
\item What units belong to the model parameters?
\item Over what domain is the model credible?
\end{enumerate}\end{promptbox}
\textbf{Workspace}\workspace
\needspace{8\baselineskip}
\subsection*{\pid{APP-004-A-01} — Advanced}
\begin{problembox}
Measurements related to ferris-wheel height, angular position and periodic geometry are 38 units initially and 56.2 units after 6 time units. Assume $Q(t)=Q_0e^{rt}$. Estimate $r$, then determine when the model first reaches 68.4 units.
\end{problembox}
\begin{promptbox}\begin{enumerate}
\item Which quantity is the independent variable?
\item What output is being modeled?
\item What units belong to the model parameters?
\item Over what domain is the model credible?
\end{enumerate}\end{promptbox}
\textbf{Workspace}\workspace
\begin{reflectionbox}\begin{enumerate}
\item Which sentence or datum in these problems supplied the main mathematical constraint?
\item Which assumptions were essential to the model, and how would the model change if one were removed?
\item When would an analytical solution become inconvenient enough to justify a numerical method?
\end{enumerate}\end{reflectionbox}
\chapter{Automobile aerodynamics, thermostat and ripple models}\label{app:A005}
\textbf{Application ID:} \pid{A005}\quad \textbf{Source locus:} P.3; 2.6\quad \textbf{Source pages:} 29–30, 153
\section*{Application Context}
This family focuses on translating observations and verbal relationships into functions that can be evaluated, compared, and differentiated. The canonical application family is \emph{Automobile aerodynamics, thermostat and ripple models}.
\section*{Calculus Toolkit}
\begin{tabularx}{\textwidth}{@{}lX@{}}\toprule
Primary technique & Functional dependence and geometry \\
Prerequisites & algebra, functions, units, graphs \\
Numerical extension & Ch. 15 (modeling of data); Ch. 3 when interpolation is used \\
Analytical/numerical class & NUMERICAL-OPTIONAL \\
\bottomrule\end{tabularx}
\section*{Modeling Questions to Ask}\begin{itemize}
\item Which quantity is the independent variable?
\item What output is being modeled?
\item What units belong to the model parameters?
\item Over what domain is the model credible?
\end{itemize}
\begin{errorbox}Do not start differentiating before the variables, units, and functional relationship are explicit. Check the domain before extrapolating a fitted model.\end{errorbox}
\section*{Problems}
\needspace{8\baselineskip}
\subsection*{\pid{APP-005-F-01} — Foundation}
\begin{problembox}
A monitoring team studying automobile aerodynamics, thermostat and ripple models records a value of 30 units at time 0 and 46 units at time 4. Assume a linear model over the next few time units. Construct the model and predict the value at time 7.
\end{problembox}
\begin{promptbox}\begin{enumerate}
\item Which quantity is the independent variable?
\item What output is being modeled?
\item What units belong to the model parameters?
\item Over what domain is the model credible?
\end{enumerate}\end{promptbox}
\textbf{Workspace}\workspace
\needspace{8\baselineskip}
\subsection*{\pid{APP-005-S-01} — Standard}
\begin{problembox}
For a short planning horizon, a quantity associated with automobile aerodynamics, thermostat and ripple models is modeled by $Q(t)=65e^{0.045t}$, where $t$ is measured in years. Find the instantaneous growth rate at $t=5$ and interpret it.
\end{problembox}
\begin{promptbox}\begin{enumerate}
\item Which quantity is the independent variable?
\item What output is being modeled?
\item What units belong to the model parameters?
\item Over what domain is the model credible?
\end{enumerate}\end{promptbox}
\textbf{Workspace}\workspace
\needspace{8\baselineskip}
\subsection*{\pid{APP-005-A-01} — Advanced}
\begin{problembox}
Measurements related to automobile aerodynamics, thermostat and ripple models are 40 units initially and 60.4 units after 6 time units. Assume $Q(t)=Q_0e^{rt}$. Estimate $r$, then determine when the model first reaches 72.0 units.
\end{problembox}
\begin{promptbox}\begin{enumerate}
\item Which quantity is the independent variable?
\item What output is being modeled?
\item What units belong to the model parameters?
\item Over what domain is the model credible?
\end{enumerate}\end{promptbox}
\textbf{Workspace}\workspace
\begin{reflectionbox}\begin{enumerate}
\item Which sentence or datum in these problems supplied the main mathematical constraint?
\item Which assumptions were essential to the model, and how would the model change if one were removed?
\item When would an analytical solution become inconvenient enough to justify a numerical method?
\end{enumerate}\end{reflectionbox}
\part{Rates and Differential Calculus}
\chapter{Average velocity, speed and acceleration}\label{app:A006}
\textbf{Application ID:} \pid{A006}\quad \textbf{Source locus:} 1; 2.1–2.3; 12.3\quad \textbf{Source pages:} 44, 93, 116, 128, 132, 160, 180, 257, 910
\section*{Application Context}
This family treats change quantitatively. The derivative connects a mathematical model to an instantaneous physical, biological, or economic rate. The canonical application family is \emph{Average velocity, speed and acceleration}.
\section*{Calculus Toolkit}
\begin{tabularx}{\textwidth}{@{}lX@{}}\toprule
Primary technique & Difference quotients, derivatives and higher derivatives \\
Prerequisites & functions, limits, derivative rules \\
Numerical extension & §5.7 (numerical derivatives) when derivative data are sampled \\
Analytical/numerical class & NUMERICAL-OPTIONAL \\
\bottomrule\end{tabularx}
\section*{Modeling Questions to Ask}\begin{itemize}
\item What quantity is changing?
\item With respect to which independent variable?
\item Is the requested rate average or instantaneous?
\item What units must the derivative have?
\end{itemize}
\begin{errorbox}Distinguish an average rate over an interval from an instantaneous derivative. Carry units through differentiation.\end{errorbox}
\section*{Problems}
\needspace{8\baselineskip}
\subsection*{\pid{APP-006-F-01} — Foundation}
\begin{problembox}
A quantity in a average velocity, speed and acceleration model is $Q(t)=5t^2+9t+10$. Find the instantaneous rate of change at $t=3$ and state its units as ‘quantity units per time unit.’
\end{problembox}
\begin{promptbox}\begin{enumerate}
\item What quantity is changing?
\item With respect to which independent variable?
\item Is the requested rate average or instantaneous?
\item What units must the derivative have?
\end{enumerate}\end{promptbox}
\textbf{Workspace}\workspace
\needspace{8\baselineskip}
\subsection*{\pid{APP-006-S-01} — Standard}
\begin{problembox}
For an application involving average velocity, speed and acceleration, let the total cost-like quantity be $C(q)=560+10q+0.030q^2$. At $q=46$, compute the marginal value $C'(q)$ and compare it with the average value $C(q)/q$.
\end{problembox}
\begin{promptbox}\begin{enumerate}
\item What quantity is changing?
\item With respect to which independent variable?
\item Is the requested rate average or instantaneous?
\item What units must the derivative have?
\end{enumerate}\end{promptbox}
\textbf{Workspace}\workspace
\needspace{8\baselineskip}
\subsection*{\pid{APP-006-A-01} — Advanced}
\begin{problembox}
A sensitivity model for average velocity, speed and acceleration contains the response $R(x)=\sqrt{x}$. Around $x=106$, the input may increase by 1.4\%. Use differentials to estimate the change in $R$, and compare with the exact change.
\end{problembox}
\begin{promptbox}\begin{enumerate}
\item What quantity is changing?
\item With respect to which independent variable?
\item Is the requested rate average or instantaneous?
\item What units must the derivative have?
\end{enumerate}\end{promptbox}
\textbf{Workspace}\workspace
\begin{reflectionbox}\begin{enumerate}
\item Which sentence or datum in these problems supplied the main mathematical constraint?
\item Which assumptions were essential to the model, and how would the model change if one were removed?
\item When would an analytical solution become inconvenient enough to justify a numerical method?
\end{enumerate}\end{reflectionbox}
\chapter{Particle, rectilinear and vertical motion}\label{app:A007}
\textbf{Application ID:} \pid{A007}\quad \textbf{Source locus:} 2; 4; 12\quad \textbf{Source pages:} 120, 132, 161, 178, 179, 254, 257, 291–295, 835–865
\section*{Application Context}
This family treats change quantitatively. The derivative connects a mathematical model to an instantaneous physical, biological, or economic rate. The canonical application family is \emph{Particle, rectilinear and vertical motion}.
\section*{Calculus Toolkit}
\begin{tabularx}{\textwidth}{@{}lX@{}}\toprule
Primary technique & Position, velocity and acceleration \\
Prerequisites & functions, limits, derivative rules \\
Numerical extension & §5.7 (numerical derivatives) when derivative data are sampled \\
Analytical/numerical class & NUMERICAL-NATURAL \\
\bottomrule\end{tabularx}
\section*{Modeling Questions to Ask}\begin{itemize}
\item What quantity is changing?
\item With respect to which independent variable?
\item Is the requested rate average or instantaneous?
\item What units must the derivative have?
\end{itemize}
\begin{errorbox}Distinguish an average rate over an interval from an instantaneous derivative. Carry units through differentiation.\end{errorbox}
\section*{Problems}
\needspace{8\baselineskip}
\subsection*{\pid{APP-007-F-01} — Foundation}
\begin{problembox}
A quantity in a particle, rectilinear and vertical motion model is $Q(t)=6t^2+10t+10$. Find the instantaneous rate of change at $t=3$ and state its units as ‘quantity units per time unit.’
\end{problembox}
\begin{promptbox}\begin{enumerate}
\item What quantity is changing?
\item With respect to which independent variable?
\item Is the requested rate average or instantaneous?
\item What units must the derivative have?
\end{enumerate}\end{promptbox}
\textbf{Workspace}\workspace
\needspace{8\baselineskip}
\subsection*{\pid{APP-007-S-01} — Standard}
\begin{problembox}
For an application involving particle, rectilinear and vertical motion, let the total cost-like quantity be $C(q)=570+11q+0.035q^2$. At $q=47$, compute the marginal value $C'(q)$ and compare it with the average value $C(q)/q$.
\end{problembox}
\begin{promptbox}\begin{enumerate}
\item What quantity is changing?
\item With respect to which independent variable?
\item Is the requested rate average or instantaneous?
\item What units must the derivative have?
\end{enumerate}\end{promptbox}
\textbf{Workspace}\workspace
\needspace{8\baselineskip}
\subsection*{\pid{APP-007-A-01} — Advanced}
\begin{problembox}
A sensitivity model for particle, rectilinear and vertical motion contains the response $R(x)=\sqrt{x}$. Around $x=107$, the input may increase by 1.5\%. Use differentials to estimate the change in $R$, and compare with the exact change.
\end{problembox}
\begin{promptbox}\begin{enumerate}
\item What quantity is changing?
\item With respect to which independent variable?
\item Is the requested rate average or instantaneous?
\item What units must the derivative have?
\end{enumerate}\end{promptbox}
\textbf{Workspace}\workspace
\begin{reflectionbox}\begin{enumerate}
\item Which sentence or datum in these problems supplied the main mathematical constraint?
\item Which assumptions were essential to the model, and how would the model change if one were removed?
\item When would an analytical solution become inconvenient enough to justify a numerical method?
\end{enumerate}\end{reflectionbox}
\chapter{Marginal cost, revenue, productivity and utility}\label{app:A011}
\textbf{Application ID:} \pid{A011}\quad \textbf{Source locus:} 13.3; 13.8–13.10\quad \textbf{Source pages:} 902–903, 949, 953, 959, 964
\section*{Application Context}
This family interprets derivatives as marginal quantities—local changes in cost, revenue, productivity, utility, or related response measures. The canonical application family is \emph{Marginal cost, revenue, productivity and utility}.
\section*{Calculus Toolkit}
\begin{tabularx}{\textwidth}{@{}lX@{}}\toprule
Primary technique & Partial derivatives and marginal analysis \\
Prerequisites & derivatives, function interpretation, units \\
Numerical extension & Ch. 10 (optimization) and Ch. 15 (model fitting) for computational extensions \\
Analytical/numerical class & NUMERICAL-NATURAL \\
\bottomrule\end{tabularx}
\section*{Modeling Questions to Ask}\begin{itemize}
\item What economic/production quantity is differentiated?
\item At what operating level is the marginal value required?
\item Is the requested quantity marginal or average?
\item What units should the derivative carry?
\end{itemize}
\begin{errorbox}A marginal quantity is a derivative evaluated at a production level; it is not the same as an average quantity.\end{errorbox}
\section*{Problems}
\needspace{8\baselineskip}
\subsection*{\pid{APP-011-F-01} — Foundation}
\begin{problembox}
A quantity in a marginal cost, revenue, productivity and utility model is $Q(t)=10t^2+9t+10$. Find the instantaneous rate of change at $t=3$ and state its units as ‘quantity units per time unit.’
\end{problembox}
\begin{promptbox}\begin{enumerate}
\item What economic/production quantity is differentiated?
\item At what operating level is the marginal value required?
\item Is the requested quantity marginal or average?
\item What units should the derivative carry?
\end{enumerate}\end{promptbox}
\textbf{Workspace}\workspace
\needspace{8\baselineskip}
\subsection*{\pid{APP-011-S-01} — Standard}
\begin{problembox}
For an application involving marginal cost, revenue, productivity and utility, let the total cost-like quantity be $C(q)=610+10q+0.040q^2$. At $q=51$, compute the marginal value $C'(q)$ and compare it with the average value $C(q)/q$.
\end{problembox}
\begin{promptbox}\begin{enumerate}
\item What economic/production quantity is differentiated?
\item At what operating level is the marginal value required?
\item Is the requested quantity marginal or average?
\item What units should the derivative carry?
\end{enumerate}\end{promptbox}
\textbf{Workspace}\workspace
\needspace{8\baselineskip}
\subsection*{\pid{APP-011-A-01} — Advanced}
\begin{problembox}
A sensitivity model for marginal cost, revenue, productivity and utility contains the response $R(x)=\sqrt{x}$. Around $x=111$, the input may increase by 1.5\%. Use differentials to estimate the change in $R$, and compare with the exact change.
\end{problembox}
\begin{promptbox}\begin{enumerate}
\item What economic/production quantity is differentiated?
\item At what operating level is the marginal value required?
\item Is the requested quantity marginal or average?
\item What units should the derivative carry?
\end{enumerate}\end{promptbox}
\textbf{Workspace}\workspace
\begin{reflectionbox}\begin{enumerate}
\item Which sentence or datum in these problems supplied the main mathematical constraint?
\item Which assumptions were essential to the model, and how would the model change if one were removed?
\item When would an analytical solution become inconvenient enough to justify a numerical method?
\end{enumerate}\end{reflectionbox}
\chapter{Sales growth, declining sales, depreciation and investment growth}\label{app:A012}
\textbf{Application ID:} \pid{A012}\quad \textbf{Source locus:} 2; 5; 6; 9\quad \textbf{Source pages:} 143, 198, 243, 307, 349, 359, 418, 436, 606, 680
\section*{Application Context}
This family treats change quantitatively. The derivative connects a mathematical model to an instantaneous physical, biological, or economic rate. The canonical application family is \emph{Sales growth, declining sales, depreciation and investment growth}.
\section*{Calculus Toolkit}
\begin{tabularx}{\textwidth}{@{}lX@{}}\toprule
Primary technique & Rates of change and exponential models \\
Prerequisites & functions, limits, derivative rules \\
Numerical extension & §5.7 (numerical derivatives) when derivative data are sampled \\
Analytical/numerical class & NUMERICAL-NATURAL \\
\bottomrule\end{tabularx}
\section*{Modeling Questions to Ask}\begin{itemize}
\item What quantity is changing?
\item With respect to which independent variable?
\item Is the requested rate average or instantaneous?
\item What units must the derivative have?
\end{itemize}
\begin{errorbox}Distinguish an average rate over an interval from an instantaneous derivative. Carry units through differentiation.\end{errorbox}
\section*{Problems}
\needspace{8\baselineskip}
\subsection*{\pid{APP-012-F-01} — Foundation}
\begin{problembox}
A quantity in a sales growth, declining sales, depreciation and investment growth model is $Q(t)=5t^2+10t+10$. Find the instantaneous rate of change at $t=3$ and state its units as ‘quantity units per time unit.’
\end{problembox}
\begin{promptbox}\begin{enumerate}
\item What quantity is changing?
\item With respect to which independent variable?
\item Is the requested rate average or instantaneous?
\item What units must the derivative have?
\end{enumerate}\end{promptbox}
\textbf{Workspace}\workspace
\needspace{8\baselineskip}
\subsection*{\pid{APP-012-S-01} — Standard}
\begin{problembox}
For an application involving sales growth, declining sales, depreciation and investment growth, let the total cost-like quantity be $C(q)=620+11q+0.030q^2$. At $q=52$, compute the marginal value $C'(q)$ and compare it with the average value $C(q)/q$.
\end{problembox}
\begin{promptbox}\begin{enumerate}
\item What quantity is changing?
\item With respect to which independent variable?
\item Is the requested rate average or instantaneous?
\item What units must the derivative have?
\end{enumerate}\end{promptbox}
\textbf{Workspace}\workspace
\needspace{8\baselineskip}
\subsection*{\pid{APP-012-A-01} — Advanced}
\begin{problembox}
A sensitivity model for sales growth, declining sales, depreciation and investment growth contains the response $R(x)=\sqrt{x}$. Around $x=112$, the input may increase by 1.2\%. Use differentials to estimate the change in $R$, and compare with the exact change.
\end{problembox}
\begin{promptbox}\begin{enumerate}
\item What quantity is changing?
\item With respect to which independent variable?
\item Is the requested rate average or instantaneous?
\item What units must the derivative have?
\end{enumerate}\end{promptbox}
\textbf{Workspace}\workspace
\begin{reflectionbox}\begin{enumerate}
\item Which sentence or datum in these problems supplied the main mathematical constraint?
\item Which assumptions were essential to the model, and how would the model change if one were removed?
\item When would an analytical solution become inconvenient enough to justify a numerical method?
\end{enumerate}\end{reflectionbox}
\chapter{Concentration, trachea contraction and physiological rate-of-change models}\label{app:A013}
\textbf{Application ID:} \pid{A013}\quad \textbf{Source locus:} 3.9; 6\quad \textbf{Source pages:} 189, 234, 442
\section*{Application Context}
This family treats change quantitatively. The derivative connects a mathematical model to an instantaneous physical, biological, or economic rate. The canonical application family is \emph{Concentration, trachea contraction and physiological rate-of-change models}.
\section*{Calculus Toolkit}
\begin{tabularx}{\textwidth}{@{}lX@{}}\toprule
Primary technique & Differentials, sensitivity and rate models \\
Prerequisites & functions, limits, derivative rules \\
Numerical extension & §5.7 (numerical derivatives) when derivative data are sampled \\
Analytical/numerical class & NUMERICAL-NATURAL \\
\bottomrule\end{tabularx}
\section*{Modeling Questions to Ask}\begin{itemize}
\item What quantity is changing?
\item With respect to which independent variable?
\item Is the requested rate average or instantaneous?
\item What units must the derivative have?
\end{itemize}
\begin{errorbox}Distinguish an average rate over an interval from an instantaneous derivative. Carry units through differentiation.\end{errorbox}
\section*{Problems}
\needspace{8\baselineskip}
\subsection*{\pid{APP-013-F-01} — Foundation}
\begin{problembox}
A quantity in a concentration, trachea contraction and physiological rate-of-change models model is $Q(t)=6t^2+11t+10$. Find the instantaneous rate of change at $t=3$ and state its units as ‘quantity units per time unit.’
\end{problembox}
\begin{promptbox}\begin{enumerate}
\item What quantity is changing?
\item With respect to which independent variable?
\item Is the requested rate average or instantaneous?
\item What units must the derivative have?
\end{enumerate}\end{promptbox}
\textbf{Workspace}\workspace
\needspace{8\baselineskip}
\subsection*{\pid{APP-013-S-01} — Standard}
\begin{problembox}
For an application involving concentration, trachea contraction and physiological rate-of-change models, let the total cost-like quantity be $C(q)=630+12q+0.035q^2$. At $q=53$, compute the marginal value $C'(q)$ and compare it with the average value $C(q)/q$.
\end{problembox}
\begin{promptbox}\begin{enumerate}
\item What quantity is changing?
\item With respect to which independent variable?
\item Is the requested rate average or instantaneous?
\item What units must the derivative have?
\end{enumerate}\end{promptbox}
\textbf{Workspace}\workspace
\needspace{8\baselineskip}
\subsection*{\pid{APP-013-A-01} — Advanced}
\begin{problembox}
A sensitivity model for concentration, trachea contraction and physiological rate-of-change models contains the response $R(x)=\sqrt{x}$. Around $x=113$, the input may increase by 1.3\%. Use differentials to estimate the change in $R$, and compare with the exact change.
\end{problembox}
\begin{promptbox}\begin{enumerate}
\item What quantity is changing?
\item With respect to which independent variable?
\item Is the requested rate average or instantaneous?
\item What units must the derivative have?
\end{enumerate}\end{promptbox}
\textbf{Workspace}\workspace
\begin{reflectionbox}\begin{enumerate}
\item Which sentence or datum in these problems supplied the main mathematical constraint?
\item Which assumptions were essential to the model, and how would the model change if one were removed?
\item When would an analytical solution become inconvenient enough to justify a numerical method?
\end{enumerate}\end{reflectionbox}
\chapter{Error and differential approximation}\label{app:A014}
\textbf{Application ID:} \pid{A014}\quad \textbf{Source locus:} 3.9; 13.4\quad \textbf{Source pages:} 237, 241, 907, 909
\section*{Application Context}
This family treats change quantitatively. The derivative connects a mathematical model to an instantaneous physical, biological, or economic rate. The canonical application family is \emph{Error and differential approximation}.
\section*{Calculus Toolkit}
\begin{tabularx}{\textwidth}{@{}lX@{}}\toprule
Primary technique & Differentials and propagated error \\
Prerequisites & functions, limits, derivative rules \\
Numerical extension & §5.7 (numerical derivatives) when derivative data are sampled \\
Analytical/numerical class & NUMERICAL-NATURAL \\
\bottomrule\end{tabularx}
\section*{Modeling Questions to Ask}\begin{itemize}
\item What quantity is changing?
\item With respect to which independent variable?
\item Is the requested rate average or instantaneous?
\item What units must the derivative have?
\end{itemize}
\begin{errorbox}Distinguish an average rate over an interval from an instantaneous derivative. Carry units through differentiation.\end{errorbox}
\section*{Problems}
\needspace{8\baselineskip}
\subsection*{\pid{APP-014-F-01} — Foundation}
\begin{problembox}
A quantity in a error and differential approximation model is $Q(t)=7t^2+12t+10$. Find the instantaneous rate of change at $t=3$ and state its units as ‘quantity units per time unit.’
\end{problembox}
\begin{promptbox}\begin{enumerate}
\item What quantity is changing?
\item With respect to which independent variable?
\item Is the requested rate average or instantaneous?
\item What units must the derivative have?
\end{enumerate}\end{promptbox}
\textbf{Workspace}\workspace
\needspace{8\baselineskip}
\subsection*{\pid{APP-014-S-01} — Standard}
\begin{problembox}
For an application involving error and differential approximation, let the total cost-like quantity be $C(q)=640+13q+0.040q^2$. At $q=54$, compute the marginal value $C'(q)$ and compare it with the average value $C(q)/q$.
\end{problembox}
\begin{promptbox}\begin{enumerate}
\item What quantity is changing?
\item With respect to which independent variable?
\item Is the requested rate average or instantaneous?
\item What units must the derivative have?
\end{enumerate}\end{promptbox}
\textbf{Workspace}\workspace
\needspace{8\baselineskip}
\subsection*{\pid{APP-014-A-01} — Advanced}
\begin{problembox}
A sensitivity model for error and differential approximation contains the response $R(x)=\sqrt{x}$. Around $x=114$, the input may increase by 1.4\%. Use differentials to estimate the change in $R$, and compare with the exact change.
\end{problembox}
\begin{promptbox}\begin{enumerate}
\item What quantity is changing?
\item With respect to which independent variable?
\item Is the requested rate average or instantaneous?
\item What units must the derivative have?
\end{enumerate}\end{promptbox}
\textbf{Workspace}\workspace
\begin{reflectionbox}\begin{enumerate}
\item Which sentence or datum in these problems supplied the main mathematical constraint?
\item Which assumptions were essential to the model, and how would the model change if one were removed?
\item When would an analytical solution become inconvenient enough to justify a numerical method?
\end{enumerate}\end{reflectionbox}
\part{Related Rates}
\chapter{Angular speed and angular rate of change}\label{app:A008}
\textbf{Application ID:} \pid{A008}\quad \textbf{Source locus:} P.4; 2.6; 5.8\quad \textbf{Source pages:} 38, 160, 162, 381
\section*{Application Context}
This family links two or more changing quantities through a constraint. Implicit differentiation converts the constraint into a relationship among rates. The canonical application family is \emph{Angular speed and angular rate of change}.
\section*{Calculus Toolkit}
\begin{tabularx}{\textwidth}{@{}lX@{}}\toprule
Primary technique & Chain rule and related rates \\
Prerequisites & implicit differentiation, geometry, derivative rules \\
Numerical extension & §5.7 for numerical derivatives; Ch. 9 if an implicit state equation must be solved \\
Analytical/numerical class & NUMERICAL-OPTIONAL \\
\bottomrule\end{tabularx}
\section*{Modeling Questions to Ask}\begin{itemize}
\item Which quantities vary with time?
\item What equation connects them?
\item Which rate is known and which rate is requested?
\item At what instant should numerical values be substituted?
\end{itemize}
\begin{errorbox}Treat all changing quantities as functions of time, differentiate the relationship first, and only then substitute the instant-specific data.\end{errorbox}
\section*{Problems}
\needspace{8\baselineskip}
\subsection*{\pid{APP-008-F-01} — Foundation}
\begin{problembox}
In a angular speed and angular rate of change scenario, a circular quantity has radius $r=4$ m increasing at $dr/dt=0.40$ m/s. How fast is its area changing at that instant?
\end{problembox}
\begin{promptbox}\begin{enumerate}
\item Which quantities vary with time?
\item What equation connects them?
\item Which rate is known and which rate is requested?
\item At what instant should numerical values be substituted?
\end{enumerate}\end{promptbox}
\textbf{Workspace}\workspace
\needspace{8\baselineskip}
\subsection*{\pid{APP-008-S-01} — Standard}
\begin{problembox}
A rigid segment of length 10 m is used in a angular speed and angular rate of change model. One endpoint moves horizontally away from a reference point at 1.00 m/s while the other endpoint moves vertically. When the horizontal distance is 6 m, how fast is the vertical distance changing?
\end{problembox}
\begin{promptbox}\begin{enumerate}
\item Which quantities vary with time?
\item What equation connects them?
\item Which rate is known and which rate is requested?
\item At what instant should numerical values be substituted?
\end{enumerate}\end{promptbox}
\textbf{Workspace}\workspace
\needspace{8\baselineskip}
\subsection*{\pid{APP-008-A-01} — Advanced}
\begin{problembox}
A nearly spherical volume in a angular speed and angular rate of change system is 1400 cm$^3$ and increasing at 21 cm$^3$/s. Assuming $V=\frac43\pi r^3$, estimate $dr/dt$ at that instant.
\end{problembox}
\begin{promptbox}\begin{enumerate}
\item Which quantities vary with time?
\item What equation connects them?
\item Which rate is known and which rate is requested?
\item At what instant should numerical values be substituted?
\end{enumerate}\end{promptbox}
\textbf{Workspace}\workspace
\begin{reflectionbox}\begin{enumerate}
\item Which sentence or datum in these problems supplied the main mathematical constraint?
\item Which assumptions were essential to the model, and how would the model change if one were removed?
\item When would an analytical solution become inconvenient enough to justify a numerical method?
\end{enumerate}\end{reflectionbox}
\chapter{Moving ladder, moving shadow, air-traffic and angle-of-elevation problems}\label{app:A009}
\textbf{Application ID:} \pid{A009}\quad \textbf{Source locus:} 2.6; 10; 12\quad \textbf{Source pages:} 93, 155–164, 749, 854
\section*{Application Context}
This family links two or more changing quantities through a constraint. Implicit differentiation converts the constraint into a relationship among rates. The canonical application family is \emph{Moving ladder, moving shadow, air-traffic and angle-of-elevation problems}.
\section*{Calculus Toolkit}
\begin{tabularx}{\textwidth}{@{}lX@{}}\toprule
Primary technique & Implicit differentiation and related rates \\
Prerequisites & implicit differentiation, geometry, derivative rules \\
Numerical extension & §5.7 for numerical derivatives; Ch. 9 if an implicit state equation must be solved \\
Analytical/numerical class & ANALYTICAL \\
\bottomrule\end{tabularx}
\section*{Modeling Questions to Ask}\begin{itemize}
\item Which quantities vary with time?
\item What equation connects them?
\item Which rate is known and which rate is requested?
\item At what instant should numerical values be substituted?
\end{itemize}
\begin{errorbox}Treat all changing quantities as functions of time, differentiate the relationship first, and only then substitute the instant-specific data.\end{errorbox}
\section*{Problems}
\needspace{8\baselineskip}
\subsection*{\pid{APP-009-F-01} — Foundation}
\begin{problembox}
In a moving ladder, moving shadow, air-traffic and angle-of-elevation problems scenario, a circular quantity has radius $r=5$ m increasing at $dr/dt=0.30$ m/s. How fast is its area changing at that instant?
\end{problembox}
\begin{promptbox}\begin{enumerate}
\item Which quantities vary with time?
\item What equation connects them?
\item Which rate is known and which rate is requested?
\item At what instant should numerical values be substituted?
\end{enumerate}\end{promptbox}
\textbf{Workspace}\workspace
\needspace{8\baselineskip}
\subsection*{\pid{APP-009-S-01} — Standard}
\begin{problembox}
A rigid segment of length 11 m is used in a moving ladder, moving shadow, air-traffic and angle-of-elevation problems model. One endpoint moves horizontally away from a reference point at 0.80 m/s while the other endpoint moves vertically. When the horizontal distance is 6 m, how fast is the vertical distance changing?
\end{problembox}
\begin{promptbox}\begin{enumerate}
\item Which quantities vary with time?
\item What equation connects them?
\item Which rate is known and which rate is requested?
\item At what instant should numerical values be substituted?
\end{enumerate}\end{promptbox}
\textbf{Workspace}\workspace
\needspace{8\baselineskip}
\subsection*{\pid{APP-009-A-01} — Advanced}
\begin{problembox}
A nearly spherical volume in a moving ladder, moving shadow, air-traffic and angle-of-elevation problems system is 1450 cm$^3$ and increasing at 23 cm$^3$/s. Assuming $V=\frac43\pi r^3$, estimate $dr/dt$ at that instant.
\end{problembox}
\begin{promptbox}\begin{enumerate}
\item Which quantities vary with time?
\item What equation connects them?
\item Which rate is known and which rate is requested?
\item At what instant should numerical values be substituted?
\end{enumerate}\end{promptbox}
\textbf{Workspace}\workspace
\begin{reflectionbox}\begin{enumerate}
\item Which sentence or datum in these problems supplied the main mathematical constraint?
\item Which assumptions were essential to the model, and how would the model change if one were removed?
\item When would an analytical solution become inconvenient enough to justify a numerical method?
\end{enumerate}\end{reflectionbox}
\chapter{Flow rate, oil leak and changing-volume systems}\label{app:A010}
\textbf{Application ID:} \pid{A010}\quad \textbf{Source locus:} 4.4–4.5; 5; 15\quad \textbf{Source pages:} 290, 294, 307, 351, 1109
\section*{Application Context}
This family links two or more changing quantities through a constraint. Implicit differentiation converts the constraint into a relationship among rates. The canonical application family is \emph{Flow rate, oil leak and changing-volume systems}.
\section*{Calculus Toolkit}
\begin{tabularx}{\textwidth}{@{}lX@{}}\toprule
Primary technique & Derivative/integral rate relations \\
Prerequisites & implicit differentiation, geometry, derivative rules \\
Numerical extension & §5.7 for numerical derivatives; Ch. 9 if an implicit state equation must be solved \\
Analytical/numerical class & NUMERICAL-OPTIONAL \\
\bottomrule\end{tabularx}
\section*{Modeling Questions to Ask}\begin{itemize}
\item Which quantities vary with time?
\item What equation connects them?
\item Which rate is known and which rate is requested?
\item At what instant should numerical values be substituted?
\end{itemize}
\begin{errorbox}Treat all changing quantities as functions of time, differentiate the relationship first, and only then substitute the instant-specific data.\end{errorbox}
\section*{Problems}
\needspace{8\baselineskip}
\subsection*{\pid{APP-010-F-01} — Foundation}
\begin{problembox}
In a flow rate, oil leak and changing-volume systems scenario, a circular quantity has radius $r=6$ m increasing at $dr/dt=0.35$ m/s. How fast is its area changing at that instant?
\end{problembox}
\begin{promptbox}\begin{enumerate}
\item Which quantities vary with time?
\item What equation connects them?
\item Which rate is known and which rate is requested?
\item At what instant should numerical values be substituted?
\end{enumerate}\end{promptbox}
\textbf{Workspace}\workspace
\needspace{8\baselineskip}
\subsection*{\pid{APP-010-S-01} — Standard}
\begin{problembox}
A rigid segment of length 12 m is used in a flow rate, oil leak and changing-volume systems model. One endpoint moves horizontally away from a reference point at 0.90 m/s while the other endpoint moves vertically. When the horizontal distance is 6 m, how fast is the vertical distance changing?
\end{problembox}
\begin{promptbox}\begin{enumerate}
\item Which quantities vary with time?
\item What equation connects them?
\item Which rate is known and which rate is requested?
\item At what instant should numerical values be substituted?
\end{enumerate}\end{promptbox}
\textbf{Workspace}\workspace
\needspace{8\baselineskip}
\subsection*{\pid{APP-010-A-01} — Advanced}
\begin{problembox}
A nearly spherical volume in a flow rate, oil leak and changing-volume systems system is 1500 cm$^3$ and increasing at 15 cm$^3$/s. Assuming $V=\frac43\pi r^3$, estimate $dr/dt$ at that instant.
\end{problembox}
\begin{promptbox}\begin{enumerate}
\item Which quantities vary with time?
\item What equation connects them?
\item Which rate is known and which rate is requested?
\item At what instant should numerical values be substituted?
\end{enumerate}\end{promptbox}
\textbf{Workspace}\workspace
\begin{reflectionbox}\begin{enumerate}
\item Which sentence or datum in these problems supplied the main mathematical constraint?
\item Which assumptions were essential to the model, and how would the model change if one were removed?
\item When would an analytical solution become inconvenient enough to justify a numerical method?
\end{enumerate}\end{reflectionbox}
\part{Nonlinear Equations and Root Finding}
\chapter{Break-even point and cost–revenue intersections}\label{app:A015}
\textbf{Application ID:} \pid{A015}\quad \textbf{Source locus:} P.2; 7.1\quad \textbf{Source pages:} 9, 41, 452
\section*{Application Context}
This family reduces an application condition to a scalar or low-dimensional nonlinear equation. Feasible roots must satisfy both the equation and the application domain. The canonical application family is \emph{Break-even point and cost–revenue intersections}.
\section*{Calculus Toolkit}
\begin{tabularx}{\textwidth}{@{}lX@{}}\toprule
Primary technique & Zeros/intersections of economic functions \\
Prerequisites & continuity, algebra, derivatives for Newton iteration \\
Numerical extension & Ch. 9, especially §§9.1–9.4 (bracketing, Brent, Newton) \\
Analytical/numerical class & NUMERICAL-NATURAL \\
\bottomrule\end{tabularx}
\section*{Modeling Questions to Ask}\begin{itemize}
\item What equation must equal zero?
\item What interval is physically feasible?
\item Can a root be bracketed?
\item How will the numerical answer be checked?
\end{itemize}
\begin{errorbox}A numerical root is meaningful only if the equation models the application and the root lies in the feasible domain. Check the residual.\end{errorbox}
\section*{Problems}
\needspace{8\baselineskip}
\subsection*{\pid{APP-015-F-01} — Foundation}
\begin{problembox}
A simplified equation arising in break-even point and cost–revenue intersections is $x^2+x-13=0$ with $x>0$. Show that the positive root is bracketed between two consecutive integers and compute it exactly.
\end{problembox}
\begin{promptbox}\begin{enumerate}
\item What equation must equal zero?
\item What interval is physically feasible?
\item Can a root be bracketed?
\item How will the numerical answer be checked?
\end{enumerate}\end{promptbox}
\textbf{Workspace}\workspace
\needspace{8\baselineskip}
\subsection*{\pid{APP-015-S-01} — Standard}
\begin{problembox}
A calibration equation for break-even point and cost–revenue intersections is $e^{-x}=x/2.0$. Use bisection on $[0,2]$ until the interval width is below $10^{-4}$, and report the solution.
\end{problembox}
\begin{promptbox}\begin{enumerate}
\item What equation must equal zero?
\item What interval is physically feasible?
\item Can a root be bracketed?
\item How will the numerical answer be checked?
\end{enumerate}\end{promptbox}
\textbf{Workspace}\workspace
\needspace{8\baselineskip}
\subsection*{\pid{APP-015-A-01} — Advanced}
\begin{problembox}
For a nonlinear model associated with break-even point and cost–revenue intersections, solve $\cos x=x/3$ using Newton's method from $x_0=1$. Carry out four iterations and check the residual.
\end{problembox}
\begin{promptbox}\begin{enumerate}
\item What equation must equal zero?
\item What interval is physically feasible?
\item Can a root be bracketed?
\item How will the numerical answer be checked?
\end{enumerate}\end{promptbox}
\textbf{Workspace}\workspace
\begin{reflectionbox}\begin{enumerate}
\item Which sentence or datum in these problems supplied the main mathematical constraint?
\item Which assumptions were essential to the model, and how would the model change if one were removed?
\item When would an analytical solution become inconvenient enough to justify a numerical method?
\end{enumerate}\end{reflectionbox}
\chapter{Boiling temperature, atmospheric-pressure/altitude inversions and calibration problems}\label{app:A016}
\textbf{Application ID:} \pid{A016}\quad \textbf{Source locus:} 5; 6; 13\quad \textbf{Source pages:} 323, 349, 439, 955
\section*{Application Context}
This family reduces an application condition to a scalar or low-dimensional nonlinear equation. Feasible roots must satisfy both the equation and the application domain. The canonical application family is \emph{Boiling temperature, atmospheric-pressure/altitude inversions and calibration problems}.
\section*{Calculus Toolkit}
\begin{tabularx}{\textwidth}{@{}lX@{}}\toprule
Primary technique & Solve nonlinear scalar equations \\
Prerequisites & continuity, algebra, derivatives for Newton iteration \\
Numerical extension & Ch. 9, especially §§9.1–9.4 (bracketing, Brent, Newton) \\
Analytical/numerical class & NUMERICAL-NATURAL \\
\bottomrule\end{tabularx}
\section*{Modeling Questions to Ask}\begin{itemize}
\item What equation must equal zero?
\item What interval is physically feasible?
\item Can a root be bracketed?
\item How will the numerical answer be checked?
\end{itemize}
\begin{errorbox}A numerical root is meaningful only if the equation models the application and the root lies in the feasible domain. Check the residual.\end{errorbox}
\section*{Problems}
\needspace{8\baselineskip}
\subsection*{\pid{APP-016-F-01} — Foundation}
\begin{problembox}
A simplified equation arising in boiling temperature, atmospheric-pressure/altitude inversions and calibration problems is $x^2+x-14=0$ with $x>0$. Show that the positive root is bracketed between two consecutive integers and compute it exactly.
\end{problembox}
\begin{promptbox}\begin{enumerate}
\item What equation must equal zero?
\item What interval is physically feasible?
\item Can a root be bracketed?
\item How will the numerical answer be checked?
\end{enumerate}\end{promptbox}
\textbf{Workspace}\workspace
\needspace{8\baselineskip}
\subsection*{\pid{APP-016-S-01} — Standard}
\begin{problembox}
A calibration equation for boiling temperature, atmospheric-pressure/altitude inversions and calibration problems is $e^{-x}=x/2.1$. Use bisection on $[0,2]$ until the interval width is below $10^{-4}$, and report the solution.
\end{problembox}
\begin{promptbox}\begin{enumerate}
\item What equation must equal zero?
\item What interval is physically feasible?
\item Can a root be bracketed?
\item How will the numerical answer be checked?
\end{enumerate}\end{promptbox}
\textbf{Workspace}\workspace
\needspace{8\baselineskip}
\subsection*{\pid{APP-016-A-01} — Advanced}
\begin{problembox}
For a nonlinear model associated with boiling temperature, atmospheric-pressure/altitude inversions and calibration problems, solve $\cos x=x/4$ using Newton's method from $x_0=1$. Carry out four iterations and check the residual.
\end{problembox}
\begin{promptbox}\begin{enumerate}
\item What equation must equal zero?
\item What interval is physically feasible?
\item Can a root be bracketed?
\item How will the numerical answer be checked?
\end{enumerate}\end{promptbox}
\textbf{Workspace}\workspace
\begin{reflectionbox}\begin{enumerate}
\item Which sentence or datum in these problems supplied the main mathematical constraint?
\item Which assumptions were essential to the model, and how would the model change if one were removed?
\item When would an analytical solution become inconvenient enough to justify a numerical method?
\end{enumerate}\end{reflectionbox}
\chapter{Optical/refraction and geometric intersection problems}\label{app:A017}
\textbf{Application ID:} \pid{A017}\quad \textbf{Source locus:} 3; 10; 13\quad \textbf{Source pages:} 197, 688, 695, 699, 963
\section*{Application Context}
This family reduces an application condition to a scalar or low-dimensional nonlinear equation. Feasible roots must satisfy both the equation and the application domain. The canonical application family is \emph{Optical/refraction and geometric intersection problems}.
\section*{Calculus Toolkit}
\begin{tabularx}{\textwidth}{@{}lX@{}}\toprule
Primary technique & Nonlinear equations from geometry \\
Prerequisites & continuity, algebra, derivatives for Newton iteration \\
Numerical extension & Ch. 9, especially §§9.1–9.4 (bracketing, Brent, Newton) \\
Analytical/numerical class & NUMERICAL-NATURAL \\
\bottomrule\end{tabularx}
\section*{Modeling Questions to Ask}\begin{itemize}
\item What equation must equal zero?
\item What interval is physically feasible?
\item Can a root be bracketed?
\item How will the numerical answer be checked?
\end{itemize}
\begin{errorbox}A numerical root is meaningful only if the equation models the application and the root lies in the feasible domain. Check the residual.\end{errorbox}
\section*{Problems}
\needspace{8\baselineskip}
\subsection*{\pid{APP-017-F-01} — Foundation}
\begin{problembox}
A simplified equation arising in optical/refraction and geometric intersection problems is $x^2+x-15=0$ with $x>0$. Show that the positive root is bracketed between two consecutive integers and compute it exactly.
\end{problembox}
\begin{promptbox}\begin{enumerate}
\item What equation must equal zero?
\item What interval is physically feasible?
\item Can a root be bracketed?
\item How will the numerical answer be checked?
\end{enumerate}\end{promptbox}
\textbf{Workspace}\workspace
\needspace{8\baselineskip}
\subsection*{\pid{APP-017-S-01} — Standard}
\begin{problembox}
A calibration equation for optical/refraction and geometric intersection problems is $e^{-x}=x/2.2$. Use bisection on $[0,2]$ until the interval width is below $10^{-4}$, and report the solution.
\end{problembox}
\begin{promptbox}\begin{enumerate}
\item What equation must equal zero?
\item What interval is physically feasible?
\item Can a root be bracketed?
\item How will the numerical answer be checked?
\end{enumerate}\end{promptbox}
\textbf{Workspace}\workspace
\needspace{8\baselineskip}
\subsection*{\pid{APP-017-A-01} — Advanced}
\begin{problembox}
For a nonlinear model associated with optical/refraction and geometric intersection problems, solve $\cos x=x/5$ using Newton's method from $x_0=1$. Carry out four iterations and check the residual.
\end{problembox}
\begin{promptbox}\begin{enumerate}
\item What equation must equal zero?
\item What interval is physically feasible?
\item Can a root be bracketed?
\item How will the numerical answer be checked?
\end{enumerate}\end{promptbox}
\textbf{Workspace}\workspace
\begin{reflectionbox}\begin{enumerate}
\item Which sentence or datum in these problems supplied the main mathematical constraint?
\item Which assumptions were essential to the model, and how would the model change if one were removed?
\item When would an analytical solution become inconvenient enough to justify a numerical method?
\end{enumerate}\end{reflectionbox}
\chapter{Polynomial design/intersection and machine-part equations}\label{app:A018}
\textbf{Application ID:} \pid{A018}\quad \textbf{Source locus:} 3; 7\quad \textbf{Source pages:} 190, 463, 471
\section*{Application Context}
This family reduces an application condition to a scalar or low-dimensional nonlinear equation. Feasible roots must satisfy both the equation and the application domain. The canonical application family is \emph{Polynomial design/intersection and machine-part equations}.
\section*{Calculus Toolkit}
\begin{tabularx}{\textwidth}{@{}lX@{}}\toprule
Primary technique & Polynomial zeros \\
Prerequisites & continuity, algebra, derivatives for Newton iteration \\
Numerical extension & Ch. 9, especially §§9.1–9.4 (bracketing, Brent, Newton) \\
Analytical/numerical class & NUMERICAL-NATURAL \\
\bottomrule\end{tabularx}
\section*{Modeling Questions to Ask}\begin{itemize}
\item What equation must equal zero?
\item What interval is physically feasible?
\item Can a root be bracketed?
\item How will the numerical answer be checked?
\end{itemize}
\begin{errorbox}A numerical root is meaningful only if the equation models the application and the root lies in the feasible domain. Check the residual.\end{errorbox}
\section*{Problems}
\needspace{8\baselineskip}
\subsection*{\pid{APP-018-F-01} — Foundation}
\begin{problembox}
A simplified equation arising in polynomial design/intersection and machine-part equations is $x^2+x-10=0$ with $x>0$. Show that the positive root is bracketed between two consecutive integers and compute it exactly.
\end{problembox}
\begin{promptbox}\begin{enumerate}
\item What equation must equal zero?
\item What interval is physically feasible?
\item Can a root be bracketed?
\item How will the numerical answer be checked?
\end{enumerate}\end{promptbox}
\textbf{Workspace}\workspace
\needspace{8\baselineskip}
\subsection*{\pid{APP-018-S-01} — Standard}
\begin{problembox}
A calibration equation for polynomial design/intersection and machine-part equations is $e^{-x}=x/2.3$. Use bisection on $[0,2]$ until the interval width is below $10^{-4}$, and report the solution.
\end{problembox}
\begin{promptbox}\begin{enumerate}
\item What equation must equal zero?
\item What interval is physically feasible?
\item Can a root be bracketed?
\item How will the numerical answer be checked?
\end{enumerate}\end{promptbox}
\textbf{Workspace}\workspace
\needspace{8\baselineskip}
\subsection*{\pid{APP-018-A-01} — Advanced}
\begin{problembox}
For a nonlinear model associated with polynomial design/intersection and machine-part equations, solve $\cos x=x/3$ using Newton's method from $x_0=1$. Carry out four iterations and check the residual.
\end{problembox}
\begin{promptbox}\begin{enumerate}
\item What equation must equal zero?
\item What interval is physically feasible?
\item Can a root be bracketed?
\item How will the numerical answer be checked?
\end{enumerate}\end{promptbox}
\textbf{Workspace}\workspace
\begin{reflectionbox}\begin{enumerate}
\item Which sentence or datum in these problems supplied the main mathematical constraint?
\item Which assumptions were essential to the model, and how would the model change if one were removed?
\item When would an analytical solution become inconvenient enough to justify a numerical method?
\end{enumerate}\end{reflectionbox}
\part{Optimization}
\chapter{Maximum area, minimum area, perimeter, length and surface area}\label{app:A019}
\textbf{Application ID:} \pid{A019}\quad \textbf{Source locus:} 3.7; 13.8–13.10\quad \textbf{Source pages:} 43, 221–226, 244, 968
\section*{Application Context}
This family turns design or decision questions into objective functions subject to geometric, physical, or economic constraints. The canonical application family is \emph{Maximum area, minimum area, perimeter, length and surface area}.
\section*{Calculus Toolkit}
\begin{tabularx}{\textwidth}{@{}lX@{}}\toprule
Primary technique & Scalar and constrained optimization \\
Prerequisites & derivatives, critical points, constraints \\
Numerical extension & Ch. 10, especially §§10.2–10.9 \\
Analytical/numerical class & NUMERICAL-NATURAL \\
\bottomrule\end{tabularx}
\section*{Modeling Questions to Ask}\begin{itemize}
\item What is the objective function?
\item What constraints define the feasible set?
\item Can the objective be reduced to one variable?
\item How will the extremum be verified?
\end{itemize}
\begin{errorbox}A critical point is only a candidate. Enforce constraints and compare endpoints or otherwise establish that the requested extremum has been found.\end{errorbox}
\section*{Problems}
\needspace{8\baselineskip}
\subsection*{\pid{APP-019-F-01} — Foundation}
\begin{problembox}
A design related to maximum area, minimum area, perimeter, length and surface area uses a rectangular region with total perimeter 68 m. Determine the dimensions that maximize area.
\end{problembox}
\begin{promptbox}\begin{enumerate}
\item What is the objective function?
\item What constraints define the feasible set?
\item Can the objective be reduced to one variable?
\item How will the extremum be verified?
\end{enumerate}\end{promptbox}
\textbf{Workspace}\workspace
\needspace{8\baselineskip}
\subsection*{\pid{APP-019-S-01} — Standard}
\begin{problembox}
A cost model for maximum area, minimum area, perimeter, length and surface area is $C(x)=2950/x+88+0.6x$ for $x>0$. Find the operating value $x$ that minimizes cost and the minimum cost.
\end{problembox}
\begin{promptbox}\begin{enumerate}
\item What is the objective function?
\item What constraints define the feasible set?
\item Can the objective be reduced to one variable?
\item How will the extremum be verified?
\end{enumerate}\end{promptbox}
\textbf{Workspace}\workspace
\needspace{8\baselineskip}
\subsection*{\pid{APP-019-A-01} — Advanced}
\begin{problembox}
A two-factor design in maximum area, minimum area, perimeter, length and surface area has nonnegative decision variables $x$ and $y$ constrained by $x+y=24$. Maximize the performance measure $P=xy$ using a Lagrange-multiplier argument.
\end{problembox}
\begin{promptbox}\begin{enumerate}
\item What is the objective function?
\item What constraints define the feasible set?
\item Can the objective be reduced to one variable?
\item How will the extremum be verified?
\end{enumerate}\end{promptbox}
\textbf{Workspace}\workspace
\begin{reflectionbox}\begin{enumerate}
\item Which sentence or datum in these problems supplied the main mathematical constraint?
\item Which assumptions were essential to the model, and how would the model change if one were removed?
\item When would an analytical solution become inconvenient enough to justify a numerical method?
\end{enumerate}\end{reflectionbox}
\chapter{Maximum volume and packaging/design optimization}\label{app:A020}
\textbf{Application ID:} \pid{A020}\quad \textbf{Source locus:} 3.7; 13.8–13.10\quad \textbf{Source pages:} 219–227, 948, 953, 962, 963
\section*{Application Context}
This family turns design or decision questions into objective functions subject to geometric, physical, or economic constraints. The canonical application family is \emph{Maximum volume and packaging/design optimization}.
\section*{Calculus Toolkit}
\begin{tabularx}{\textwidth}{@{}lX@{}}\toprule
Primary technique & Extrema in one and several variables \\
Prerequisites & derivatives, critical points, constraints \\
Numerical extension & Ch. 10, especially §§10.2–10.9 \\
Analytical/numerical class & NUMERICAL-NATURAL \\
\bottomrule\end{tabularx}
\section*{Modeling Questions to Ask}\begin{itemize}
\item What is the objective function?
\item What constraints define the feasible set?
\item Can the objective be reduced to one variable?
\item How will the extremum be verified?
\end{itemize}
\begin{errorbox}A critical point is only a candidate. Enforce constraints and compare endpoints or otherwise establish that the requested extremum has been found.\end{errorbox}
\section*{Problems}
\needspace{8\baselineskip}
\subsection*{\pid{APP-020-F-01} — Foundation}
\begin{problembox}
A design related to maximum volume and packaging/design optimization uses a rectangular region with total perimeter 60 m. Determine the dimensions that maximize area.
\end{problembox}
\begin{promptbox}\begin{enumerate}
\item What is the objective function?
\item What constraints define the feasible set?
\item Can the objective be reduced to one variable?
\item How will the extremum be verified?
\end{enumerate}\end{promptbox}
\textbf{Workspace}\workspace
\needspace{8\baselineskip}
\subsection*{\pid{APP-020-S-01} — Standard}
\begin{problembox}
A cost model for maximum volume and packaging/design optimization is $C(x)=3000/x+80+0.7x$ for $x>0$. Find the operating value $x$ that minimizes cost and the minimum cost.
\end{problembox}
\begin{promptbox}\begin{enumerate}
\item What is the objective function?
\item What constraints define the feasible set?
\item Can the objective be reduced to one variable?
\item How will the extremum be verified?
\end{enumerate}\end{promptbox}
\textbf{Workspace}\workspace
\needspace{8\baselineskip}
\subsection*{\pid{APP-020-A-01} — Advanced}
\begin{problembox}
A two-factor design in maximum volume and packaging/design optimization has nonnegative decision variables $x$ and $y$ constrained by $x+y=20$. Maximize the performance measure $P=xy$ using a Lagrange-multiplier argument.
\end{problembox}
\begin{promptbox}\begin{enumerate}
\item What is the objective function?
\item What constraints define the feasible set?
\item Can the objective be reduced to one variable?
\item How will the extremum be verified?
\end{enumerate}\end{promptbox}
\textbf{Workspace}\workspace
\begin{reflectionbox}\begin{enumerate}
\item Which sentence or datum in these problems supplied the main mathematical constraint?
\item Which assumptions were essential to the model, and how would the model change if one were removed?
\item When would an analytical solution become inconvenient enough to justify a numerical method?
\end{enumerate}\end{reflectionbox}
\chapter{Minimum distance, time and travel-path optimization}\label{app:A021}
\textbf{Application ID:} \pid{A021}\quad \textbf{Source locus:} 3.7; 10; 13\quad \textbf{Source pages:} 221–228, 234, 245
\section*{Application Context}
This family turns design or decision questions into objective functions subject to geometric, physical, or economic constraints. The canonical application family is \emph{Minimum distance, time and travel-path optimization}.
\section*{Calculus Toolkit}
\begin{tabularx}{\textwidth}{@{}lX@{}}\toprule
Primary technique & Geometric optimization \\
Prerequisites & derivatives, critical points, constraints \\
Numerical extension & Ch. 10, especially §§10.2–10.9 \\
Analytical/numerical class & NUMERICAL-NATURAL \\
\bottomrule\end{tabularx}
\section*{Modeling Questions to Ask}\begin{itemize}
\item What is the objective function?
\item What constraints define the feasible set?
\item Can the objective be reduced to one variable?
\item How will the extremum be verified?
\end{itemize}
\begin{errorbox}A critical point is only a candidate. Enforce constraints and compare endpoints or otherwise establish that the requested extremum has been found.\end{errorbox}
\section*{Problems}
\needspace{8\baselineskip}
\subsection*{\pid{APP-021-F-01} — Foundation}
\begin{problembox}
A design related to minimum distance, time and travel-path optimization uses a rectangular region with total perimeter 62 m. Determine the dimensions that maximize area.
\end{problembox}
\begin{promptbox}\begin{enumerate}
\item What is the objective function?
\item What constraints define the feasible set?
\item Can the objective be reduced to one variable?
\item How will the extremum be verified?
\end{enumerate}\end{promptbox}
\textbf{Workspace}\workspace
\needspace{8\baselineskip}
\subsection*{\pid{APP-021-S-01} — Standard}
\begin{problembox}
A cost model for minimum distance, time and travel-path optimization is $C(x)=3050/x+82+0.5x$ for $x>0$. Find the operating value $x$ that minimizes cost and the minimum cost.
\end{problembox}
\begin{promptbox}\begin{enumerate}
\item What is the objective function?
\item What constraints define the feasible set?
\item Can the objective be reduced to one variable?
\item How will the extremum be verified?
\end{enumerate}\end{promptbox}
\textbf{Workspace}\workspace
\needspace{8\baselineskip}
\subsection*{\pid{APP-021-A-01} — Advanced}
\begin{problembox}
A two-factor design in minimum distance, time and travel-path optimization has nonnegative decision variables $x$ and $y$ constrained by $x+y=21$. Maximize the performance measure $P=xy$ using a Lagrange-multiplier argument.
\end{problembox}
\begin{promptbox}\begin{enumerate}
\item What is the objective function?
\item What constraints define the feasible set?
\item Can the objective be reduced to one variable?
\item How will the extremum be verified?
\end{enumerate}\end{promptbox}
\textbf{Workspace}\workspace
\begin{reflectionbox}\begin{enumerate}
\item Which sentence or datum in these problems supplied the main mathematical constraint?
\item Which assumptions were essential to the model, and how would the model change if one were removed?
\item When would an analytical solution become inconvenient enough to justify a numerical method?
\end{enumerate}\end{reflectionbox}
\chapter{Beam strength, illumination, engine efficiency and mechanical design}\label{app:A022}
\textbf{Application ID:} \pid{A022}\quad \textbf{Source locus:} 3.7; 13\quad \textbf{Source pages:} 207, 226, 245, 453, 797
\section*{Application Context}
This family turns design or decision questions into objective functions subject to geometric, physical, or economic constraints. The canonical application family is \emph{Beam strength, illumination, engine efficiency and mechanical design}.
\section*{Calculus Toolkit}
\begin{tabularx}{\textwidth}{@{}lX@{}}\toprule
Primary technique & Engineering objective functions \\
Prerequisites & derivatives, critical points, constraints \\
Numerical extension & Ch. 10, especially §§10.2–10.9 \\
Analytical/numerical class & NUMERICAL-NATURAL \\
\bottomrule\end{tabularx}
\section*{Modeling Questions to Ask}\begin{itemize}
\item What is the objective function?
\item What constraints define the feasible set?
\item Can the objective be reduced to one variable?
\item How will the extremum be verified?
\end{itemize}
\begin{errorbox}A critical point is only a candidate. Enforce constraints and compare endpoints or otherwise establish that the requested extremum has been found.\end{errorbox}
\section*{Problems}
\needspace{8\baselineskip}
\subsection*{\pid{APP-022-F-01} — Foundation}
\begin{problembox}
A design related to beam strength, illumination, engine efficiency and mechanical design uses a rectangular region with total perimeter 64 m. Determine the dimensions that maximize area.
\end{problembox}
\begin{promptbox}\begin{enumerate}
\item What is the objective function?
\item What constraints define the feasible set?
\item Can the objective be reduced to one variable?
\item How will the extremum be verified?
\end{enumerate}\end{promptbox}
\textbf{Workspace}\workspace
\needspace{8\baselineskip}
\subsection*{\pid{APP-022-S-01} — Standard}
\begin{problembox}
A cost model for beam strength, illumination, engine efficiency and mechanical design is $C(x)=3100/x+84+0.6x$ for $x>0$. Find the operating value $x$ that minimizes cost and the minimum cost.
\end{problembox}
\begin{promptbox}\begin{enumerate}
\item What is the objective function?
\item What constraints define the feasible set?
\item Can the objective be reduced to one variable?
\item How will the extremum be verified?
\end{enumerate}\end{promptbox}
\textbf{Workspace}\workspace
\needspace{8\baselineskip}
\subsection*{\pid{APP-022-A-01} — Advanced}
\begin{problembox}
A two-factor design in beam strength, illumination, engine efficiency and mechanical design has nonnegative decision variables $x$ and $y$ constrained by $x+y=22$. Maximize the performance measure $P=xy$ using a Lagrange-multiplier argument.
\end{problembox}
\begin{promptbox}\begin{enumerate}
\item What is the objective function?
\item What constraints define the feasible set?
\item Can the objective be reduced to one variable?
\item How will the extremum be verified?
\end{enumerate}\end{promptbox}
\textbf{Workspace}\workspace
\begin{reflectionbox}\begin{enumerate}
\item Which sentence or datum in these problems supplied the main mathematical constraint?
\item Which assumptions were essential to the model, and how would the model change if one were removed?
\item When would an analytical solution become inconvenient enough to justify a numerical method?
\end{enumerate}\end{reflectionbox}
\chapter{Highway design, honeycomb and structural geometry optimization}\label{app:A023}
\textbf{Application ID:} \pid{A023}\quad \textbf{Source locus:} 3; 12\quad \textbf{Source pages:} 173, 197, 870
\section*{Application Context}
This family turns design or decision questions into objective functions subject to geometric, physical, or economic constraints. The canonical application family is \emph{Highway design, honeycomb and structural geometry optimization}.
\section*{Calculus Toolkit}
\begin{tabularx}{\textwidth}{@{}lX@{}}\toprule
Primary technique & Extrema and curvature constraints \\
Prerequisites & derivatives, critical points, constraints \\
Numerical extension & Ch. 10, especially §§10.2–10.9 \\
Analytical/numerical class & NUMERICAL-NATURAL \\
\bottomrule\end{tabularx}
\section*{Modeling Questions to Ask}\begin{itemize}
\item What is the objective function?
\item What constraints define the feasible set?
\item Can the objective be reduced to one variable?
\item How will the extremum be verified?
\end{itemize}
\begin{errorbox}A critical point is only a candidate. Enforce constraints and compare endpoints or otherwise establish that the requested extremum has been found.\end{errorbox}
\section*{Problems}
\needspace{8\baselineskip}
\subsection*{\pid{APP-023-F-01} — Foundation}
\begin{problembox}
A design related to highway design, honeycomb and structural geometry optimization uses a rectangular region with total perimeter 66 m. Determine the dimensions that maximize area.
\end{problembox}
\begin{promptbox}\begin{enumerate}
\item What is the objective function?
\item What constraints define the feasible set?
\item Can the objective be reduced to one variable?
\item How will the extremum be verified?
\end{enumerate}\end{promptbox}
\textbf{Workspace}\workspace
\needspace{8\baselineskip}
\subsection*{\pid{APP-023-S-01} — Standard}
\begin{problembox}
A cost model for highway design, honeycomb and structural geometry optimization is $C(x)=3150/x+86+0.7x$ for $x>0$. Find the operating value $x$ that minimizes cost and the minimum cost.
\end{problembox}
\begin{promptbox}\begin{enumerate}
\item What is the objective function?
\item What constraints define the feasible set?
\item Can the objective be reduced to one variable?
\item How will the extremum be verified?
\end{enumerate}\end{promptbox}
\textbf{Workspace}\workspace
\needspace{8\baselineskip}
\subsection*{\pid{APP-023-A-01} — Advanced}
\begin{problembox}
A two-factor design in highway design, honeycomb and structural geometry optimization has nonnegative decision variables $x$ and $y$ constrained by $x+y=23$. Maximize the performance measure $P=xy$ using a Lagrange-multiplier argument.
\end{problembox}
\begin{promptbox}\begin{enumerate}
\item What is the objective function?
\item What constraints define the feasible set?
\item Can the objective be reduced to one variable?
\item How will the extremum be verified?
\end{enumerate}\end{promptbox}
\textbf{Workspace}\workspace
\begin{reflectionbox}\begin{enumerate}
\item Which sentence or datum in these problems supplied the main mathematical constraint?
\item Which assumptions were essential to the model, and how would the model change if one were removed?
\item When would an analytical solution become inconvenient enough to justify a numerical method?
\end{enumerate}\end{reflectionbox}
\chapter{Inventory management/replenishment and manufacturing optimization}\label{app:A024}
\textbf{Application ID:} \pid{A024}\quad \textbf{Source locus:} 1; 2; 7\quad \textbf{Source pages:} 85, 121, 130, 459, 463
\section*{Application Context}
This family turns design or decision questions into objective functions subject to geometric, physical, or economic constraints. The canonical application family is \emph{Inventory management/replenishment and manufacturing optimization}.
\section*{Calculus Toolkit}
\begin{tabularx}{\textwidth}{@{}lX@{}}\toprule
Primary technique & Cost/objective optimization \\
Prerequisites & derivatives, critical points, constraints \\
Numerical extension & Ch. 10, especially §§10.2–10.9 \\
Analytical/numerical class & NUMERICAL-NATURAL \\
\bottomrule\end{tabularx}
\section*{Modeling Questions to Ask}\begin{itemize}
\item What is the objective function?
\item What constraints define the feasible set?
\item Can the objective be reduced to one variable?
\item How will the extremum be verified?
\end{itemize}
\begin{errorbox}A critical point is only a candidate. Enforce constraints and compare endpoints or otherwise establish that the requested extremum has been found.\end{errorbox}
\section*{Problems}
\needspace{8\baselineskip}
\subsection*{\pid{APP-024-F-01} — Foundation}
\begin{problembox}
A design related to inventory management/replenishment and manufacturing optimization uses a rectangular region with total perimeter 68 m. Determine the dimensions that maximize area.
\end{problembox}
\begin{promptbox}\begin{enumerate}
\item What is the objective function?
\item What constraints define the feasible set?
\item Can the objective be reduced to one variable?
\item How will the extremum be verified?
\end{enumerate}\end{promptbox}
\textbf{Workspace}\workspace
\needspace{8\baselineskip}
\subsection*{\pid{APP-024-S-01} — Standard}
\begin{problembox}
A cost model for inventory management/replenishment and manufacturing optimization is $C(x)=3200/x+88+0.5x$ for $x>0$. Find the operating value $x$ that minimizes cost and the minimum cost.
\end{problembox}
\begin{promptbox}\begin{enumerate}
\item What is the objective function?
\item What constraints define the feasible set?
\item Can the objective be reduced to one variable?
\item How will the extremum be verified?
\end{enumerate}\end{promptbox}
\textbf{Workspace}\workspace
\needspace{8\baselineskip}
\subsection*{\pid{APP-024-A-01} — Advanced}
\begin{problembox}
A two-factor design in inventory management/replenishment and manufacturing optimization has nonnegative decision variables $x$ and $y$ constrained by $x+y=24$. Maximize the performance measure $P=xy$ using a Lagrange-multiplier argument.
\end{problembox}
\begin{promptbox}\begin{enumerate}
\item What is the objective function?
\item What constraints define the feasible set?
\item Can the objective be reduced to one variable?
\item How will the extremum be verified?
\end{enumerate}\end{promptbox}
\textbf{Workspace}\workspace
\begin{reflectionbox}\begin{enumerate}
\item Which sentence or datum in these problems supplied the main mathematical constraint?
\item Which assumptions were essential to the model, and how would the model change if one were removed?
\item When would an analytical solution become inconvenient enough to justify a numerical method?
\end{enumerate}\end{reflectionbox}
\chapter{Maximum profit/revenue and minimum cost}\label{app:A025}
\textbf{Application ID:} \pid{A025}\quad \textbf{Source locus:} 3.7; 13.9–13.10\quad \textbf{Source pages:} 226–227, 949, 953, 963, 966–967
\section*{Application Context}
This family turns design or decision questions into objective functions subject to geometric, physical, or economic constraints. The canonical application family is \emph{Maximum profit/revenue and minimum cost}.
\section*{Calculus Toolkit}
\begin{tabularx}{\textwidth}{@{}lX@{}}\toprule
Primary technique & Business optimization and constraints \\
Prerequisites & derivatives, critical points, constraints \\
Numerical extension & Ch. 10, especially §§10.2–10.9 \\
Analytical/numerical class & NUMERICAL-NATURAL \\
\bottomrule\end{tabularx}
\section*{Modeling Questions to Ask}\begin{itemize}
\item What is the objective function?
\item What constraints define the feasible set?
\item Can the objective be reduced to one variable?
\item How will the extremum be verified?
\end{itemize}
\begin{errorbox}A critical point is only a candidate. Enforce constraints and compare endpoints or otherwise establish that the requested extremum has been found.\end{errorbox}
\section*{Problems}
\needspace{8\baselineskip}
\subsection*{\pid{APP-025-F-01} — Foundation}
\begin{problembox}
A design related to maximum profit/revenue and minimum cost uses a rectangular region with total perimeter 60 m. Determine the dimensions that maximize area.
\end{problembox}
\begin{promptbox}\begin{enumerate}
\item What is the objective function?
\item What constraints define the feasible set?
\item Can the objective be reduced to one variable?
\item How will the extremum be verified?
\end{enumerate}\end{promptbox}
\textbf{Workspace}\workspace
\needspace{8\baselineskip}
\subsection*{\pid{APP-025-S-01} — Standard}
\begin{problembox}
A cost model for maximum profit/revenue and minimum cost is $C(x)=3250/x+80+0.6x$ for $x>0$. Find the operating value $x$ that minimizes cost and the minimum cost.
\end{problembox}
\begin{promptbox}\begin{enumerate}
\item What is the objective function?
\item What constraints define the feasible set?
\item Can the objective be reduced to one variable?
\item How will the extremum be verified?
\end{enumerate}\end{promptbox}
\textbf{Workspace}\workspace
\needspace{8\baselineskip}
\subsection*{\pid{APP-025-A-01} — Advanced}
\begin{problembox}
A two-factor design in maximum profit/revenue and minimum cost has nonnegative decision variables $x$ and $y$ constrained by $x+y=20$. Maximize the performance measure $P=xy$ using a Lagrange-multiplier argument.
\end{problembox}
\begin{promptbox}\begin{enumerate}
\item What is the objective function?
\item What constraints define the feasible set?
\item Can the objective be reduced to one variable?
\item How will the extremum be verified?
\end{enumerate}\end{promptbox}
\textbf{Workspace}\workspace
\begin{reflectionbox}\begin{enumerate}
\item Which sentence or datum in these problems supplied the main mathematical constraint?
\item Which assumptions were essential to the model, and how would the model change if one were removed?
\item When would an analytical solution become inconvenient enough to justify a numerical method?
\end{enumerate}\end{reflectionbox}
\chapter{Building a pipeline}\label{app:A026}
\textbf{Application ID:} \pid{A026}\quad \textbf{Source locus:} 13.9\quad \textbf{Source pages:} 955
\section*{Application Context}
This family turns design or decision questions into objective functions subject to geometric, physical, or economic constraints. The canonical application family is \emph{Building a pipeline}.
\section*{Calculus Toolkit}
\begin{tabularx}{\textwidth}{@{}lX@{}}\toprule
Primary technique & Multivariable geometric optimization \\
Prerequisites & derivatives, critical points, constraints \\
Numerical extension & Ch. 10, especially §§10.2–10.9 \\
Analytical/numerical class & NUMERICAL-NATURAL \\
\bottomrule\end{tabularx}
\section*{Modeling Questions to Ask}\begin{itemize}
\item What is the objective function?
\item What constraints define the feasible set?
\item Can the objective be reduced to one variable?
\item How will the extremum be verified?
\end{itemize}
\begin{errorbox}A critical point is only a candidate. Enforce constraints and compare endpoints or otherwise establish that the requested extremum has been found.\end{errorbox}
\section*{Problems}
\needspace{8\baselineskip}
\subsection*{\pid{APP-026-F-01} — Foundation}
\begin{problembox}
A design related to building a pipeline uses a rectangular region with total perimeter 62 m. Determine the dimensions that maximize area.
\end{problembox}
\begin{promptbox}\begin{enumerate}
\item What is the objective function?
\item What constraints define the feasible set?
\item Can the objective be reduced to one variable?
\item How will the extremum be verified?
\end{enumerate}\end{promptbox}
\textbf{Workspace}\workspace
\needspace{8\baselineskip}
\subsection*{\pid{APP-026-S-01} — Standard}
\begin{problembox}
A cost model for building a pipeline is $C(x)=3300/x+82+0.7x$ for $x>0$. Find the operating value $x$ that minimizes cost and the minimum cost.
\end{problembox}
\begin{promptbox}\begin{enumerate}
\item What is the objective function?
\item What constraints define the feasible set?
\item Can the objective be reduced to one variable?
\item How will the extremum be verified?
\end{enumerate}\end{promptbox}
\textbf{Workspace}\workspace
\needspace{8\baselineskip}
\subsection*{\pid{APP-026-A-01} — Advanced}
\begin{problembox}
A two-factor design in building a pipeline has nonnegative decision variables $x$ and $y$ constrained by $x+y=21$. Maximize the performance measure $P=xy$ using a Lagrange-multiplier argument.
\end{problembox}
\begin{promptbox}\begin{enumerate}
\item What is the objective function?
\item What constraints define the feasible set?
\item Can the objective be reduced to one variable?
\item How will the extremum be verified?
\end{enumerate}\end{promptbox}
\textbf{Workspace}\workspace
\begin{reflectionbox}\begin{enumerate}
\item Which sentence or datum in these problems supplied the main mathematical constraint?
\item Which assumptions were essential to the model, and how would the model change if one were removed?
\item When would an analytical solution become inconvenient enough to justify a numerical method?
\end{enumerate}\end{reflectionbox}
\part{Integral Applications}
\chapter{Area and area between curves}\label{app:A027}
\textbf{Application ID:} \pid{A027}\quad \textbf{Source locus:} 4; 7.1\quad \textbf{Source pages:} 120, 130, 157, 240, 260, 444–453, 607, 678
\section*{Application Context}
This family models totals by accumulating infinitesimal contributions. Correct bounds and units are as important as evaluating the antiderivative or quadrature. The canonical application family is \emph{Area and area between curves}.
\section*{Calculus Toolkit}
\begin{tabularx}{\textwidth}{@{}lX@{}}\toprule
Primary technique & Definite integration and accumulation \\
Prerequisites & definite integrals, units, geometric/physical modeling \\
Numerical extension & Ch. 4, especially §§4.1–4.8 \\
Analytical/numerical class & NUMERICAL-NATURAL \\
\bottomrule\end{tabularx}
\section*{Modeling Questions to Ask}\begin{itemize}
\item What infinitesimal contribution is accumulated?
\item What variable should be integrated over?
\item What are the correct bounds?
\item What units should the final integral have?
\end{itemize}
\begin{errorbox}Build the infinitesimal contribution before integrating. Check bounds and units of the integrand and final accumulated quantity.\end{errorbox}
\section*{Problems}
\needspace{8\baselineskip}
\subsection*{\pid{APP-027-F-01} — Foundation}
\begin{problembox}
A cumulative quantity in area and area between curves is modeled by the rate $r(x)=x^2+2$ over $0\le x\le 9$. Find the total accumulated amount.
\end{problembox}
\begin{promptbox}\begin{enumerate}
\item What infinitesimal contribution is accumulated?
\item What variable should be integrated over?
\item What are the correct bounds?
\item What units should the final integral have?
\end{enumerate}\end{promptbox}
\textbf{Workspace}\workspace
\needspace{8\baselineskip}
\subsection*{\pid{APP-027-S-01} — Standard}
\begin{problembox}
In a area and area between curves application, liquid fills a vertical prismatic tank of height 8 m and horizontal cross-sectional area 2 m$^2$. Use a weight density of 9.8 kN/m$^3$. Find the work required to pump all liquid just to the top edge.
\end{problembox}
\begin{promptbox}\begin{enumerate}
\item What infinitesimal contribution is accumulated?
\item What variable should be integrated over?
\item What are the correct bounds?
\item What units should the final integral have?
\end{enumerate}\end{promptbox}
\textbf{Workspace}\workspace
\needspace{8\baselineskip}
\subsection*{\pid{APP-027-A-01} — Advanced}
\begin{problembox}
A non-elementary accumulation arising in area and area between curves is $I=\int_0^1 e^{-x^2}\,dx$. Approximate $I$ with composite Simpson's rule using four equal subintervals, and explain why a numerical method is appropriate.
\end{problembox}
\begin{promptbox}\begin{enumerate}
\item What infinitesimal contribution is accumulated?
\item What variable should be integrated over?
\item What are the correct bounds?
\item What units should the final integral have?
\end{enumerate}\end{promptbox}
\textbf{Workspace}\workspace
\begin{reflectionbox}\begin{enumerate}
\item Which sentence or datum in these problems supplied the main mathematical constraint?
\item Which assumptions were essential to the model, and how would the model change if one were removed?
\item When would an analytical solution become inconvenient enough to justify a numerical method?
\end{enumerate}\end{reflectionbox}
\chapter{Volumes by disks/shells and general solids}\label{app:A028}
\textbf{Application ID:} \pid{A028}\quad \textbf{Source locus:} 7.2–7.3; 14\quad \textbf{Source pages:} 30, 130, 157, 454–473, 548, 997
\section*{Application Context}
This family models totals by accumulating infinitesimal contributions. Correct bounds and units are as important as evaluating the antiderivative or quadrature. The canonical application family is \emph{Volumes by disks/shells and general solids}.
\section*{Calculus Toolkit}
\begin{tabularx}{\textwidth}{@{}lX@{}}\toprule
Primary technique & One- and multiple-integral volume \\
Prerequisites & definite integrals, units, geometric/physical modeling \\
Numerical extension & Ch. 4, especially §§4.1–4.8 \\
Analytical/numerical class & NUMERICAL-NATURAL \\
\bottomrule\end{tabularx}
\section*{Modeling Questions to Ask}\begin{itemize}
\item What infinitesimal contribution is accumulated?
\item What variable should be integrated over?
\item What are the correct bounds?
\item What units should the final integral have?
\end{itemize}
\begin{errorbox}Build the infinitesimal contribution before integrating. Check bounds and units of the integrand and final accumulated quantity.\end{errorbox}
\section*{Problems}
\needspace{8\baselineskip}
\subsection*{\pid{APP-028-F-01} — Foundation}
\begin{problembox}
A cumulative quantity in volumes by disks/shells and general solids is modeled by the rate $r(x)=x^2+3$ over $0\le x\le 6$. Find the total accumulated amount.
\end{problembox}
\begin{promptbox}\begin{enumerate}
\item What infinitesimal contribution is accumulated?
\item What variable should be integrated over?
\item What are the correct bounds?
\item What units should the final integral have?
\end{enumerate}\end{promptbox}
\textbf{Workspace}\workspace
\needspace{8\baselineskip}
\subsection*{\pid{APP-028-S-01} — Standard}
\begin{problembox}
In a volumes by disks/shells and general solids application, liquid fills a vertical prismatic tank of height 5 m and horizontal cross-sectional area 3 m$^2$. Use a weight density of 9.8 kN/m$^3$. Find the work required to pump all liquid just to the top edge.
\end{problembox}
\begin{promptbox}\begin{enumerate}
\item What infinitesimal contribution is accumulated?
\item What variable should be integrated over?
\item What are the correct bounds?
\item What units should the final integral have?
\end{enumerate}\end{promptbox}
\textbf{Workspace}\workspace
\needspace{8\baselineskip}
\subsection*{\pid{APP-028-A-01} — Advanced}
\begin{problembox}
A non-elementary accumulation arising in volumes by disks/shells and general solids is $I=\int_0^1 e^{-x^2}\,dx$. Approximate $I$ with composite Simpson's rule using four equal subintervals, and explain why a numerical method is appropriate.
\end{problembox}
\begin{promptbox}\begin{enumerate}
\item What infinitesimal contribution is accumulated?
\item What variable should be integrated over?
\item What are the correct bounds?
\item What units should the final integral have?
\end{enumerate}\end{promptbox}
\textbf{Workspace}\workspace
\begin{reflectionbox}\begin{enumerate}
\item Which sentence or datum in these problems supplied the main mathematical constraint?
\item Which assumptions were essential to the model, and how would the model change if one were removed?
\item When would an analytical solution become inconvenient enough to justify a numerical method?
\end{enumerate}\end{reflectionbox}
\chapter{Arc length, cable length, Gateway Arch, pursuit and stream length}\label{app:A029}
\textbf{Application ID:} \pid{A029}\quad \textbf{Source locus:} 7.4\quad \textbf{Source pages:} 477–484
\section*{Application Context}
This family models totals by accumulating infinitesimal contributions. Correct bounds and units are as important as evaluating the antiderivative or quadrature. The canonical application family is \emph{Arc length, cable length, Gateway Arch, pursuit and stream length}.
\section*{Calculus Toolkit}
\begin{tabularx}{\textwidth}{@{}lX@{}}\toprule
Primary technique & Arc-length integrals \\
Prerequisites & definite integrals, units, geometric/physical modeling \\
Numerical extension & Ch. 4, especially §§4.1–4.8 \\
Analytical/numerical class & NUMERICAL-NATURAL \\
\bottomrule\end{tabularx}
\section*{Modeling Questions to Ask}\begin{itemize}
\item What infinitesimal contribution is accumulated?
\item What variable should be integrated over?
\item What are the correct bounds?
\item What units should the final integral have?
\end{itemize}
\begin{errorbox}Build the infinitesimal contribution before integrating. Check bounds and units of the integrand and final accumulated quantity.\end{errorbox}
\section*{Problems}
\needspace{8\baselineskip}
\subsection*{\pid{APP-029-F-01} — Foundation}
\begin{problembox}
A cumulative quantity in arc length, cable length, gateway arch, pursuit and stream length is modeled by the rate $r(x)=x^2+4$ over $0\le x\le 7$. Find the total accumulated amount.
\end{problembox}
\begin{promptbox}\begin{enumerate}
\item What infinitesimal contribution is accumulated?
\item What variable should be integrated over?
\item What are the correct bounds?
\item What units should the final integral have?
\end{enumerate}\end{promptbox}
\textbf{Workspace}\workspace
\needspace{8\baselineskip}
\subsection*{\pid{APP-029-S-01} — Standard}
\begin{problembox}
In a arc length, cable length, gateway arch, pursuit and stream length application, liquid fills a vertical prismatic tank of height 6 m and horizontal cross-sectional area 4 m$^2$. Use a weight density of 9.8 kN/m$^3$. Find the work required to pump all liquid just to the top edge.
\end{problembox}
\begin{promptbox}\begin{enumerate}
\item What infinitesimal contribution is accumulated?
\item What variable should be integrated over?
\item What are the correct bounds?
\item What units should the final integral have?
\end{enumerate}\end{promptbox}
\textbf{Workspace}\workspace
\needspace{8\baselineskip}
\subsection*{\pid{APP-029-A-01} — Advanced}
\begin{problembox}
A non-elementary accumulation arising in arc length, cable length, gateway arch, pursuit and stream length is $I=\int_0^1 e^{-x^2}\,dx$. Approximate $I$ with composite Simpson's rule using four equal subintervals, and explain why a numerical method is appropriate.
\end{problembox}
\begin{promptbox}\begin{enumerate}
\item What infinitesimal contribution is accumulated?
\item What variable should be integrated over?
\item What are the correct bounds?
\item What units should the final integral have?
\end{enumerate}\end{promptbox}
\textbf{Workspace}\workspace
\begin{reflectionbox}\begin{enumerate}
\item Which sentence or datum in these problems supplied the main mathematical constraint?
\item Which assumptions were essential to the model, and how would the model change if one were removed?
\item When would an analytical solution become inconvenient enough to justify a numerical method?
\end{enumerate}\end{reflectionbox}
\chapter{Surface area and surfaces of revolution}\label{app:A030}
\textbf{Application ID:} \pid{A030}\quad \textbf{Source locus:} 7.4; 14.5\quad \textbf{Source pages:} 157, 162, 241, 474–484, 578, 1006–1012
\section*{Application Context}
This family models totals by accumulating infinitesimal contributions. Correct bounds and units are as important as evaluating the antiderivative or quadrature. The canonical application family is \emph{Surface area and surfaces of revolution}.
\section*{Calculus Toolkit}
\begin{tabularx}{\textwidth}{@{}lX@{}}\toprule
Primary technique & Surface-area integrals \\
Prerequisites & definite integrals, units, geometric/physical modeling \\
Numerical extension & Ch. 4, especially §§4.1–4.8 \\
Analytical/numerical class & NUMERICAL-NATURAL \\
\bottomrule\end{tabularx}
\section*{Modeling Questions to Ask}\begin{itemize}
\item What infinitesimal contribution is accumulated?
\item What variable should be integrated over?
\item What are the correct bounds?
\item What units should the final integral have?
\end{itemize}
\begin{errorbox}Build the infinitesimal contribution before integrating. Check bounds and units of the integrand and final accumulated quantity.\end{errorbox}
\section*{Problems}
\needspace{8\baselineskip}
\subsection*{\pid{APP-030-F-01} — Foundation}
\begin{problembox}
A cumulative quantity in surface area and surfaces of revolution is modeled by the rate $r(x)=x^2+2$ over $0\le x\le 8$. Find the total accumulated amount.
\end{problembox}
\begin{promptbox}\begin{enumerate}
\item What infinitesimal contribution is accumulated?
\item What variable should be integrated over?
\item What are the correct bounds?
\item What units should the final integral have?
\end{enumerate}\end{promptbox}
\textbf{Workspace}\workspace
\needspace{8\baselineskip}
\subsection*{\pid{APP-030-S-01} — Standard}
\begin{problembox}
In a surface area and surfaces of revolution application, liquid fills a vertical prismatic tank of height 7 m and horizontal cross-sectional area 2 m$^2$. Use a weight density of 9.8 kN/m$^3$. Find the work required to pump all liquid just to the top edge.
\end{problembox}
\begin{promptbox}\begin{enumerate}
\item What infinitesimal contribution is accumulated?
\item What variable should be integrated over?
\item What are the correct bounds?
\item What units should the final integral have?
\end{enumerate}\end{promptbox}
\textbf{Workspace}\workspace
\needspace{8\baselineskip}
\subsection*{\pid{APP-030-A-01} — Advanced}
\begin{problembox}
A non-elementary accumulation arising in surface area and surfaces of revolution is $I=\int_0^1 e^{-x^2}\,dx$. Approximate $I$ with composite Simpson's rule using four equal subintervals, and explain why a numerical method is appropriate.
\end{problembox}
\begin{promptbox}\begin{enumerate}
\item What infinitesimal contribution is accumulated?
\item What variable should be integrated over?
\item What are the correct bounds?
\item What units should the final integral have?
\end{enumerate}\end{promptbox}
\textbf{Workspace}\workspace
\begin{reflectionbox}\begin{enumerate}
\item Which sentence or datum in these problems supplied the main mathematical constraint?
\item Which assumptions were essential to the model, and how would the model change if one were removed?
\item When would an analytical solution become inconvenient enough to justify a numerical method?
\end{enumerate}\end{reflectionbox}
\chapter{Work: lifting, pumping, towing, springs, gases, rockets and force fields}\label{app:A031}
\textbf{Application ID:} \pid{A031}\quad \textbf{Source locus:} 7.5; 15\quad \textbf{Source pages:} 485–493, 512, 565, 580, 776, 815, 1066–1128
\section*{Application Context}
This family models totals by accumulating infinitesimal contributions. Correct bounds and units are as important as evaluating the antiderivative or quadrature. The canonical application family is \emph{Work: lifting, pumping, towing, springs, gases, rockets and force fields}.
\section*{Calculus Toolkit}
\begin{tabularx}{\textwidth}{@{}lX@{}}\toprule
Primary technique & Work integrals and line integrals \\
Prerequisites & definite integrals, units, geometric/physical modeling \\
Numerical extension & Ch. 4, especially §§4.1–4.8 \\
Analytical/numerical class & NUMERICAL-NATURAL \\
\bottomrule\end{tabularx}
\section*{Modeling Questions to Ask}\begin{itemize}
\item What infinitesimal contribution is accumulated?
\item What variable should be integrated over?
\item What are the correct bounds?
\item What units should the final integral have?
\end{itemize}
\begin{errorbox}Build the infinitesimal contribution before integrating. Check bounds and units of the integrand and final accumulated quantity.\end{errorbox}
\section*{Problems}
\needspace{8\baselineskip}
\subsection*{\pid{APP-031-F-01} — Foundation}
\begin{problembox}
A cumulative quantity in work: lifting, pumping, towing, springs, gases, rockets and force fields is modeled by the rate $r(x)=x^2+3$ over $0\le x\le 9$. Find the total accumulated amount.
\end{problembox}
\begin{promptbox}\begin{enumerate}
\item What infinitesimal contribution is accumulated?
\item What variable should be integrated over?
\item What are the correct bounds?
\item What units should the final integral have?
\end{enumerate}\end{promptbox}
\textbf{Workspace}\workspace
\needspace{8\baselineskip}
\subsection*{\pid{APP-031-S-01} — Standard}
\begin{problembox}
In a work: lifting, pumping, towing, springs, gases, rockets and force fields application, liquid fills a vertical prismatic tank of height 8 m and horizontal cross-sectional area 3 m$^2$. Use a weight density of 9.8 kN/m$^3$. Find the work required to pump all liquid just to the top edge.
\end{problembox}
\begin{promptbox}\begin{enumerate}
\item What infinitesimal contribution is accumulated?
\item What variable should be integrated over?
\item What are the correct bounds?
\item What units should the final integral have?
\end{enumerate}\end{promptbox}
\textbf{Workspace}\workspace
\needspace{8\baselineskip}
\subsection*{\pid{APP-031-A-01} — Advanced}
\begin{problembox}
A non-elementary accumulation arising in work: lifting, pumping, towing, springs, gases, rockets and force fields is $I=\int_0^1 e^{-x^2}\,dx$. Approximate $I$ with composite Simpson's rule using four equal subintervals, and explain why a numerical method is appropriate.
\end{problembox}
\begin{promptbox}\begin{enumerate}
\item What infinitesimal contribution is accumulated?
\item What variable should be integrated over?
\item What are the correct bounds?
\item What units should the final integral have?
\end{enumerate}\end{promptbox}
\textbf{Workspace}\workspace
\begin{reflectionbox}\begin{enumerate}
\item Which sentence or datum in these problems supplied the main mathematical constraint?
\item Which assumptions were essential to the model, and how would the model change if one were removed?
\item When would an analytical solution become inconvenient enough to justify a numerical method?
\end{enumerate}\end{reflectionbox}
\chapter{Fluid pressure and fluid force}\label{app:A032}
\textbf{Application ID:} \pid{A032}\quad \textbf{Source locus:} 7.7\quad \textbf{Source pages:} 506–514, 546, 549
\section*{Application Context}
This family models totals by accumulating infinitesimal contributions. Correct bounds and units are as important as evaluating the antiderivative or quadrature. The canonical application family is \emph{Fluid pressure and fluid force}.
\section*{Calculus Toolkit}
\begin{tabularx}{\textwidth}{@{}lX@{}}\toprule
Primary technique & Pressure-force integration \\
Prerequisites & definite integrals, units, geometric/physical modeling \\
Numerical extension & Ch. 4, especially §§4.1–4.8 \\
Analytical/numerical class & NUMERICAL-NATURAL \\
\bottomrule\end{tabularx}
\section*{Modeling Questions to Ask}\begin{itemize}
\item What infinitesimal contribution is accumulated?
\item What variable should be integrated over?
\item What are the correct bounds?
\item What units should the final integral have?
\end{itemize}
\begin{errorbox}Build the infinitesimal contribution before integrating. Check bounds and units of the integrand and final accumulated quantity.\end{errorbox}
\section*{Problems}
\needspace{8\baselineskip}
\subsection*{\pid{APP-032-F-01} — Foundation}
\begin{problembox}
A cumulative quantity in fluid pressure and fluid force is modeled by the rate $r(x)=x^2+4$ over $0\le x\le 6$. Find the total accumulated amount.
\end{problembox}
\begin{promptbox}\begin{enumerate}
\item What infinitesimal contribution is accumulated?
\item What variable should be integrated over?
\item What are the correct bounds?
\item What units should the final integral have?
\end{enumerate}\end{promptbox}
\textbf{Workspace}\workspace
\needspace{8\baselineskip}
\subsection*{\pid{APP-032-S-01} — Standard}
\begin{problembox}
In a fluid pressure and fluid force application, liquid fills a vertical prismatic tank of height 5 m and horizontal cross-sectional area 4 m$^2$. Use a weight density of 9.8 kN/m$^3$. Find the work required to pump all liquid just to the top edge.
\end{problembox}
\begin{promptbox}\begin{enumerate}
\item What infinitesimal contribution is accumulated?
\item What variable should be integrated over?
\item What are the correct bounds?
\item What units should the final integral have?
\end{enumerate}\end{promptbox}
\textbf{Workspace}\workspace
\needspace{8\baselineskip}
\subsection*{\pid{APP-032-A-01} — Advanced}
\begin{problembox}
A non-elementary accumulation arising in fluid pressure and fluid force is $I=\int_0^1 e^{-x^2}\,dx$. Approximate $I$ with composite Simpson's rule using four equal subintervals, and explain why a numerical method is appropriate.
\end{problembox}
\begin{promptbox}\begin{enumerate}
\item What infinitesimal contribution is accumulated?
\item What variable should be integrated over?
\item What are the correct bounds?
\item What units should the final integral have?
\end{enumerate}\end{promptbox}
\textbf{Workspace}\workspace
\begin{reflectionbox}\begin{enumerate}
\item Which sentence or datum in these problems supplied the main mathematical constraint?
\item Which assumptions were essential to the model, and how would the model change if one were removed?
\item When would an analytical solution become inconvenient enough to justify a numerical method?
\end{enumerate}\end{reflectionbox}
\chapter{Center of mass, centroid, moments and inertia}\label{app:A033}
\textbf{Application ID:} \pid{A033}\quad \textbf{Source locus:} 7.6; 14.4\quad \textbf{Source pages:} 494–504, 527, 998–1005
\section*{Application Context}
This family models totals by accumulating infinitesimal contributions. Correct bounds and units are as important as evaluating the antiderivative or quadrature. The canonical application family is \emph{Center of mass, centroid, moments and inertia}.
\section*{Calculus Toolkit}
\begin{tabularx}{\textwidth}{@{}lX@{}}\toprule
Primary technique & Moment integrals \\
Prerequisites & definite integrals, units, geometric/physical modeling \\
Numerical extension & Ch. 4, especially §§4.1–4.8 \\
Analytical/numerical class & NUMERICAL-NATURAL \\
\bottomrule\end{tabularx}
\section*{Modeling Questions to Ask}\begin{itemize}
\item What infinitesimal contribution is accumulated?
\item What variable should be integrated over?
\item What are the correct bounds?
\item What units should the final integral have?
\end{itemize}
\begin{errorbox}Build the infinitesimal contribution before integrating. Check bounds and units of the integrand and final accumulated quantity.\end{errorbox}
\section*{Problems}
\needspace{8\baselineskip}
\subsection*{\pid{APP-033-F-01} — Foundation}
\begin{problembox}
A cumulative quantity in center of mass, centroid, moments and inertia is modeled by the rate $r(x)=x^2+2$ over $0\le x\le 7$. Find the total accumulated amount.
\end{problembox}
\begin{promptbox}\begin{enumerate}
\item What infinitesimal contribution is accumulated?
\item What variable should be integrated over?
\item What are the correct bounds?
\item What units should the final integral have?
\end{enumerate}\end{promptbox}
\textbf{Workspace}\workspace
\needspace{8\baselineskip}
\subsection*{\pid{APP-033-S-01} — Standard}
\begin{problembox}
In a center of mass, centroid, moments and inertia application, liquid fills a vertical prismatic tank of height 6 m and horizontal cross-sectional area 2 m$^2$. Use a weight density of 9.8 kN/m$^3$. Find the work required to pump all liquid just to the top edge.
\end{problembox}
\begin{promptbox}\begin{enumerate}
\item What infinitesimal contribution is accumulated?
\item What variable should be integrated over?
\item What are the correct bounds?
\item What units should the final integral have?
\end{enumerate}\end{promptbox}
\textbf{Workspace}\workspace
\needspace{8\baselineskip}
\subsection*{\pid{APP-033-A-01} — Advanced}
\begin{problembox}
A non-elementary accumulation arising in center of mass, centroid, moments and inertia is $I=\int_0^1 e^{-x^2}\,dx$. Approximate $I$ with composite Simpson's rule using four equal subintervals, and explain why a numerical method is appropriate.
\end{problembox}
\begin{promptbox}\begin{enumerate}
\item What infinitesimal contribution is accumulated?
\item What variable should be integrated over?
\item What are the correct bounds?
\item What units should the final integral have?
\end{enumerate}\end{promptbox}
\textbf{Workspace}\workspace
\begin{reflectionbox}\begin{enumerate}
\item Which sentence or datum in these problems supplied the main mathematical constraint?
\item Which assumptions were essential to the model, and how would the model change if one were removed?
\item When would an analytical solution become inconvenient enough to justify a numerical method?
\end{enumerate}\end{reflectionbox}
\chapter{Consumer and producer surplus}\label{app:A034}
\textbf{Application ID:} \pid{A034}\quad \textbf{Source locus:} 7.7 / applied integration\quad \textbf{Source pages:} 514
\section*{Application Context}
This family models totals by accumulating infinitesimal contributions. Correct bounds and units are as important as evaluating the antiderivative or quadrature. The canonical application family is \emph{Consumer and producer surplus}.
\section*{Calculus Toolkit}
\begin{tabularx}{\textwidth}{@{}lX@{}}\toprule
Primary technique & Area between economic curves \\
Prerequisites & definite integrals, units, geometric/physical modeling \\
Numerical extension & Ch. 4, especially §§4.1–4.8 \\
Analytical/numerical class & NUMERICAL-NATURAL \\
\bottomrule\end{tabularx}
\section*{Modeling Questions to Ask}\begin{itemize}
\item What infinitesimal contribution is accumulated?
\item What variable should be integrated over?
\item What are the correct bounds?
\item What units should the final integral have?
\end{itemize}
\begin{errorbox}Build the infinitesimal contribution before integrating. Check bounds and units of the integrand and final accumulated quantity.\end{errorbox}
\section*{Problems}
\needspace{8\baselineskip}
\subsection*{\pid{APP-034-F-01} — Foundation}
\begin{problembox}
A cumulative quantity in consumer and producer surplus is modeled by the rate $r(x)=x^2+3$ over $0\le x\le 8$. Find the total accumulated amount.
\end{problembox}
\begin{promptbox}\begin{enumerate}
\item What infinitesimal contribution is accumulated?
\item What variable should be integrated over?
\item What are the correct bounds?
\item What units should the final integral have?
\end{enumerate}\end{promptbox}
\textbf{Workspace}\workspace
\needspace{8\baselineskip}
\subsection*{\pid{APP-034-S-01} — Standard}
\begin{problembox}
In a consumer and producer surplus application, liquid fills a vertical prismatic tank of height 7 m and horizontal cross-sectional area 3 m$^2$. Use a weight density of 9.8 kN/m$^3$. Find the work required to pump all liquid just to the top edge.
\end{problembox}
\begin{promptbox}\begin{enumerate}
\item What infinitesimal contribution is accumulated?
\item What variable should be integrated over?
\item What are the correct bounds?
\item What units should the final integral have?
\end{enumerate}\end{promptbox}
\textbf{Workspace}\workspace
\needspace{8\baselineskip}
\subsection*{\pid{APP-034-A-01} — Advanced}
\begin{problembox}
A non-elementary accumulation arising in consumer and producer surplus is $I=\int_0^1 e^{-x^2}\,dx$. Approximate $I$ with composite Simpson's rule using four equal subintervals, and explain why a numerical method is appropriate.
\end{problembox}
\begin{promptbox}\begin{enumerate}
\item What infinitesimal contribution is accumulated?
\item What variable should be integrated over?
\item What are the correct bounds?
\item What units should the final integral have?
\end{enumerate}\end{promptbox}
\textbf{Workspace}\workspace
\begin{reflectionbox}\begin{enumerate}
\item Which sentence or datum in these problems supplied the main mathematical constraint?
\item Which assumptions were essential to the model, and how would the model change if one were removed?
\item When would an analytical solution become inconvenient enough to justify a numerical method?
\end{enumerate}\end{reflectionbox}
\chapter{Average temperature, production and profit}\label{app:A035}
\textbf{Application ID:} \pid{A035}\quad \textbf{Source locus:} 14\quad \textbf{Source pages:} 988, 1038
\section*{Application Context}
This family models totals by accumulating infinitesimal contributions. Correct bounds and units are as important as evaluating the antiderivative or quadrature. The canonical application family is \emph{Average temperature, production and profit}.
\section*{Calculus Toolkit}
\begin{tabularx}{\textwidth}{@{}lX@{}}\toprule
Primary technique & Average-value/multiple integration \\
Prerequisites & definite integrals, units, geometric/physical modeling \\
Numerical extension & Ch. 4, especially §§4.1–4.8 \\
Analytical/numerical class & NUMERICAL-NATURAL \\
\bottomrule\end{tabularx}
\section*{Modeling Questions to Ask}\begin{itemize}
\item What infinitesimal contribution is accumulated?
\item What variable should be integrated over?
\item What are the correct bounds?
\item What units should the final integral have?
\end{itemize}
\begin{errorbox}Build the infinitesimal contribution before integrating. Check bounds and units of the integrand and final accumulated quantity.\end{errorbox}
\section*{Problems}
\needspace{8\baselineskip}
\subsection*{\pid{APP-035-F-01} — Foundation}
\begin{problembox}
A cumulative quantity in average temperature, production and profit is modeled by the rate $r(x)=x^2+4$ over $0\le x\le 9$. Find the total accumulated amount.
\end{problembox}
\begin{promptbox}\begin{enumerate}
\item What infinitesimal contribution is accumulated?
\item What variable should be integrated over?
\item What are the correct bounds?
\item What units should the final integral have?
\end{enumerate}\end{promptbox}
\textbf{Workspace}\workspace
\needspace{8\baselineskip}
\subsection*{\pid{APP-035-S-01} — Standard}
\begin{problembox}
In a average temperature, production and profit application, liquid fills a vertical prismatic tank of height 8 m and horizontal cross-sectional area 4 m$^2$. Use a weight density of 9.8 kN/m$^3$. Find the work required to pump all liquid just to the top edge.
\end{problembox}
\begin{promptbox}\begin{enumerate}
\item What infinitesimal contribution is accumulated?
\item What variable should be integrated over?
\item What are the correct bounds?
\item What units should the final integral have?
\end{enumerate}\end{promptbox}
\textbf{Workspace}\workspace
\needspace{8\baselineskip}
\subsection*{\pid{APP-035-A-01} — Advanced}
\begin{problembox}
A non-elementary accumulation arising in average temperature, production and profit is $I=\int_0^1 e^{-x^2}\,dx$. Approximate $I$ with composite Simpson's rule using four equal subintervals, and explain why a numerical method is appropriate.
\end{problembox}
\begin{promptbox}\begin{enumerate}
\item What infinitesimal contribution is accumulated?
\item What variable should be integrated over?
\item What are the correct bounds?
\item What units should the final integral have?
\end{enumerate}\end{promptbox}
\textbf{Workspace}\workspace
\begin{reflectionbox}\begin{enumerate}
\item Which sentence or datum in these problems supplied the main mathematical constraint?
\item Which assumptions were essential to the model, and how would the model change if one were removed?
\item When would an analytical solution become inconvenient enough to justify a numerical method?
\end{enumerate}\end{reflectionbox}
\chapter{Probability integrals and expectation-like calculations}\label{app:A036}
\textbf{Application ID:} \pid{A036}\quad \textbf{Source locus:} 4–9; 14\quad \textbf{Source pages:} 308, 351, 581, 584, 606, 607, 667, 678, 989, 997
\section*{Application Context}
This family models totals by accumulating infinitesimal contributions. Correct bounds and units are as important as evaluating the antiderivative or quadrature. The canonical application family is \emph{Probability integrals and expectation-like calculations}.
\section*{Calculus Toolkit}
\begin{tabularx}{\textwidth}{@{}lX@{}}\toprule
Primary technique & Integration of probability models \\
Prerequisites & definite integrals, units, geometric/physical modeling \\
Numerical extension & Ch. 4, especially §§4.1–4.8 \\
Analytical/numerical class & NUMERICAL-NATURAL \\
\bottomrule\end{tabularx}
\section*{Modeling Questions to Ask}\begin{itemize}
\item What infinitesimal contribution is accumulated?
\item What variable should be integrated over?
\item What are the correct bounds?
\item What units should the final integral have?
\end{itemize}
\begin{errorbox}Build the infinitesimal contribution before integrating. Check bounds and units of the integrand and final accumulated quantity.\end{errorbox}
\section*{Problems}
\needspace{8\baselineskip}
\subsection*{\pid{APP-036-F-01} — Foundation}
\begin{problembox}
A cumulative quantity in probability integrals and expectation-like calculations is modeled by the rate $r(x)=x^2+2$ over $0\le x\le 6$. Find the total accumulated amount.
\end{problembox}
\begin{promptbox}\begin{enumerate}
\item What infinitesimal contribution is accumulated?
\item What variable should be integrated over?
\item What are the correct bounds?
\item What units should the final integral have?
\end{enumerate}\end{promptbox}
\textbf{Workspace}\workspace
\needspace{8\baselineskip}
\subsection*{\pid{APP-036-S-01} — Standard}
\begin{problembox}
In a probability integrals and expectation-like calculations application, liquid fills a vertical prismatic tank of height 5 m and horizontal cross-sectional area 2 m$^2$. Use a weight density of 9.8 kN/m$^3$. Find the work required to pump all liquid just to the top edge.
\end{problembox}
\begin{promptbox}\begin{enumerate}
\item What infinitesimal contribution is accumulated?
\item What variable should be integrated over?
\item What are the correct bounds?
\item What units should the final integral have?
\end{enumerate}\end{promptbox}
\textbf{Workspace}\workspace
\needspace{8\baselineskip}
\subsection*{\pid{APP-036-A-01} — Advanced}
\begin{problembox}
A non-elementary accumulation arising in probability integrals and expectation-like calculations is $I=\int_0^1 e^{-x^2}\,dx$. Approximate $I$ with composite Simpson's rule using four equal subintervals, and explain why a numerical method is appropriate.
\end{problembox}
\begin{promptbox}\begin{enumerate}
\item What infinitesimal contribution is accumulated?
\item What variable should be integrated over?
\item What are the correct bounds?
\item What units should the final integral have?
\end{enumerate}\end{promptbox}
\textbf{Workspace}\workspace
\begin{reflectionbox}\begin{enumerate}
\item Which sentence or datum in these problems supplied the main mathematical constraint?
\item Which assumptions were essential to the model, and how would the model change if one were removed?
\item When would an analytical solution become inconvenient enough to justify a numerical method?
\end{enumerate}\end{reflectionbox}
\part{Approximation, Series, and Special Models}
\chapter{Topographic maps, terrain and scattered spatial data}\label{app:A037}
\textbf{Application ID:} \pid{A037}\quad \textbf{Source locus:} 2; 13\quad \textbf{Source pages:} 150, 875, 929, 930
\section*{Application Context}
This family uses interpolation, series, transformations, and approximation to evaluate functions or represent data when direct formulas are inconvenient. The canonical application family is \emph{Topographic maps, terrain and scattered spatial data}.
\section*{Calculus Toolkit}
\begin{tabularx}{\textwidth}{@{}lX@{}}\toprule
Primary technique & Interpolation, gradients and surfaces \\
Prerequisites & functions, series or interpolation, approximation error \\
Numerical extension & Ch. 3 (interpolation), Ch. 5 (function evaluation), Ch. 6 (special functions) \\
Analytical/numerical class & NUMERICAL-NATURAL \\
\bottomrule\end{tabularx}
\section*{Modeling Questions to Ask}\begin{itemize}
\item What data or series is being approximated?
\item Is the evaluation interpolation or extrapolation?
\item What smoothness or convergence assumptions are used?
\item How will approximation quality be checked?
\end{itemize}
\begin{errorbox}Interpolation is not unrestricted extrapolation. State the interval on which the approximation is intended to be used and check approximation error when possible.\end{errorbox}
\section*{Problems}
\needspace{8\baselineskip}
\subsection*{\pid{APP-037-F-01} — Foundation}
\begin{problembox}
Two measurements in topographic maps, terrain and scattered spatial data are $(0,22)$ and $(10,37)$. Estimate the value at $x=4$ by linear interpolation.
\end{problembox}
\begin{promptbox}\begin{enumerate}
\item What data or series is being approximated?
\item Is the evaluation interpolation or extrapolation?
\item What smoothness or convergence assumptions are used?
\item How will approximation quality be checked?
\end{enumerate}\end{promptbox}
\textbf{Workspace}\workspace
\needspace{8\baselineskip}
\subsection*{\pid{APP-037-S-01} — Standard}
\begin{problembox}
A repeated contribution in topographic maps, terrain and scattered spatial data has magnitudes $285, 185.2, 120.4,\ldots$, with common ratio $r=0.65$. Determine the infinite accumulated total and state the convergence condition used.
\end{problembox}
\begin{promptbox}\begin{enumerate}
\item What data or series is being approximated?
\item Is the evaluation interpolation or extrapolation?
\item What smoothness or convergence assumptions are used?
\item How will approximation quality be checked?
\end{enumerate}\end{promptbox}
\textbf{Workspace}\workspace
\needspace{8\baselineskip}
\subsection*{\pid{APP-037-A-01} — Advanced}
\begin{problembox}
A local approximation useful in topographic maps, terrain and scattered spatial data requires $\ln(1+x)$ at $x=0.2$. Use the cubic Taylor polynomial $x-x^2/2+x^3/3$, compare with the exact value, and report the absolute error.
\end{problembox}
\begin{promptbox}\begin{enumerate}
\item What data or series is being approximated?
\item Is the evaluation interpolation or extrapolation?
\item What smoothness or convergence assumptions are used?
\item How will approximation quality be checked?
\end{enumerate}\end{promptbox}
\textbf{Workspace}\workspace
\begin{reflectionbox}\begin{enumerate}
\item Which sentence or datum in these problems supplied the main mathematical constraint?
\item Which assumptions were essential to the model, and how would the model change if one were removed?
\item When would an analytical solution become inconvenient enough to justify a numerical method?
\end{enumerate}\end{reflectionbox}
\chapter{Mercator map and coordinate-transform applications}\label{app:A038}
\textbf{Application ID:} \pid{A038}\quad \textbf{Source locus:} 5.9\quad \textbf{Source pages:} 399
\section*{Application Context}
This family uses interpolation, series, transformations, and approximation to evaluate functions or represent data when direct formulas are inconvenient. The canonical application family is \emph{Mercator map and coordinate-transform applications}.
\section*{Calculus Toolkit}
\begin{tabularx}{\textwidth}{@{}lX@{}}\toprule
Primary technique & Transcendental functions and integration \\
Prerequisites & functions, series or interpolation, approximation error \\
Numerical extension & Ch. 3 (interpolation), Ch. 5 (function evaluation), Ch. 6 (special functions) \\
Analytical/numerical class & NUMERICAL-OPTIONAL \\
\bottomrule\end{tabularx}
\section*{Modeling Questions to Ask}\begin{itemize}
\item What data or series is being approximated?
\item Is the evaluation interpolation or extrapolation?
\item What smoothness or convergence assumptions are used?
\item How will approximation quality be checked?
\end{itemize}
\begin{errorbox}Interpolation is not unrestricted extrapolation. State the interval on which the approximation is intended to be used and check approximation error when possible.\end{errorbox}
\section*{Problems}
\needspace{8\baselineskip}
\subsection*{\pid{APP-038-F-01} — Foundation}
\begin{problembox}
Two measurements in mercator map and coordinate-transform applications are $(0,23)$ and $(10,38)$. Estimate the value at $x=4$ by linear interpolation.
\end{problembox}
\begin{promptbox}\begin{enumerate}
\item What data or series is being approximated?
\item Is the evaluation interpolation or extrapolation?
\item What smoothness or convergence assumptions are used?
\item How will approximation quality be checked?
\end{enumerate}\end{promptbox}
\textbf{Workspace}\workspace
\needspace{8\baselineskip}
\subsection*{\pid{APP-038-S-01} — Standard}
\begin{problembox}
A repeated contribution in mercator map and coordinate-transform applications has magnitudes $290, 203.0, 142.1,\ldots$, with common ratio $r=0.70$. Determine the infinite accumulated total and state the convergence condition used.
\end{problembox}
\begin{promptbox}\begin{enumerate}
\item What data or series is being approximated?
\item Is the evaluation interpolation or extrapolation?
\item What smoothness or convergence assumptions are used?
\item How will approximation quality be checked?
\end{enumerate}\end{promptbox}
\textbf{Workspace}\workspace
\needspace{8\baselineskip}
\subsection*{\pid{APP-038-A-01} — Advanced}
\begin{problembox}
A local approximation useful in mercator map and coordinate-transform applications requires $\ln(1+x)$ at $x=0.2$. Use the cubic Taylor polynomial $x-x^2/2+x^3/3$, compare with the exact value, and report the absolute error.
\end{problembox}
\begin{promptbox}\begin{enumerate}
\item What data or series is being approximated?
\item Is the evaluation interpolation or extrapolation?
\item What smoothness or convergence assumptions are used?
\item How will approximation quality be checked?
\end{enumerate}\end{promptbox}
\textbf{Workspace}\workspace
\begin{reflectionbox}\begin{enumerate}
\item Which sentence or datum in these problems supplied the main mathematical constraint?
\item Which assumptions were essential to the model, and how would the model change if one were removed?
\item When would an analytical solution become inconvenient enough to justify a numerical method?
\end{enumerate}\end{reflectionbox}
\chapter{Tractrix and implicit/transcendental curve models}\label{app:A039}
\textbf{Application ID:} \pid{A039}\quad \textbf{Source locus:} 5; 8; 10\quad \textbf{Source pages:} 323, 395, 399, 549, 718
\section*{Application Context}
This family uses interpolation, series, transformations, and approximation to evaluate functions or represent data when direct formulas are inconvenient. The canonical application family is \emph{Tractrix and implicit/transcendental curve models}.
\section*{Calculus Toolkit}
\begin{tabularx}{\textwidth}{@{}lX@{}}\toprule
Primary technique & Implicit curves, special functions and quadrature \\
Prerequisites & functions, series or interpolation, approximation error \\
Numerical extension & Ch. 3 (interpolation), Ch. 5 (function evaluation), Ch. 6 (special functions) \\
Analytical/numerical class & NUMERICAL-NATURAL \\
\bottomrule\end{tabularx}
\section*{Modeling Questions to Ask}\begin{itemize}
\item What data or series is being approximated?
\item Is the evaluation interpolation or extrapolation?
\item What smoothness or convergence assumptions are used?
\item How will approximation quality be checked?
\end{itemize}
\begin{errorbox}Interpolation is not unrestricted extrapolation. State the interval on which the approximation is intended to be used and check approximation error when possible.\end{errorbox}
\section*{Problems}
\needspace{8\baselineskip}
\subsection*{\pid{APP-039-F-01} — Foundation}
\begin{problembox}
Two measurements in tractrix and implicit/transcendental curve models are $(0,24)$ and $(10,39)$. Estimate the value at $x=4$ by linear interpolation.
\end{problembox}
\begin{promptbox}\begin{enumerate}
\item What data or series is being approximated?
\item Is the evaluation interpolation or extrapolation?
\item What smoothness or convergence assumptions are used?
\item How will approximation quality be checked?
\end{enumerate}\end{promptbox}
\textbf{Workspace}\workspace
\needspace{8\baselineskip}
\subsection*{\pid{APP-039-S-01} — Standard}
\begin{problembox}
A repeated contribution in tractrix and implicit/transcendental curve models has magnitudes $295, 177.0, 106.2,\ldots$, with common ratio $r=0.60$. Determine the infinite accumulated total and state the convergence condition used.
\end{problembox}
\begin{promptbox}\begin{enumerate}
\item What data or series is being approximated?
\item Is the evaluation interpolation or extrapolation?
\item What smoothness or convergence assumptions are used?
\item How will approximation quality be checked?
\end{enumerate}\end{promptbox}
\textbf{Workspace}\workspace
\needspace{8\baselineskip}
\subsection*{\pid{APP-039-A-01} — Advanced}
\begin{problembox}
A local approximation useful in tractrix and implicit/transcendental curve models requires $\ln(1+x)$ at $x=0.2$. Use the cubic Taylor polynomial $x-x^2/2+x^3/3$, compare with the exact value, and report the absolute error.
\end{problembox}
\begin{promptbox}\begin{enumerate}
\item What data or series is being approximated?
\item Is the evaluation interpolation or extrapolation?
\item What smoothness or convergence assumptions are used?
\item How will approximation quality be checked?
\end{enumerate}\end{promptbox}
\textbf{Workspace}\workspace
\begin{reflectionbox}\begin{enumerate}
\item Which sentence or datum in these problems supplied the main mathematical constraint?
\item Which assumptions were essential to the model, and how would the model change if one were removed?
\item When would an analytical solution become inconvenient enough to justify a numerical method?
\end{enumerate}\end{reflectionbox}
\chapter{Sound intensity, decibel-like and exponential/logarithmic models}\label{app:A040}
\textbf{Application ID:} \pid{A040}\quad \textbf{Source locus:} 1; 5; 6\quad \textbf{Source pages:} 44, 323, 422
\section*{Application Context}
This family uses interpolation, series, transformations, and approximation to evaluate functions or represent data when direct formulas are inconvenient. The canonical application family is \emph{Sound intensity, decibel-like and exponential/logarithmic models}.
\section*{Calculus Toolkit}
\begin{tabularx}{\textwidth}{@{}lX@{}}\toprule
Primary technique & Logarithmic/exponential functions \\
Prerequisites & functions, series or interpolation, approximation error \\
Numerical extension & Ch. 3 (interpolation), Ch. 5 (function evaluation), Ch. 6 (special functions) \\
Analytical/numerical class & NUMERICAL-OPTIONAL \\
\bottomrule\end{tabularx}
\section*{Modeling Questions to Ask}\begin{itemize}
\item What data or series is being approximated?
\item Is the evaluation interpolation or extrapolation?
\item What smoothness or convergence assumptions are used?
\item How will approximation quality be checked?
\end{itemize}
\begin{errorbox}Interpolation is not unrestricted extrapolation. State the interval on which the approximation is intended to be used and check approximation error when possible.\end{errorbox}
\section*{Problems}
\needspace{8\baselineskip}
\subsection*{\pid{APP-040-F-01} — Foundation}
\begin{problembox}
Two measurements in sound intensity, decibel-like and exponential/logarithmic models are $(0,25)$ and $(10,40)$. Estimate the value at $x=4$ by linear interpolation.
\end{problembox}
\begin{promptbox}\begin{enumerate}
\item What data or series is being approximated?
\item Is the evaluation interpolation or extrapolation?
\item What smoothness or convergence assumptions are used?
\item How will approximation quality be checked?
\end{enumerate}\end{promptbox}
\textbf{Workspace}\workspace
\needspace{8\baselineskip}
\subsection*{\pid{APP-040-S-01} — Standard}
\begin{problembox}
A repeated contribution in sound intensity, decibel-like and exponential/logarithmic models has magnitudes $300, 195.0, 126.8,\ldots$, with common ratio $r=0.65$. Determine the infinite accumulated total and state the convergence condition used.
\end{problembox}
\begin{promptbox}\begin{enumerate}
\item What data or series is being approximated?
\item Is the evaluation interpolation or extrapolation?
\item What smoothness or convergence assumptions are used?
\item How will approximation quality be checked?
\end{enumerate}\end{promptbox}
\textbf{Workspace}\workspace
\needspace{8\baselineskip}
\subsection*{\pid{APP-040-A-01} — Advanced}
\begin{problembox}
A local approximation useful in sound intensity, decibel-like and exponential/logarithmic models requires $\ln(1+x)$ at $x=0.2$. Use the cubic Taylor polynomial $x-x^2/2+x^3/3$, compare with the exact value, and report the absolute error.
\end{problembox}
\begin{promptbox}\begin{enumerate}
\item What data or series is being approximated?
\item Is the evaluation interpolation or extrapolation?
\item What smoothness or convergence assumptions are used?
\item How will approximation quality be checked?
\end{enumerate}\end{promptbox}
\textbf{Workspace}\workspace
\begin{reflectionbox}\begin{enumerate}
\item Which sentence or datum in these problems supplied the main mathematical constraint?
\item Which assumptions were essential to the model, and how would the model change if one were removed?
\item When would an analytical solution become inconvenient enough to justify a numerical method?
\end{enumerate}\end{reflectionbox}
\chapter{Compound interest, annuities, home mortgage, present value and capitalized cost}\label{app:A041}
\textbf{Application ID:} \pid{A041}\quad \textbf{Source locus:} 5; 6; 9\quad \textbf{Source pages:} 323, 356–360, 371, 401, 421, 440, 531, 581, 584, 597, 607, 680
\section*{Application Context}
This family uses interpolation, series, transformations, and approximation to evaluate functions or represent data when direct formulas are inconvenient. The canonical application family is \emph{Compound interest, annuities, home mortgage, present value and capitalized cost}.
\section*{Calculus Toolkit}
\begin{tabularx}{\textwidth}{@{}lX@{}}\toprule
Primary technique & Exponential models, series and accumulation \\
Prerequisites & functions, series or interpolation, approximation error \\
Numerical extension & Ch. 3 (interpolation), Ch. 5 (function evaluation), Ch. 6 (special functions) \\
Analytical/numerical class & NUMERICAL-NATURAL \\
\bottomrule\end{tabularx}
\section*{Modeling Questions to Ask}\begin{itemize}
\item What data or series is being approximated?
\item Is the evaluation interpolation or extrapolation?
\item What smoothness or convergence assumptions are used?
\item How will approximation quality be checked?
\end{itemize}
\begin{errorbox}Interpolation is not unrestricted extrapolation. State the interval on which the approximation is intended to be used and check approximation error when possible.\end{errorbox}
\section*{Problems}
\needspace{8\baselineskip}
\subsection*{\pid{APP-041-F-01} — Foundation}
\begin{problembox}
Two measurements in compound interest, annuities, home mortgage, present value and capitalized cost are $(0,26)$ and $(10,41)$. Estimate the value at $x=4$ by linear interpolation.
\end{problembox}
\begin{promptbox}\begin{enumerate}
\item What data or series is being approximated?
\item Is the evaluation interpolation or extrapolation?
\item What smoothness or convergence assumptions are used?
\item How will approximation quality be checked?
\end{enumerate}\end{promptbox}
\textbf{Workspace}\workspace
\needspace{8\baselineskip}
\subsection*{\pid{APP-041-S-01} — Standard}
\begin{problembox}
A repeated contribution in compound interest, annuities, home mortgage, present value and capitalized cost has magnitudes $305, 213.5, 149.4,\ldots$, with common ratio $r=0.70$. Determine the infinite accumulated total and state the convergence condition used.
\end{problembox}
\begin{promptbox}\begin{enumerate}
\item What data or series is being approximated?
\item Is the evaluation interpolation or extrapolation?
\item What smoothness or convergence assumptions are used?
\item How will approximation quality be checked?
\end{enumerate}\end{promptbox}
\textbf{Workspace}\workspace
\needspace{8\baselineskip}
\subsection*{\pid{APP-041-A-01} — Advanced}
\begin{problembox}
A local approximation useful in compound interest, annuities, home mortgage, present value and capitalized cost requires $\ln(1+x)$ at $x=0.2$. Use the cubic Taylor polynomial $x-x^2/2+x^3/3$, compare with the exact value, and report the absolute error.
\end{problembox}
\begin{promptbox}\begin{enumerate}
\item What data or series is being approximated?
\item Is the evaluation interpolation or extrapolation?
\item What smoothness or convergence assumptions are used?
\item How will approximation quality be checked?
\end{enumerate}\end{promptbox}
\textbf{Workspace}\workspace
\begin{reflectionbox}\begin{enumerate}
\item Which sentence or datum in these problems supplied the main mathematical constraint?
\item Which assumptions were essential to the model, and how would the model change if one were removed?
\item When would an analytical solution become inconvenient enough to justify a numerical method?
\end{enumerate}\end{reflectionbox}
\chapter{Cantor set, disappearing table, bouncing-ball distance and geometric series}\label{app:A042}
\textbf{Application ID:} \pid{A042}\quad \textbf{Source locus:} 9\quad \textbf{Source pages:} 604–608, 681, 683
\section*{Application Context}
This family uses interpolation, series, transformations, and approximation to evaluate functions or represent data when direct formulas are inconvenient. The canonical application family is \emph{Cantor set, disappearing table, bouncing-ball distance and geometric series}.
\section*{Calculus Toolkit}
\begin{tabularx}{\textwidth}{@{}lX@{}}\toprule
Primary technique & Sequences, series and convergence \\
Prerequisites & functions, series or interpolation, approximation error \\
Numerical extension & Ch. 3 (interpolation), Ch. 5 (function evaluation), Ch. 6 (special functions) \\
Analytical/numerical class & ANALYTICAL \\
\bottomrule\end{tabularx}
\section*{Modeling Questions to Ask}\begin{itemize}
\item What data or series is being approximated?
\item Is the evaluation interpolation or extrapolation?
\item What smoothness or convergence assumptions are used?
\item How will approximation quality be checked?
\end{itemize}
\begin{errorbox}Interpolation is not unrestricted extrapolation. State the interval on which the approximation is intended to be used and check approximation error when possible.\end{errorbox}
\section*{Problems}
\needspace{8\baselineskip}
\subsection*{\pid{APP-042-F-01} — Foundation}
\begin{problembox}
Two measurements in cantor set, disappearing table, bouncing-ball distance and geometric series are $(0,20)$ and $(10,35)$. Estimate the value at $x=4$ by linear interpolation.
\end{problembox}
\begin{promptbox}\begin{enumerate}
\item What data or series is being approximated?
\item Is the evaluation interpolation or extrapolation?
\item What smoothness or convergence assumptions are used?
\item How will approximation quality be checked?
\end{enumerate}\end{promptbox}
\textbf{Workspace}\workspace
\needspace{8\baselineskip}
\subsection*{\pid{APP-042-S-01} — Standard}
\begin{problembox}
A repeated contribution in cantor set, disappearing table, bouncing-ball distance and geometric series has magnitudes $310, 186.0, 111.6,\ldots$, with common ratio $r=0.60$. Determine the infinite accumulated total and state the convergence condition used.
\end{problembox}
\begin{promptbox}\begin{enumerate}
\item What data or series is being approximated?
\item Is the evaluation interpolation or extrapolation?
\item What smoothness or convergence assumptions are used?
\item How will approximation quality be checked?
\end{enumerate}\end{promptbox}
\textbf{Workspace}\workspace
\needspace{8\baselineskip}
\subsection*{\pid{APP-042-A-01} — Advanced}
\begin{problembox}
A local approximation useful in cantor set, disappearing table, bouncing-ball distance and geometric series requires $\ln(1+x)$ at $x=0.2$. Use the cubic Taylor polynomial $x-x^2/2+x^3/3$, compare with the exact value, and report the absolute error.
\end{problembox}
\begin{promptbox}\begin{enumerate}
\item What data or series is being approximated?
\item Is the evaluation interpolation or extrapolation?
\item What smoothness or convergence assumptions are used?
\item How will approximation quality be checked?
\end{enumerate}\end{promptbox}
\textbf{Workspace}\workspace
\begin{reflectionbox}\begin{enumerate}
\item Which sentence or datum in these problems supplied the main mathematical constraint?
\item Which assumptions were essential to the model, and how would the model change if one were removed?
\item When would an analytical solution become inconvenient enough to justify a numerical method?
\end{enumerate}\end{reflectionbox}
\chapter{Special-function physics: pendulum refinements, elliptic geometry and spherical harmonics}\label{app:A043}
\textbf{Application ID:} \pid{A043}\quad \textbf{Source locus:} 10; 12; advanced\quad \textbf{Source pages:} 241, 691, 910
\section*{Application Context}
This family uses interpolation, series, transformations, and approximation to evaluate functions or represent data when direct formulas are inconvenient. The canonical application family is \emph{Special-function physics: pendulum refinements, elliptic geometry and spherical harmonics}.
\section*{Calculus Toolkit}
\begin{tabularx}{\textwidth}{@{}lX@{}}\toprule
Primary technique & Special functions and nonlinear models \\
Prerequisites & functions, series or interpolation, approximation error \\
Numerical extension & Ch. 3 (interpolation), Ch. 5 (function evaluation), Ch. 6 (special functions) \\
Analytical/numerical class & NUMERICAL-NATURAL \\
\bottomrule\end{tabularx}
\section*{Modeling Questions to Ask}\begin{itemize}
\item What data or series is being approximated?
\item Is the evaluation interpolation or extrapolation?
\item What smoothness or convergence assumptions are used?
\item How will approximation quality be checked?
\end{itemize}
\begin{errorbox}Interpolation is not unrestricted extrapolation. State the interval on which the approximation is intended to be used and check approximation error when possible.\end{errorbox}
\section*{Problems}
\needspace{8\baselineskip}
\subsection*{\pid{APP-043-F-01} — Foundation}
\begin{problembox}
Two measurements in special-function physics: pendulum refinements, elliptic geometry and spherical harmonics are $(0,21)$ and $(10,36)$. Estimate the value at $x=4$ by linear interpolation.
\end{problembox}
\begin{promptbox}\begin{enumerate}
\item What data or series is being approximated?
\item Is the evaluation interpolation or extrapolation?
\item What smoothness or convergence assumptions are used?
\item How will approximation quality be checked?
\end{enumerate}\end{promptbox}
\textbf{Workspace}\workspace
\needspace{8\baselineskip}
\subsection*{\pid{APP-043-S-01} — Standard}
\begin{problembox}
A repeated contribution in special-function physics: pendulum refinements, elliptic geometry and spherical harmonics has magnitudes $315, 204.8, 133.1,\ldots$, with common ratio $r=0.65$. Determine the infinite accumulated total and state the convergence condition used.
\end{problembox}
\begin{promptbox}\begin{enumerate}
\item What data or series is being approximated?
\item Is the evaluation interpolation or extrapolation?
\item What smoothness or convergence assumptions are used?
\item How will approximation quality be checked?
\end{enumerate}\end{promptbox}
\textbf{Workspace}\workspace
\needspace{8\baselineskip}
\subsection*{\pid{APP-043-A-01} — Advanced}
\begin{problembox}
A local approximation useful in special-function physics: pendulum refinements, elliptic geometry and spherical harmonics requires $\ln(1+x)$ at $x=0.2$. Use the cubic Taylor polynomial $x-x^2/2+x^3/3$, compare with the exact value, and report the absolute error.
\end{problembox}
\begin{promptbox}\begin{enumerate}
\item What data or series is being approximated?
\item Is the evaluation interpolation or extrapolation?
\item What smoothness or convergence assumptions are used?
\item How will approximation quality be checked?
\end{enumerate}\end{promptbox}
\textbf{Workspace}\workspace
\begin{reflectionbox}\begin{enumerate}
\item Which sentence or datum in these problems supplied the main mathematical constraint?
\item Which assumptions were essential to the model, and how would the model change if one were removed?
\item When would an analytical solution become inconvenient enough to justify a numerical method?
\end{enumerate}\end{reflectionbox}
\part{Vectors and Multivariable Applications}
\chapter{Resultant force/velocity and vector mechanics}\label{app:A044}
\textbf{Application ID:} \pid{A044}\quad \textbf{Source locus:} 11.1–11.2\quad \textbf{Source pages:} 758–761
\section*{Application Context}
This family models directional and spatial quantities using vectors, gradients, linear systems, and multiple integrals. The canonical application family is \emph{Resultant force/velocity and vector mechanics}.
\section*{Calculus Toolkit}
\begin{tabularx}{\textwidth}{@{}lX@{}}\toprule
Primary technique & Vector addition and components \\
Prerequisites & vectors, partial derivatives, linear algebra \\
Numerical extension & Ch. 2 (linear algebra), Ch. 10 (multivariable minimization), Ch. 11 (eigensystems) \\
Analytical/numerical class & NUMERICAL-OPTIONAL \\
\bottomrule\end{tabularx}
\section*{Modeling Questions to Ask}\begin{itemize}
\item Which quantities are scalars and which are vectors?
\item What coordinate system is being used?
\item Is a dot product, cross product, gradient, or multiple integral appropriate?
\item What geometric meaning should the result have?
\end{itemize}
\begin{errorbox}Keep scalar, vector, and coordinate quantities distinct. For gradients, state the point and normalize a direction vector before taking a directional derivative.\end{errorbox}
\section*{Problems}
\needspace{8\baselineskip}
\subsection*{\pid{APP-044-F-01} — Foundation}
\begin{problembox}
Two vector quantities in resultant force/velocity and vector mechanics are $\mathbf a=(4,-1,3)$ and $\mathbf b=(1,2,-2)$. Compute $\mathbf a\cdot\mathbf b$ and $\mathbf a\times\mathbf b$, and state one geometric interpretation of each.
\end{problembox}
\begin{promptbox}\begin{enumerate}
\item Which quantities are scalars and which are vectors?
\item What coordinate system is being used?
\item Is a dot product, cross product, gradient, or multiple integral appropriate?
\item What geometric meaning should the result have?
\end{enumerate}\end{promptbox}
\textbf{Workspace}\workspace
\needspace{8\baselineskip}
\subsection*{\pid{APP-044-S-01} — Standard}
\begin{problembox}
A scalar field used in resultant force/velocity and vector mechanics is $F(x,y,z)=x^2+2y^2+3z$. At $P=(3,2,1)$, find the gradient and the directional derivative toward the direction $(1,2,1)$.
\end{problembox}
\begin{promptbox}\begin{enumerate}
\item Which quantities are scalars and which are vectors?
\item What coordinate system is being used?
\item Is a dot product, cross product, gradient, or multiple integral appropriate?
\item What geometric meaning should the result have?
\end{enumerate}\end{promptbox}
\textbf{Workspace}\workspace
\needspace{8\baselineskip}
\subsection*{\pid{APP-044-A-01} — Advanced}
\begin{problembox}
A planar density model related to resultant force/velocity and vector mechanics is $\rho(x,y)=x+y$ on the rectangle $0\le x\le 4$, $0\le y\le 3$. Compute the total mass.
\end{problembox}
\begin{promptbox}\begin{enumerate}
\item Which quantities are scalars and which are vectors?
\item What coordinate system is being used?
\item Is a dot product, cross product, gradient, or multiple integral appropriate?
\item What geometric meaning should the result have?
\end{enumerate}\end{promptbox}
\textbf{Workspace}\workspace
\begin{reflectionbox}\begin{enumerate}
\item Which sentence or datum in these problems supplied the main mathematical constraint?
\item Which assumptions were essential to the model, and how would the model change if one were removed?
\item When would an analytical solution become inconvenient enough to justify a numerical method?
\end{enumerate}\end{reflectionbox}
\chapter{Torque, work from vector forces and rigid-mechanics geometry}\label{app:A045}
\textbf{Application ID:} \pid{A045}\quad \textbf{Source locus:} 11.3–11.5\quad \textbf{Source pages:} 770–816
\section*{Application Context}
This family models directional and spatial quantities using vectors, gradients, linear systems, and multiple integrals. The canonical application family is \emph{Torque, work from vector forces and rigid-mechanics geometry}.
\section*{Calculus Toolkit}
\begin{tabularx}{\textwidth}{@{}lX@{}}\toprule
Primary technique & Dot/cross products and vector geometry \\
Prerequisites & vectors, partial derivatives, linear algebra \\
Numerical extension & Ch. 2 (linear algebra), Ch. 10 (multivariable minimization), Ch. 11 (eigensystems) \\
Analytical/numerical class & NUMERICAL-NATURAL \\
\bottomrule\end{tabularx}
\section*{Modeling Questions to Ask}\begin{itemize}
\item Which quantities are scalars and which are vectors?
\item What coordinate system is being used?
\item Is a dot product, cross product, gradient, or multiple integral appropriate?
\item What geometric meaning should the result have?
\end{itemize}
\begin{errorbox}Keep scalar, vector, and coordinate quantities distinct. For gradients, state the point and normalize a direction vector before taking a directional derivative.\end{errorbox}
\section*{Problems}
\needspace{8\baselineskip}
\subsection*{\pid{APP-045-F-01} — Foundation}
\begin{problembox}
Two vector quantities in torque, work from vector forces and rigid-mechanics geometry are $\mathbf a=(2,-2,3)$ and $\mathbf b=(1,3,-2)$. Compute $\mathbf a\cdot\mathbf b$ and $\mathbf a\times\mathbf b$, and state one geometric interpretation of each.
\end{problembox}
\begin{promptbox}\begin{enumerate}
\item Which quantities are scalars and which are vectors?
\item What coordinate system is being used?
\item Is a dot product, cross product, gradient, or multiple integral appropriate?
\item What geometric meaning should the result have?
\end{enumerate}\end{promptbox}
\textbf{Workspace}\workspace
\needspace{8\baselineskip}
\subsection*{\pid{APP-045-S-01} — Standard}
\begin{problembox}
A scalar field used in torque, work from vector forces and rigid-mechanics geometry is $F(x,y,z)=x^2+2y^2+3z$. At $P=(1,2,1)$, find the gradient and the directional derivative toward the direction $(1,2,1)$.
\end{problembox}
\begin{promptbox}\begin{enumerate}
\item Which quantities are scalars and which are vectors?
\item What coordinate system is being used?
\item Is a dot product, cross product, gradient, or multiple integral appropriate?
\item What geometric meaning should the result have?
\end{enumerate}\end{promptbox}
\textbf{Workspace}\workspace
\needspace{8\baselineskip}
\subsection*{\pid{APP-045-A-01} — Advanced}
\begin{problembox}
A planar density model related to torque, work from vector forces and rigid-mechanics geometry is $\rho(x,y)=x+y$ on the rectangle $0\le x\le 2$, $0\le y\le 4$. Compute the total mass.
\end{problembox}
\begin{promptbox}\begin{enumerate}
\item Which quantities are scalars and which are vectors?
\item What coordinate system is being used?
\item Is a dot product, cross product, gradient, or multiple integral appropriate?
\item What geometric meaning should the result have?
\end{enumerate}\end{promptbox}
\textbf{Workspace}\workspace
\begin{reflectionbox}\begin{enumerate}
\item Which sentence or datum in these problems supplied the main mathematical constraint?
\item Which assumptions were essential to the model, and how would the model change if one were removed?
\item When would an analytical solution become inconvenient enough to justify a numerical method?
\end{enumerate}\end{reflectionbox}
\chapter{Cable tension, load supports, navigation and surveying}\label{app:A046}
\textbf{Application ID:} \pid{A046}\quad \textbf{Source locus:} 10–11\quad \textbf{Source pages:} 565, 699, 761, 769, 818
\section*{Application Context}
This family models directional and spatial quantities using vectors, gradients, linear systems, and multiple integrals. The canonical application family is \emph{Cable tension, load supports, navigation and surveying}.
\section*{Calculus Toolkit}
\begin{tabularx}{\textwidth}{@{}lX@{}}\toprule
Primary technique & Linear/vector systems \\
Prerequisites & vectors, partial derivatives, linear algebra \\
Numerical extension & Ch. 2 (linear algebra), Ch. 10 (multivariable minimization), Ch. 11 (eigensystems) \\
Analytical/numerical class & NUMERICAL-NATURAL \\
\bottomrule\end{tabularx}
\section*{Modeling Questions to Ask}\begin{itemize}
\item Which quantities are scalars and which are vectors?
\item What coordinate system is being used?
\item Is a dot product, cross product, gradient, or multiple integral appropriate?
\item What geometric meaning should the result have?
\end{itemize}
\begin{errorbox}Keep scalar, vector, and coordinate quantities distinct. For gradients, state the point and normalize a direction vector before taking a directional derivative.\end{errorbox}
\section*{Problems}
\needspace{8\baselineskip}
\subsection*{\pid{APP-046-F-01} — Foundation}
\begin{problembox}
Two vector quantities in cable tension, load supports, navigation and surveying are $\mathbf a=(3,-1,3)$ and $\mathbf b=(1,4,-2)$. Compute $\mathbf a\cdot\mathbf b$ and $\mathbf a\times\mathbf b$, and state one geometric interpretation of each.
\end{problembox}
\begin{promptbox}\begin{enumerate}
\item Which quantities are scalars and which are vectors?
\item What coordinate system is being used?
\item Is a dot product, cross product, gradient, or multiple integral appropriate?
\item What geometric meaning should the result have?
\end{enumerate}\end{promptbox}
\textbf{Workspace}\workspace
\needspace{8\baselineskip}
\subsection*{\pid{APP-046-S-01} — Standard}
\begin{problembox}
A scalar field used in cable tension, load supports, navigation and surveying is $F(x,y,z)=x^2+2y^2+3z$. At $P=(2,2,1)$, find the gradient and the directional derivative toward the direction $(1,2,1)$.
\end{problembox}
\begin{promptbox}\begin{enumerate}
\item Which quantities are scalars and which are vectors?
\item What coordinate system is being used?
\item Is a dot product, cross product, gradient, or multiple integral appropriate?
\item What geometric meaning should the result have?
\end{enumerate}\end{promptbox}
\textbf{Workspace}\workspace
\needspace{8\baselineskip}
\subsection*{\pid{APP-046-A-01} — Advanced}
\begin{problembox}
A planar density model related to cable tension, load supports, navigation and surveying is $\rho(x,y)=x+y$ on the rectangle $0\le x\le 3$, $0\le y\le 3$. Compute the total mass.
\end{problembox}
\begin{promptbox}\begin{enumerate}
\item Which quantities are scalars and which are vectors?
\item What coordinate system is being used?
\item Is a dot product, cross product, gradient, or multiple integral appropriate?
\item What geometric meaning should the result have?
\end{enumerate}\end{promptbox}
\textbf{Workspace}\workspace
\begin{reflectionbox}\begin{enumerate}
\item Which sentence or datum in these problems supplied the main mathematical constraint?
\item Which assumptions were essential to the model, and how would the model change if one were removed?
\item When would an analytical solution become inconvenient enough to justify a numerical method?
\end{enumerate}\end{reflectionbox}
\chapter{Ideal-gas, meteorology and apparent-temperature multivariable models}\label{app:A047}
\textbf{Application ID:} \pid{A047}\quad \textbf{Source locus:} 13.2–13.6\quad \textbf{Source pages:} 883, 903, 918
\section*{Application Context}
This family models directional and spatial quantities using vectors, gradients, linear systems, and multiple integrals. The canonical application family is \emph{Ideal-gas, meteorology and apparent-temperature multivariable models}.
\section*{Calculus Toolkit}
\begin{tabularx}{\textwidth}{@{}lX@{}}\toprule
Primary technique & Partial derivatives and gradients \\
Prerequisites & vectors, partial derivatives, linear algebra \\
Numerical extension & Ch. 2 (linear algebra), Ch. 10 (multivariable minimization), Ch. 11 (eigensystems) \\
Analytical/numerical class & NUMERICAL-NATURAL \\
\bottomrule\end{tabularx}
\section*{Modeling Questions to Ask}\begin{itemize}
\item Which quantities are scalars and which are vectors?
\item What coordinate system is being used?
\item Is a dot product, cross product, gradient, or multiple integral appropriate?
\item What geometric meaning should the result have?
\end{itemize}
\begin{errorbox}Keep scalar, vector, and coordinate quantities distinct. For gradients, state the point and normalize a direction vector before taking a directional derivative.\end{errorbox}
\section*{Problems}
\needspace{8\baselineskip}
\subsection*{\pid{APP-047-F-01} — Foundation}
\begin{problembox}
Two vector quantities in ideal-gas, meteorology and apparent-temperature multivariable models are $\mathbf a=(4,-2,3)$ and $\mathbf b=(1,5,-2)$. Compute $\mathbf a\cdot\mathbf b$ and $\mathbf a\times\mathbf b$, and state one geometric interpretation of each.
\end{problembox}
\begin{promptbox}\begin{enumerate}
\item Which quantities are scalars and which are vectors?
\item What coordinate system is being used?
\item Is a dot product, cross product, gradient, or multiple integral appropriate?
\item What geometric meaning should the result have?
\end{enumerate}\end{promptbox}
\textbf{Workspace}\workspace
\needspace{8\baselineskip}
\subsection*{\pid{APP-047-S-01} — Standard}
\begin{problembox}
A scalar field used in ideal-gas, meteorology and apparent-temperature multivariable models is $F(x,y,z)=x^2+2y^2+3z$. At $P=(3,2,1)$, find the gradient and the directional derivative toward the direction $(1,2,1)$.
\end{problembox}
\begin{promptbox}\begin{enumerate}
\item Which quantities are scalars and which are vectors?
\item What coordinate system is being used?
\item Is a dot product, cross product, gradient, or multiple integral appropriate?
\item What geometric meaning should the result have?
\end{enumerate}\end{promptbox}
\textbf{Workspace}\workspace
\needspace{8\baselineskip}
\subsection*{\pid{APP-047-A-01} — Advanced}
\begin{problembox}
A planar density model related to ideal-gas, meteorology and apparent-temperature multivariable models is $\rho(x,y)=x+y$ on the rectangle $0\le x\le 4$, $0\le y\le 4$. Compute the total mass.
\end{problembox}
\begin{promptbox}\begin{enumerate}
\item Which quantities are scalars and which are vectors?
\item What coordinate system is being used?
\item Is a dot product, cross product, gradient, or multiple integral appropriate?
\item What geometric meaning should the result have?
\end{enumerate}\end{promptbox}
\textbf{Workspace}\workspace
\begin{reflectionbox}\begin{enumerate}
\item Which sentence or datum in these problems supplied the main mathematical constraint?
\item Which assumptions were essential to the model, and how would the model change if one were removed?
\item When would an analytical solution become inconvenient enough to justify a numerical method?
\end{enumerate}\end{reflectionbox}
\chapter{Cobb–Douglas production, marginal analysis and economic surfaces}\label{app:A048}
\textbf{Application ID:} \pid{A048}\quad \textbf{Source locus:} 13\quad \textbf{Source pages:} 877, 882, 883, 959, 964, 967
\section*{Application Context}
This family models directional and spatial quantities using vectors, gradients, linear systems, and multiple integrals. The canonical application family is \emph{Cobb–Douglas production, marginal analysis and economic surfaces}.
\section*{Calculus Toolkit}
\begin{tabularx}{\textwidth}{@{}lX@{}}\toprule
Primary technique & Partial derivatives, gradients and constrained extrema \\
Prerequisites & vectors, partial derivatives, linear algebra \\
Numerical extension & Ch. 2 (linear algebra), Ch. 10 (multivariable minimization), Ch. 11 (eigensystems) \\
Analytical/numerical class & NUMERICAL-NATURAL \\
\bottomrule\end{tabularx}
\section*{Modeling Questions to Ask}\begin{itemize}
\item Which quantities are scalars and which are vectors?
\item What coordinate system is being used?
\item Is a dot product, cross product, gradient, or multiple integral appropriate?
\item What geometric meaning should the result have?
\end{itemize}
\begin{errorbox}Keep scalar, vector, and coordinate quantities distinct. For gradients, state the point and normalize a direction vector before taking a directional derivative.\end{errorbox}
\section*{Problems}
\needspace{8\baselineskip}
\subsection*{\pid{APP-048-F-01} — Foundation}
\begin{problembox}
Two vector quantities in cobb–douglas production, marginal analysis and economic surfaces are $\mathbf a=(2,-1,3)$ and $\mathbf b=(1,2,-2)$. Compute $\mathbf a\cdot\mathbf b$ and $\mathbf a\times\mathbf b$, and state one geometric interpretation of each.
\end{problembox}
\begin{promptbox}\begin{enumerate}
\item Which quantities are scalars and which are vectors?
\item What coordinate system is being used?
\item Is a dot product, cross product, gradient, or multiple integral appropriate?
\item What geometric meaning should the result have?
\end{enumerate}\end{promptbox}
\textbf{Workspace}\workspace
\needspace{8\baselineskip}
\subsection*{\pid{APP-048-S-01} — Standard}
\begin{problembox}
A scalar field used in cobb–douglas production, marginal analysis and economic surfaces is $F(x,y,z)=x^2+2y^2+3z$. At $P=(1,2,1)$, find the gradient and the directional derivative toward the direction $(1,2,1)$.
\end{problembox}
\begin{promptbox}\begin{enumerate}
\item Which quantities are scalars and which are vectors?
\item What coordinate system is being used?
\item Is a dot product, cross product, gradient, or multiple integral appropriate?
\item What geometric meaning should the result have?
\end{enumerate}\end{promptbox}
\textbf{Workspace}\workspace
\needspace{8\baselineskip}
\subsection*{\pid{APP-048-A-01} — Advanced}
\begin{problembox}
A planar density model related to cobb–douglas production, marginal analysis and economic surfaces is $\rho(x,y)=x+y$ on the rectangle $0\le x\le 2$, $0\le y\le 3$. Compute the total mass.
\end{problembox}
\begin{promptbox}\begin{enumerate}
\item Which quantities are scalars and which are vectors?
\item What coordinate system is being used?
\item Is a dot product, cross product, gradient, or multiple integral appropriate?
\item What geometric meaning should the result have?
\end{enumerate}\end{promptbox}
\textbf{Workspace}\workspace
\begin{reflectionbox}\begin{enumerate}
\item Which sentence or datum in these problems supplied the main mathematical constraint?
\item Which assumptions were essential to the model, and how would the model change if one were removed?
\item When would an analytical solution become inconvenient enough to justify a numerical method?
\end{enumerate}\end{reflectionbox}
\chapter{Hardy–Weinberg law, Shannon diversity and constrained population composition}\label{app:A049}
\textbf{Application ID:} \pid{A049}\quad \textbf{Source locus:} 13.9–13.10\quad \textbf{Source pages:} 953
\section*{Application Context}
This family models directional and spatial quantities using vectors, gradients, linear systems, and multiple integrals. The canonical application family is \emph{Hardy–Weinberg law, Shannon diversity and constrained population composition}.
\section*{Calculus Toolkit}
\begin{tabularx}{\textwidth}{@{}lX@{}}\toprule
Primary technique & Constrained multivariable optimization \\
Prerequisites & vectors, partial derivatives, linear algebra \\
Numerical extension & Ch. 2 (linear algebra), Ch. 10 (multivariable minimization), Ch. 11 (eigensystems) \\
Analytical/numerical class & NUMERICAL-NATURAL \\
\bottomrule\end{tabularx}
\section*{Modeling Questions to Ask}\begin{itemize}
\item Which quantities are scalars and which are vectors?
\item What coordinate system is being used?
\item Is a dot product, cross product, gradient, or multiple integral appropriate?
\item What geometric meaning should the result have?
\end{itemize}
\begin{errorbox}Keep scalar, vector, and coordinate quantities distinct. For gradients, state the point and normalize a direction vector before taking a directional derivative.\end{errorbox}
\section*{Problems}
\needspace{8\baselineskip}
\subsection*{\pid{APP-049-F-01} — Foundation}
\begin{problembox}
Two vector quantities in hardy–weinberg law, shannon diversity and constrained population composition are $\mathbf a=(3,-2,3)$ and $\mathbf b=(1,3,-2)$. Compute $\mathbf a\cdot\mathbf b$ and $\mathbf a\times\mathbf b$, and state one geometric interpretation of each.
\end{problembox}
\begin{promptbox}\begin{enumerate}
\item Which quantities are scalars and which are vectors?
\item What coordinate system is being used?
\item Is a dot product, cross product, gradient, or multiple integral appropriate?
\item What geometric meaning should the result have?
\end{enumerate}\end{promptbox}
\textbf{Workspace}\workspace
\needspace{8\baselineskip}
\subsection*{\pid{APP-049-S-01} — Standard}
\begin{problembox}
A scalar field used in hardy–weinberg law, shannon diversity and constrained population composition is $F(x,y,z)=x^2+2y^2+3z$. At $P=(2,2,1)$, find the gradient and the directional derivative toward the direction $(1,2,1)$.
\end{problembox}
\begin{promptbox}\begin{enumerate}
\item Which quantities are scalars and which are vectors?
\item What coordinate system is being used?
\item Is a dot product, cross product, gradient, or multiple integral appropriate?
\item What geometric meaning should the result have?
\end{enumerate}\end{promptbox}
\textbf{Workspace}\workspace
\needspace{8\baselineskip}
\subsection*{\pid{APP-049-A-01} — Advanced}
\begin{problembox}
A planar density model related to hardy–weinberg law, shannon diversity and constrained population composition is $\rho(x,y)=x+y$ on the rectangle $0\le x\le 3$, $0\le y\le 4$. Compute the total mass.
\end{problembox}
\begin{promptbox}\begin{enumerate}
\item Which quantities are scalars and which are vectors?
\item What coordinate system is being used?
\item Is a dot product, cross product, gradient, or multiple integral appropriate?
\item What geometric meaning should the result have?
\end{enumerate}\end{promptbox}
\textbf{Workspace}\workspace
\begin{reflectionbox}\begin{enumerate}
\item Which sentence or datum in these problems supplied the main mathematical constraint?
\item Which assumptions were essential to the model, and how would the model change if one were removed?
\item When would an analytical solution become inconvenient enough to justify a numerical method?
\end{enumerate}\end{reflectionbox}
\chapter{Mass, moments and density distributions in two/three dimensions}\label{app:A050}
\textbf{Application ID:} \pid{A050}\quad \textbf{Source locus:} 14; 15\quad \textbf{Source pages:} 1059, 1065, 1066
\section*{Application Context}
This family models directional and spatial quantities using vectors, gradients, linear systems, and multiple integrals. The canonical application family is \emph{Mass, moments and density distributions in two/three dimensions}.
\section*{Calculus Toolkit}
\begin{tabularx}{\textwidth}{@{}lX@{}}\toprule
Primary technique & Multiple integrals and vector fields \\
Prerequisites & vectors, partial derivatives, linear algebra \\
Numerical extension & Ch. 2 (linear algebra), Ch. 10 (multivariable minimization), Ch. 11 (eigensystems) \\
Analytical/numerical class & NUMERICAL-NATURAL \\
\bottomrule\end{tabularx}
\section*{Modeling Questions to Ask}\begin{itemize}
\item Which quantities are scalars and which are vectors?
\item What coordinate system is being used?
\item Is a dot product, cross product, gradient, or multiple integral appropriate?
\item What geometric meaning should the result have?
\end{itemize}
\begin{errorbox}Keep scalar, vector, and coordinate quantities distinct. For gradients, state the point and normalize a direction vector before taking a directional derivative.\end{errorbox}
\section*{Problems}
\needspace{8\baselineskip}
\subsection*{\pid{APP-050-F-01} — Foundation}
\begin{problembox}
Two vector quantities in mass, moments and density distributions in two/three dimensions are $\mathbf a=(4,-1,3)$ and $\mathbf b=(1,4,-2)$. Compute $\mathbf a\cdot\mathbf b$ and $\mathbf a\times\mathbf b$, and state one geometric interpretation of each.
\end{problembox}
\begin{promptbox}\begin{enumerate}
\item Which quantities are scalars and which are vectors?
\item What coordinate system is being used?
\item Is a dot product, cross product, gradient, or multiple integral appropriate?
\item What geometric meaning should the result have?
\end{enumerate}\end{promptbox}
\textbf{Workspace}\workspace
\needspace{8\baselineskip}
\subsection*{\pid{APP-050-S-01} — Standard}
\begin{problembox}
A scalar field used in mass, moments and density distributions in two/three dimensions is $F(x,y,z)=x^2+2y^2+3z$. At $P=(3,2,1)$, find the gradient and the directional derivative toward the direction $(1,2,1)$.
\end{problembox}
\begin{promptbox}\begin{enumerate}
\item Which quantities are scalars and which are vectors?
\item What coordinate system is being used?
\item Is a dot product, cross product, gradient, or multiple integral appropriate?
\item What geometric meaning should the result have?
\end{enumerate}\end{promptbox}
\textbf{Workspace}\workspace
\needspace{8\baselineskip}
\subsection*{\pid{APP-050-A-01} — Advanced}
\begin{problembox}
A planar density model related to mass, moments and density distributions in two/three dimensions is $\rho(x,y)=x+y$ on the rectangle $0\le x\le 4$, $0\le y\le 3$. Compute the total mass.
\end{problembox}
\begin{promptbox}\begin{enumerate}
\item Which quantities are scalars and which are vectors?
\item What coordinate system is being used?
\item Is a dot product, cross product, gradient, or multiple integral appropriate?
\item What geometric meaning should the result have?
\end{enumerate}\end{promptbox}
\textbf{Workspace}\workspace
\begin{reflectionbox}\begin{enumerate}
\item Which sentence or datum in these problems supplied the main mathematical constraint?
\item Which assumptions were essential to the model, and how would the model change if one were removed?
\item When would an analytical solution become inconvenient enough to justify a numerical method?
\end{enumerate}\end{reflectionbox}
\chapter{Field strength, electric potential and gravitational potential}\label{app:A051}
\textbf{Application ID:} \pid{A051}\quad \textbf{Source locus:} 7; 13; 15\quad \textbf{Source pages:} 548, 882, 1045
\section*{Application Context}
This family models directional and spatial quantities using vectors, gradients, linear systems, and multiple integrals. The canonical application family is \emph{Field strength, electric potential and gravitational potential}.
\section*{Calculus Toolkit}
\begin{tabularx}{\textwidth}{@{}lX@{}}\toprule
Primary technique & Scalar/vector fields and gradients \\
Prerequisites & vectors, partial derivatives, linear algebra \\
Numerical extension & Ch. 2 (linear algebra), Ch. 10 (multivariable minimization), Ch. 11 (eigensystems) \\
Analytical/numerical class & NUMERICAL-NATURAL \\
\bottomrule\end{tabularx}
\section*{Modeling Questions to Ask}\begin{itemize}
\item Which quantities are scalars and which are vectors?
\item What coordinate system is being used?
\item Is a dot product, cross product, gradient, or multiple integral appropriate?
\item What geometric meaning should the result have?
\end{itemize}
\begin{errorbox}Keep scalar, vector, and coordinate quantities distinct. For gradients, state the point and normalize a direction vector before taking a directional derivative.\end{errorbox}
\section*{Problems}
\needspace{8\baselineskip}
\subsection*{\pid{APP-051-F-01} — Foundation}
\begin{problembox}
Two vector quantities in field strength, electric potential and gravitational potential are $\mathbf a=(2,-2,3)$ and $\mathbf b=(1,5,-2)$. Compute $\mathbf a\cdot\mathbf b$ and $\mathbf a\times\mathbf b$, and state one geometric interpretation of each.
\end{problembox}
\begin{promptbox}\begin{enumerate}
\item Which quantities are scalars and which are vectors?
\item What coordinate system is being used?
\item Is a dot product, cross product, gradient, or multiple integral appropriate?
\item What geometric meaning should the result have?
\end{enumerate}\end{promptbox}
\textbf{Workspace}\workspace
\needspace{8\baselineskip}
\subsection*{\pid{APP-051-S-01} — Standard}
\begin{problembox}
A scalar field used in field strength, electric potential and gravitational potential is $F(x,y,z)=x^2+2y^2+3z$. At $P=(1,2,1)$, find the gradient and the directional derivative toward the direction $(1,2,1)$.
\end{problembox}
\begin{promptbox}\begin{enumerate}
\item Which quantities are scalars and which are vectors?
\item What coordinate system is being used?
\item Is a dot product, cross product, gradient, or multiple integral appropriate?
\item What geometric meaning should the result have?
\end{enumerate}\end{promptbox}
\textbf{Workspace}\workspace
\needspace{8\baselineskip}
\subsection*{\pid{APP-051-A-01} — Advanced}
\begin{problembox}
A planar density model related to field strength, electric potential and gravitational potential is $\rho(x,y)=x+y$ on the rectangle $0\le x\le 2$, $0\le y\le 4$. Compute the total mass.
\end{problembox}
\begin{promptbox}\begin{enumerate}
\item Which quantities are scalars and which are vectors?
\item What coordinate system is being used?
\item Is a dot product, cross product, gradient, or multiple integral appropriate?
\item What geometric meaning should the result have?
\end{enumerate}\end{promptbox}
\textbf{Workspace}\workspace
\begin{reflectionbox}\begin{enumerate}
\item Which sentence or datum in these problems supplied the main mathematical constraint?
\item Which assumptions were essential to the model, and how would the model change if one were removed?
\item When would an analytical solution become inconvenient enough to justify a numerical method?
\end{enumerate}\end{reflectionbox}
\part{Differential Equations}
\chapter{Radioactive decay and carbon dating}\label{app:A052}
\textbf{Application ID:} \pid{A052}\quad \textbf{Source locus:} 5.5; 6.2–6.3\quad \textbf{Source pages:} 352, 417, 421, 429, 439
\section*{Application Context}
This family models a state whose derivative is governed by the state itself, external inputs, or coupled physical laws. Initial conditions select a unique trajectory when the model permits one. The canonical application family is \emph{Radioactive decay and carbon dating}.
\section*{Calculus Toolkit}
\begin{tabularx}{\textwidth}{@{}lX@{}}\toprule
Primary technique & Exponential decay ODE \\
Prerequisites & derivatives, separation or linear ODE methods, initial conditions \\
Numerical extension & Ch. 17 (initial-value ODEs); Ch. 18 for two-point boundary values \\
Analytical/numerical class & NUMERICAL-OPTIONAL \\
\bottomrule\end{tabularx}
\section*{Modeling Questions to Ask}\begin{itemize}
\item What is the state variable?
\item What law specifies its derivative?
\item What initial/boundary data are available?
\item Can the solution be checked by substitution?
\end{itemize}
\begin{errorbox}An ODE solution is not determined until initial or boundary data are applied. Verify the solution by substitution and check the initial condition.\end{errorbox}
\section*{Problems}
\needspace{8\baselineskip}
\subsection*{\pid{APP-052-F-01} — Foundation}
\begin{problembox}
A state variable in radioactive decay and carbon dating satisfies $y'=0.08y$ with $y(0)=70$. Solve the initial-value problem and find $y(6)$.
\end{problembox}
\begin{promptbox}\begin{enumerate}
\item What is the state variable?
\item What law specifies its derivative?
\item What initial/boundary data are available?
\item Can the solution be checked by substitution?
\end{enumerate}\end{promptbox}
\textbf{Workspace}\workspace
\needspace{8\baselineskip}
\subsection*{\pid{APP-052-S-01} — Standard}
\begin{problembox}
A bounded-growth model for radioactive decay and carbon dating satisfies $y'=0.32y(1-y/540)$ with $y(0)=40$. Use the logistic solution to estimate $y(8)$.
\end{problembox}
\begin{promptbox}\begin{enumerate}
\item What is the state variable?
\item What law specifies its derivative?
\item What initial/boundary data are available?
\item Can the solution be checked by substitution?
\end{enumerate}\end{promptbox}
\textbf{Workspace}\workspace
\needspace{8\baselineskip}
\subsection*{\pid{APP-052-A-01} — Advanced}
\begin{problembox}
A numerical extension of radioactive decay and carbon dating uses $y'=-0.5y+10$, $y(0)=0$. Apply Euler's method with step $h=0.5$ to approximate $y(2)$, then compare with the exact solution $y=20(1-e^{-0.5t})$.
\end{problembox}
\begin{promptbox}\begin{enumerate}
\item What is the state variable?
\item What law specifies its derivative?
\item What initial/boundary data are available?
\item Can the solution be checked by substitution?
\end{enumerate}\end{promptbox}
\textbf{Workspace}\workspace
\begin{reflectionbox}\begin{enumerate}
\item Which sentence or datum in these problems supplied the main mathematical constraint?
\item Which assumptions were essential to the model, and how would the model change if one were removed?
\item When would an analytical solution become inconvenient enough to justify a numerical method?
\end{enumerate}\end{reflectionbox}
\chapter{Bacterial culture, tree/dog growth and population growth}\label{app:A053}
\textbf{Application ID:} \pid{A053}\quad \textbf{Source locus:} 5–6\quad \textbf{Source pages:} 130, 256, 294, 332, 357, 418, 421, 425, 428, 430, 440
\section*{Application Context}
This family models a state whose derivative is governed by the state itself, external inputs, or coupled physical laws. Initial conditions select a unique trajectory when the model permits one. The canonical application family is \emph{Bacterial culture, tree/dog growth and population growth}.
\section*{Calculus Toolkit}
\begin{tabularx}{\textwidth}{@{}lX@{}}\toprule
Primary technique & Growth/logistic ODEs \\
Prerequisites & derivatives, separation or linear ODE methods, initial conditions \\
Numerical extension & Ch. 17 (initial-value ODEs); Ch. 18 for two-point boundary values \\
Analytical/numerical class & NUMERICAL-NATURAL \\
\bottomrule\end{tabularx}
\section*{Modeling Questions to Ask}\begin{itemize}
\item What is the state variable?
\item What law specifies its derivative?
\item What initial/boundary data are available?
\item Can the solution be checked by substitution?
\end{itemize}
\begin{errorbox}An ODE solution is not determined until initial or boundary data are applied. Verify the solution by substitution and check the initial condition.\end{errorbox}
\section*{Problems}
\needspace{8\baselineskip}
\subsection*{\pid{APP-053-F-01} — Foundation}
\begin{problembox}
A state variable in bacterial culture, tree/dog growth and population growth satisfies $y'=0.09y$ with $y(0)=75$. Solve the initial-value problem and find $y(6)$.
\end{problembox}
\begin{promptbox}\begin{enumerate}
\item What is the state variable?
\item What law specifies its derivative?
\item What initial/boundary data are available?
\item Can the solution be checked by substitution?
\end{enumerate}\end{promptbox}
\textbf{Workspace}\workspace
\needspace{8\baselineskip}
\subsection*{\pid{APP-053-S-01} — Standard}
\begin{problembox}
A bounded-growth model for bacterial culture, tree/dog growth and population growth satisfies $y'=0.34y(1-y/560)$ with $y(0)=45$. Use the logistic solution to estimate $y(8)$.
\end{problembox}
\begin{promptbox}\begin{enumerate}
\item What is the state variable?
\item What law specifies its derivative?
\item What initial/boundary data are available?
\item Can the solution be checked by substitution?
\end{enumerate}\end{promptbox}
\textbf{Workspace}\workspace
\needspace{8\baselineskip}
\subsection*{\pid{APP-053-A-01} — Advanced}
\begin{problembox}
A numerical extension of bacterial culture, tree/dog growth and population growth uses $y'=-0.5y+10$, $y(0)=0$. Apply Euler's method with step $h=0.5$ to approximate $y(2)$, then compare with the exact solution $y=20(1-e^{-0.5t})$.
\end{problembox}
\begin{promptbox}\begin{enumerate}
\item What is the state variable?
\item What law specifies its derivative?
\item What initial/boundary data are available?
\item Can the solution be checked by substitution?
\end{enumerate}\end{promptbox}
\textbf{Workspace}\workspace
\begin{reflectionbox}\begin{enumerate}
\item Which sentence or datum in these problems supplied the main mathematical constraint?
\item Which assumptions were essential to the model, and how would the model change if one were removed?
\item When would an analytical solution become inconvenient enough to justify a numerical method?
\end{enumerate}\end{reflectionbox}
\chapter{Endangered species, epidemic and ecological population models}\label{app:A054}
\textbf{Application ID:} \pid{A054}\quad \textbf{Source locus:} 6; 8\quad \textbf{Source pages:} 430, 442, 558
\section*{Application Context}
This family models a state whose derivative is governed by the state itself, external inputs, or coupled physical laws. Initial conditions select a unique trajectory when the model permits one. The canonical application family is \emph{Endangered species, epidemic and ecological population models}.
\section*{Calculus Toolkit}
\begin{tabularx}{\textwidth}{@{}lX@{}}\toprule
Primary technique & Coupled/separable population ODEs \\
Prerequisites & derivatives, separation or linear ODE methods, initial conditions \\
Numerical extension & Ch. 17 (initial-value ODEs); Ch. 18 for two-point boundary values \\
Analytical/numerical class & NUMERICAL-NATURAL \\
\bottomrule\end{tabularx}
\section*{Modeling Questions to Ask}\begin{itemize}
\item What is the state variable?
\item What law specifies its derivative?
\item What initial/boundary data are available?
\item Can the solution be checked by substitution?
\end{itemize}
\begin{errorbox}An ODE solution is not determined until initial or boundary data are applied. Verify the solution by substitution and check the initial condition.\end{errorbox}
\section*{Problems}
\needspace{8\baselineskip}
\subsection*{\pid{APP-054-F-01} — Foundation}
\begin{problembox}
A state variable in endangered species, epidemic and ecological population models satisfies $y'=0.10y$ with $y(0)=50$. Solve the initial-value problem and find $y(6)$.
\end{problembox}
\begin{promptbox}\begin{enumerate}
\item What is the state variable?
\item What law specifies its derivative?
\item What initial/boundary data are available?
\item Can the solution be checked by substitution?
\end{enumerate}\end{promptbox}
\textbf{Workspace}\workspace
\needspace{8\baselineskip}
\subsection*{\pid{APP-054-S-01} — Standard}
\begin{problembox}
A bounded-growth model for endangered species, epidemic and ecological population models satisfies $y'=0.30y(1-y/580)$ with $y(0)=50$. Use the logistic solution to estimate $y(8)$.
\end{problembox}
\begin{promptbox}\begin{enumerate}
\item What is the state variable?
\item What law specifies its derivative?
\item What initial/boundary data are available?
\item Can the solution be checked by substitution?
\end{enumerate}\end{promptbox}
\textbf{Workspace}\workspace
\needspace{8\baselineskip}
\subsection*{\pid{APP-054-A-01} — Advanced}
\begin{problembox}
A numerical extension of endangered species, epidemic and ecological population models uses $y'=-0.5y+10$, $y(0)=0$. Apply Euler's method with step $h=0.5$ to approximate $y(2)$, then compare with the exact solution $y=20(1-e^{-0.5t})$.
\end{problembox}
\begin{promptbox}\begin{enumerate}
\item What is the state variable?
\item What law specifies its derivative?
\item What initial/boundary data are available?
\item Can the solution be checked by substitution?
\end{enumerate}\end{promptbox}
\textbf{Workspace}\workspace
\begin{reflectionbox}\begin{enumerate}
\item Which sentence or datum in these problems supplied the main mathematical constraint?
\item Which assumptions were essential to the model, and how would the model change if one were removed?
\item When would an analytical solution become inconvenient enough to justify a numerical method?
\end{enumerate}\end{reflectionbox}
\chapter{Chemical mixture, tracer-drug and intravenous-feeding models}\label{app:A055}
\textbf{Application ID:} \pid{A055}\quad \textbf{Source locus:} 6.4\quad \textbf{Source pages:} 435, 437, 442
\section*{Application Context}
This family models a state whose derivative is governed by the state itself, external inputs, or coupled physical laws. Initial conditions select a unique trajectory when the model permits one. The canonical application family is \emph{Chemical mixture, tracer-drug and intravenous-feeding models}.
\section*{Calculus Toolkit}
\begin{tabularx}{\textwidth}{@{}lX@{}}\toprule
Primary technique & First-order linear compartment ODEs \\
Prerequisites & derivatives, separation or linear ODE methods, initial conditions \\
Numerical extension & Ch. 17 (initial-value ODEs); Ch. 18 for two-point boundary values \\
Analytical/numerical class & NUMERICAL-NATURAL \\
\bottomrule\end{tabularx}
\section*{Modeling Questions to Ask}\begin{itemize}
\item What is the state variable?
\item What law specifies its derivative?
\item What initial/boundary data are available?
\item Can the solution be checked by substitution?
\end{itemize}
\begin{errorbox}An ODE solution is not determined until initial or boundary data are applied. Verify the solution by substitution and check the initial condition.\end{errorbox}
\section*{Problems}
\needspace{8\baselineskip}
\subsection*{\pid{APP-055-F-01} — Foundation}
\begin{problembox}
A state variable in chemical mixture, tracer-drug and intravenous-feeding models satisfies $y'=0.11y$ with $y(0)=55$. Solve the initial-value problem and find $y(6)$.
\end{problembox}
\begin{promptbox}\begin{enumerate}
\item What is the state variable?
\item What law specifies its derivative?
\item What initial/boundary data are available?
\item Can the solution be checked by substitution?
\end{enumerate}\end{promptbox}
\textbf{Workspace}\workspace
\needspace{8\baselineskip}
\subsection*{\pid{APP-055-S-01} — Standard}
\begin{problembox}
A bounded-growth model for chemical mixture, tracer-drug and intravenous-feeding models satisfies $y'=0.32y(1-y/500)$ with $y(0)=55$. Use the logistic solution to estimate $y(8)$.
\end{problembox}
\begin{promptbox}\begin{enumerate}
\item What is the state variable?
\item What law specifies its derivative?
\item What initial/boundary data are available?
\item Can the solution be checked by substitution?
\end{enumerate}\end{promptbox}
\textbf{Workspace}\workspace
\needspace{8\baselineskip}
\subsection*{\pid{APP-055-A-01} — Advanced}
\begin{problembox}
A numerical extension of chemical mixture, tracer-drug and intravenous-feeding models uses $y'=-0.5y+10$, $y(0)=0$. Apply Euler's method with step $h=0.5$ to approximate $y(2)$, then compare with the exact solution $y=20(1-e^{-0.5t})$.
\end{problembox}
\begin{promptbox}\begin{enumerate}
\item What is the state variable?
\item What law specifies its derivative?
\item What initial/boundary data are available?
\item Can the solution be checked by substitution?
\end{enumerate}\end{promptbox}
\textbf{Workspace}\workspace
\begin{reflectionbox}\begin{enumerate}
\item Which sentence or datum in these problems supplied the main mathematical constraint?
\item Which assumptions were essential to the model, and how would the model change if one were removed?
\item When would an analytical solution become inconvenient enough to justify a numerical method?
\end{enumerate}\end{reflectionbox}
\chapter{Newton cooling, heat transfer and thermal transients}\label{app:A056}
\textbf{Application ID:} \pid{A056}\quad \textbf{Source locus:} 6.2\quad \textbf{Source pages:} 332, 419, 422
\section*{Application Context}
This family models a state whose derivative is governed by the state itself, external inputs, or coupled physical laws. Initial conditions select a unique trajectory when the model permits one. The canonical application family is \emph{Newton cooling, heat transfer and thermal transients}.
\section*{Calculus Toolkit}
\begin{tabularx}{\textwidth}{@{}lX@{}}\toprule
Primary technique & First-order thermal ODEs \\
Prerequisites & derivatives, separation or linear ODE methods, initial conditions \\
Numerical extension & Ch. 17 (initial-value ODEs); Ch. 18 for two-point boundary values \\
Analytical/numerical class & NUMERICAL-NATURAL \\
\bottomrule\end{tabularx}
\section*{Modeling Questions to Ask}\begin{itemize}
\item What is the state variable?
\item What law specifies its derivative?
\item What initial/boundary data are available?
\item Can the solution be checked by substitution?
\end{itemize}
\begin{errorbox}An ODE solution is not determined until initial or boundary data are applied. Verify the solution by substitution and check the initial condition.\end{errorbox}
\section*{Problems}
\needspace{8\baselineskip}
\subsection*{\pid{APP-056-F-01} — Foundation}
\begin{problembox}
A state variable in newton cooling, heat transfer and thermal transients satisfies $y'=0.08y$ with $y(0)=60$. Solve the initial-value problem and find $y(6)$.
\end{problembox}
\begin{promptbox}\begin{enumerate}
\item What is the state variable?
\item What law specifies its derivative?
\item What initial/boundary data are available?
\item Can the solution be checked by substitution?
\end{enumerate}\end{promptbox}
\textbf{Workspace}\workspace
\needspace{8\baselineskip}
\subsection*{\pid{APP-056-S-01} — Standard}
\begin{problembox}
A bounded-growth model for newton cooling, heat transfer and thermal transients satisfies $y'=0.34y(1-y/520)$ with $y(0)=40$. Use the logistic solution to estimate $y(8)$.
\end{problembox}
\begin{promptbox}\begin{enumerate}
\item What is the state variable?
\item What law specifies its derivative?
\item What initial/boundary data are available?
\item Can the solution be checked by substitution?
\end{enumerate}\end{promptbox}
\textbf{Workspace}\workspace
\needspace{8\baselineskip}
\subsection*{\pid{APP-056-A-01} — Advanced}
\begin{problembox}
A numerical extension of newton cooling, heat transfer and thermal transients uses $y'=-0.5y+10$, $y(0)=0$. Apply Euler's method with step $h=0.5$ to approximate $y(2)$, then compare with the exact solution $y=20(1-e^{-0.5t})$.
\end{problembox}
\begin{promptbox}\begin{enumerate}
\item What is the state variable?
\item What law specifies its derivative?
\item What initial/boundary data are available?
\item Can the solution be checked by substitution?
\end{enumerate}\end{promptbox}
\textbf{Workspace}\workspace
\begin{reflectionbox}\begin{enumerate}
\item Which sentence or datum in these problems supplied the main mathematical constraint?
\item Which assumptions were essential to the model, and how would the model change if one were removed?
\item When would an analytical solution become inconvenient enough to justify a numerical method?
\end{enumerate}\end{reflectionbox}
\chapter{Torricelli flow, emptying tanks and hydraulic transients}\label{app:A057}
\textbf{Application ID:} \pid{A057}\quad \textbf{Source locus:} 6; 7\quad \textbf{Source pages:} 441, 442, 489
\section*{Application Context}
This family models a state whose derivative is governed by the state itself, external inputs, or coupled physical laws. Initial conditions select a unique trajectory when the model permits one. The canonical application family is \emph{Torricelli flow, emptying tanks and hydraulic transients}.
\section*{Calculus Toolkit}
\begin{tabularx}{\textwidth}{@{}lX@{}}\toprule
Primary technique & Separable ODEs and fluid models \\
Prerequisites & derivatives, separation or linear ODE methods, initial conditions \\
Numerical extension & Ch. 17 (initial-value ODEs); Ch. 18 for two-point boundary values \\
Analytical/numerical class & NUMERICAL-NATURAL \\
\bottomrule\end{tabularx}
\section*{Modeling Questions to Ask}\begin{itemize}
\item What is the state variable?
\item What law specifies its derivative?
\item What initial/boundary data are available?
\item Can the solution be checked by substitution?
\end{itemize}
\begin{errorbox}An ODE solution is not determined until initial or boundary data are applied. Verify the solution by substitution and check the initial condition.\end{errorbox}
\section*{Problems}
\needspace{8\baselineskip}
\subsection*{\pid{APP-057-F-01} — Foundation}
\begin{problembox}
A state variable in torricelli flow, emptying tanks and hydraulic transients satisfies $y'=0.09y$ with $y(0)=65$. Solve the initial-value problem and find $y(6)$.
\end{problembox}
\begin{promptbox}\begin{enumerate}
\item What is the state variable?
\item What law specifies its derivative?
\item What initial/boundary data are available?
\item Can the solution be checked by substitution?
\end{enumerate}\end{promptbox}
\textbf{Workspace}\workspace
\needspace{8\baselineskip}
\subsection*{\pid{APP-057-S-01} — Standard}
\begin{problembox}
A bounded-growth model for torricelli flow, emptying tanks and hydraulic transients satisfies $y'=0.30y(1-y/540)$ with $y(0)=45$. Use the logistic solution to estimate $y(8)$.
\end{problembox}
\begin{promptbox}\begin{enumerate}
\item What is the state variable?
\item What law specifies its derivative?
\item What initial/boundary data are available?
\item Can the solution be checked by substitution?
\end{enumerate}\end{promptbox}
\textbf{Workspace}\workspace
\needspace{8\baselineskip}
\subsection*{\pid{APP-057-A-01} — Advanced}
\begin{problembox}
A numerical extension of torricelli flow, emptying tanks and hydraulic transients uses $y'=-0.5y+10$, $y(0)=0$. Apply Euler's method with step $h=0.5$ to approximate $y(2)$, then compare with the exact solution $y=20(1-e^{-0.5t})$.
\end{problembox}
\begin{promptbox}\begin{enumerate}
\item What is the state variable?
\item What law specifies its derivative?
\item What initial/boundary data are available?
\item Can the solution be checked by substitution?
\end{enumerate}\end{promptbox}
\textbf{Workspace}\workspace
\begin{reflectionbox}\begin{enumerate}
\item Which sentence or datum in these problems supplied the main mathematical constraint?
\item Which assumptions were essential to the model, and how would the model change if one were removed?
\item When would an analytical solution become inconvenient enough to justify a numerical method?
\end{enumerate}\end{reflectionbox}
\chapter{Electric circuits, resistance/inductance dynamics and control}\label{app:A058}
\textbf{Application ID:} \pid{A058}\quad \textbf{Source locus:} 5–6; 13\quad \textbf{Source pages:} 371, 414, 434, 437, 910
\section*{Application Context}
This family models a state whose derivative is governed by the state itself, external inputs, or coupled physical laws. Initial conditions select a unique trajectory when the model permits one. The canonical application family is \emph{Electric circuits, resistance/inductance dynamics and control}.
\section*{Calculus Toolkit}
\begin{tabularx}{\textwidth}{@{}lX@{}}\toprule
Primary technique & First-order and coupled circuit ODEs \\
Prerequisites & derivatives, separation or linear ODE methods, initial conditions \\
Numerical extension & Ch. 17 (initial-value ODEs); Ch. 18 for two-point boundary values \\
Analytical/numerical class & NUMERICAL-NATURAL \\
\bottomrule\end{tabularx}
\section*{Modeling Questions to Ask}\begin{itemize}
\item What is the state variable?
\item What law specifies its derivative?
\item What initial/boundary data are available?
\item Can the solution be checked by substitution?
\end{itemize}
\begin{errorbox}An ODE solution is not determined until initial or boundary data are applied. Verify the solution by substitution and check the initial condition.\end{errorbox}
\section*{Problems}
\needspace{8\baselineskip}
\subsection*{\pid{APP-058-F-01} — Foundation}
\begin{problembox}
A state variable in electric circuits, resistance/inductance dynamics and control satisfies $y'=0.10y$ with $y(0)=70$. Solve the initial-value problem and find $y(6)$.
\end{problembox}
\begin{promptbox}\begin{enumerate}
\item What is the state variable?
\item What law specifies its derivative?
\item What initial/boundary data are available?
\item Can the solution be checked by substitution?
\end{enumerate}\end{promptbox}
\textbf{Workspace}\workspace
\needspace{8\baselineskip}
\subsection*{\pid{APP-058-S-01} — Standard}
\begin{problembox}
A bounded-growth model for electric circuits, resistance/inductance dynamics and control satisfies $y'=0.32y(1-y/560)$ with $y(0)=50$. Use the logistic solution to estimate $y(8)$.
\end{problembox}
\begin{promptbox}\begin{enumerate}
\item What is the state variable?
\item What law specifies its derivative?
\item What initial/boundary data are available?
\item Can the solution be checked by substitution?
\end{enumerate}\end{promptbox}
\textbf{Workspace}\workspace
\needspace{8\baselineskip}
\subsection*{\pid{APP-058-A-01} — Advanced}
\begin{problembox}
A numerical extension of electric circuits, resistance/inductance dynamics and control uses $y'=-0.5y+10$, $y(0)=0$. Apply Euler's method with step $h=0.5$ to approximate $y(2)$, then compare with the exact solution $y=20(1-e^{-0.5t})$.
\end{problembox}
\begin{promptbox}\begin{enumerate}
\item What is the state variable?
\item What law specifies its derivative?
\item What initial/boundary data are available?
\item Can the solution be checked by substitution?
\end{enumerate}\end{promptbox}
\textbf{Workspace}\workspace
\begin{reflectionbox}\begin{enumerate}
\item Which sentence or datum in these problems supplied the main mathematical constraint?
\item Which assumptions were essential to the model, and how would the model change if one were removed?
\item When would an analytical solution become inconvenient enough to justify a numerical method?
\end{enumerate}\end{reflectionbox}
\chapter{Spring motion, pendulum and harmonic motion}\label{app:A059}
\textbf{Application ID:} \pid{A059}\quad \textbf{Source locus:} 2; 6; 12\quad \textbf{Source pages:} 142, 162, 241, 349, 531, 910
\section*{Application Context}
This family models a state whose derivative is governed by the state itself, external inputs, or coupled physical laws. Initial conditions select a unique trajectory when the model permits one. The canonical application family is \emph{Spring motion, pendulum and harmonic motion}.
\section*{Calculus Toolkit}
\begin{tabularx}{\textwidth}{@{}lX@{}}\toprule
Primary technique & Second-order ODEs; oscillation \\
Prerequisites & derivatives, separation or linear ODE methods, initial conditions \\
Numerical extension & Ch. 17 (initial-value ODEs); Ch. 18 for two-point boundary values \\
Analytical/numerical class & NUMERICAL-NATURAL \\
\bottomrule\end{tabularx}
\section*{Modeling Questions to Ask}\begin{itemize}
\item What is the state variable?
\item What law specifies its derivative?
\item What initial/boundary data are available?
\item Can the solution be checked by substitution?
\end{itemize}
\begin{errorbox}An ODE solution is not determined until initial or boundary data are applied. Verify the solution by substitution and check the initial condition.\end{errorbox}
\section*{Problems}
\needspace{8\baselineskip}
\subsection*{\pid{APP-059-F-01} — Foundation}
\begin{problembox}
A state variable in spring motion, pendulum and harmonic motion satisfies $y'=0.11y$ with $y(0)=75$. Solve the initial-value problem and find $y(6)$.
\end{problembox}
\begin{promptbox}\begin{enumerate}
\item What is the state variable?
\item What law specifies its derivative?
\item What initial/boundary data are available?
\item Can the solution be checked by substitution?
\end{enumerate}\end{promptbox}
\textbf{Workspace}\workspace
\needspace{8\baselineskip}
\subsection*{\pid{APP-059-S-01} — Standard}
\begin{problembox}
A bounded-growth model for spring motion, pendulum and harmonic motion satisfies $y'=0.34y(1-y/580)$ with $y(0)=55$. Use the logistic solution to estimate $y(8)$.
\end{problembox}
\begin{promptbox}\begin{enumerate}
\item What is the state variable?
\item What law specifies its derivative?
\item What initial/boundary data are available?
\item Can the solution be checked by substitution?
\end{enumerate}\end{promptbox}
\textbf{Workspace}\workspace
\needspace{8\baselineskip}
\subsection*{\pid{APP-059-A-01} — Advanced}
\begin{problembox}
A numerical extension of spring motion, pendulum and harmonic motion uses $y'=-0.5y+10$, $y(0)=0$. Apply Euler's method with step $h=0.5$ to approximate $y(2)$, then compare with the exact solution $y=20(1-e^{-0.5t})$.
\end{problembox}
\begin{promptbox}\begin{enumerate}
\item What is the state variable?
\item What law specifies its derivative?
\item What initial/boundary data are available?
\item Can the solution be checked by substitution?
\end{enumerate}\end{promptbox}
\textbf{Workspace}\workspace
\begin{reflectionbox}\begin{enumerate}
\item Which sentence or datum in these problems supplied the main mathematical constraint?
\item Which assumptions were essential to the model, and how would the model change if one were removed?
\item When would an analytical solution become inconvenient enough to justify a numerical method?
\end{enumerate}\end{reflectionbox}
\chapter{Falling bodies and velocity in a resisting medium}\label{app:A060}
\textbf{Application ID:} \pid{A060}\quad \textbf{Source locus:} 1; 6\quad \textbf{Source pages:} 73, 95, 311, 371, 434, 437
\section*{Application Context}
This family models a state whose derivative is governed by the state itself, external inputs, or coupled physical laws. Initial conditions select a unique trajectory when the model permits one. The canonical application family is \emph{Falling bodies and velocity in a resisting medium}.
\section*{Calculus Toolkit}
\begin{tabularx}{\textwidth}{@{}lX@{}}\toprule
Primary technique & Initial-value ODEs with drag \\
Prerequisites & derivatives, separation or linear ODE methods, initial conditions \\
Numerical extension & Ch. 17 (initial-value ODEs); Ch. 18 for two-point boundary values \\
Analytical/numerical class & NUMERICAL-NATURAL \\
\bottomrule\end{tabularx}
\section*{Modeling Questions to Ask}\begin{itemize}
\item What is the state variable?
\item What law specifies its derivative?
\item What initial/boundary data are available?
\item Can the solution be checked by substitution?
\end{itemize}
\begin{errorbox}An ODE solution is not determined until initial or boundary data are applied. Verify the solution by substitution and check the initial condition.\end{errorbox}
\section*{Problems}
\needspace{8\baselineskip}
\subsection*{\pid{APP-060-F-01} — Foundation}
\begin{problembox}
A state variable in falling bodies and velocity in a resisting medium satisfies $y'=0.08y$ with $y(0)=50$. Solve the initial-value problem and find $y(6)$.
\end{problembox}
\begin{promptbox}\begin{enumerate}
\item What is the state variable?
\item What law specifies its derivative?
\item What initial/boundary data are available?
\item Can the solution be checked by substitution?
\end{enumerate}\end{promptbox}
\textbf{Workspace}\workspace
\needspace{8\baselineskip}
\subsection*{\pid{APP-060-S-01} — Standard}
\begin{problembox}
A bounded-growth model for falling bodies and velocity in a resisting medium satisfies $y'=0.30y(1-y/500)$ with $y(0)=40$. Use the logistic solution to estimate $y(8)$.
\end{problembox}
\begin{promptbox}\begin{enumerate}
\item What is the state variable?
\item What law specifies its derivative?
\item What initial/boundary data are available?
\item Can the solution be checked by substitution?
\end{enumerate}\end{promptbox}
\textbf{Workspace}\workspace
\needspace{8\baselineskip}
\subsection*{\pid{APP-060-A-01} — Advanced}
\begin{problembox}
A numerical extension of falling bodies and velocity in a resisting medium uses $y'=-0.5y+10$, $y(0)=0$. Apply Euler's method with step $h=0.5$ to approximate $y(2)$, then compare with the exact solution $y=20(1-e^{-0.5t})$.
\end{problembox}
\begin{promptbox}\begin{enumerate}
\item What is the state variable?
\item What law specifies its derivative?
\item What initial/boundary data are available?
\item Can the solution be checked by substitution?
\end{enumerate}\end{promptbox}
\textbf{Workspace}\workspace
\begin{reflectionbox}\begin{enumerate}
\item Which sentence or datum in these problems supplied the main mathematical constraint?
\item Which assumptions were essential to the model, and how would the model change if one were removed?
\item When would an analytical solution become inconvenient enough to justify a numerical method?
\end{enumerate}\end{reflectionbox}
\chapter{Projectile trajectories, baseball, football, and shot paths; tossing bales}\label{app:A061}
\textbf{Application ID:} \pid{A061}\quad \textbf{Source locus:} 3; 10; 12–13\quad \textbf{Source pages:} 164, 186, 241, 679, 706, 709, 716, 761, 840–869, 917, 968
\section*{Application Context}
This family models a state whose derivative is governed by the state itself, external inputs, or coupled physical laws. Initial conditions select a unique trajectory when the model permits one. The canonical application family is \emph{Projectile trajectories, baseball, football, and shot paths; tossing bales}.
\section*{Calculus Toolkit}
\begin{tabularx}{\textwidth}{@{}lX@{}}\toprule
Primary technique & Parametric/vector ODE trajectory models \\
Prerequisites & derivatives, separation or linear ODE methods, initial conditions \\
Numerical extension & Ch. 17 (initial-value ODEs); Ch. 18 for two-point boundary values \\
Analytical/numerical class & NUMERICAL-NATURAL \\
\bottomrule\end{tabularx}
\section*{Modeling Questions to Ask}\begin{itemize}
\item What is the state variable?
\item What law specifies its derivative?
\item What initial/boundary data are available?
\item Can the solution be checked by substitution?
\end{itemize}
\begin{errorbox}An ODE solution is not determined until initial or boundary data are applied. Verify the solution by substitution and check the initial condition.\end{errorbox}
\section*{Problems}
\needspace{8\baselineskip}
\subsection*{\pid{APP-061-F-01} — Foundation}
\begin{problembox}
A state variable in projectile trajectories, baseball, football, and shot paths; tossing bales satisfies $y'=0.09y$ with $y(0)=55$. Solve the initial-value problem and find $y(6)$.
\end{problembox}
\begin{promptbox}\begin{enumerate}
\item What is the state variable?
\item What law specifies its derivative?
\item What initial/boundary data are available?
\item Can the solution be checked by substitution?
\end{enumerate}\end{promptbox}
\textbf{Workspace}\workspace
\needspace{8\baselineskip}
\subsection*{\pid{APP-061-S-01} — Standard}
\begin{problembox}
A bounded-growth model for projectile trajectories, baseball, football, and shot paths; tossing bales satisfies $y'=0.32y(1-y/520)$ with $y(0)=45$. Use the logistic solution to estimate $y(8)$.
\end{problembox}
\begin{promptbox}\begin{enumerate}
\item What is the state variable?
\item What law specifies its derivative?
\item What initial/boundary data are available?
\item Can the solution be checked by substitution?
\end{enumerate}\end{promptbox}
\textbf{Workspace}\workspace
\needspace{8\baselineskip}
\subsection*{\pid{APP-061-A-01} — Advanced}
\begin{problembox}
A numerical extension of projectile trajectories, baseball, football, and shot paths; tossing bales uses $y'=-0.5y+10$, $y(0)=0$. Apply Euler's method with step $h=0.5$ to approximate $y(2)$, then compare with the exact solution $y=20(1-e^{-0.5t})$.
\end{problembox}
\begin{promptbox}\begin{enumerate}
\item What is the state variable?
\item What law specifies its derivative?
\item What initial/boundary data are available?
\item Can the solution be checked by substitution?
\end{enumerate}\end{promptbox}
\textbf{Workspace}\workspace
\begin{reflectionbox}\begin{enumerate}
\item Which sentence or datum in these problems supplied the main mathematical constraint?
\item Which assumptions were essential to the model, and how would the model change if one were removed?
\item When would an analytical solution become inconvenient enough to justify a numerical method?
\end{enumerate}\end{reflectionbox}
\chapter{Orbital motion, satellites, planetary motion and Kepler-law simulation}\label{app:A062}
\textbf{Application ID:} \pid{A062}\quad \textbf{Source locus:} 10–12\quad \textbf{Source pages:} 690–698, 731, 741–745, 854, 866, 870
\section*{Application Context}
This family models a state whose derivative is governed by the state itself, external inputs, or coupled physical laws. Initial conditions select a unique trajectory when the model permits one. The canonical application family is \emph{Orbital motion, satellites, planetary motion and Kepler-law simulation}.
\section*{Calculus Toolkit}
\begin{tabularx}{\textwidth}{@{}lX@{}}\toprule
Primary technique & Second-order gravitational ODEs \\
Prerequisites & derivatives, separation or linear ODE methods, initial conditions \\
Numerical extension & Ch. 17 (initial-value ODEs); Ch. 18 for two-point boundary values \\
Analytical/numerical class & NUMERICAL-NATURAL \\
\bottomrule\end{tabularx}
\section*{Modeling Questions to Ask}\begin{itemize}
\item What is the state variable?
\item What law specifies its derivative?
\item What initial/boundary data are available?
\item Can the solution be checked by substitution?
\end{itemize}
\begin{errorbox}An ODE solution is not determined until initial or boundary data are applied. Verify the solution by substitution and check the initial condition.\end{errorbox}
\section*{Problems}
\needspace{8\baselineskip}
\subsection*{\pid{APP-062-F-01} — Foundation}
\begin{problembox}
A state variable in orbital motion, satellites, planetary motion and kepler-law simulation satisfies $y'=0.10y$ with $y(0)=60$. Solve the initial-value problem and find $y(6)$.
\end{problembox}
\begin{promptbox}\begin{enumerate}
\item What is the state variable?
\item What law specifies its derivative?
\item What initial/boundary data are available?
\item Can the solution be checked by substitution?
\end{enumerate}\end{promptbox}
\textbf{Workspace}\workspace
\needspace{8\baselineskip}
\subsection*{\pid{APP-062-S-01} — Standard}
\begin{problembox}
A bounded-growth model for orbital motion, satellites, planetary motion and kepler-law simulation satisfies $y'=0.34y(1-y/540)$ with $y(0)=50$. Use the logistic solution to estimate $y(8)$.
\end{problembox}
\begin{promptbox}\begin{enumerate}
\item What is the state variable?
\item What law specifies its derivative?
\item What initial/boundary data are available?
\item Can the solution be checked by substitution?
\end{enumerate}\end{promptbox}
\textbf{Workspace}\workspace
\needspace{8\baselineskip}
\subsection*{\pid{APP-062-A-01} — Advanced}
\begin{problembox}
A numerical extension of orbital motion, satellites, planetary motion and kepler-law simulation uses $y'=-0.5y+10$, $y(0)=0$. Apply Euler's method with step $h=0.5$ to approximate $y(2)$, then compare with the exact solution $y=20(1-e^{-0.5t})$.
\end{problembox}
\begin{promptbox}\begin{enumerate}
\item What is the state variable?
\item What law specifies its derivative?
\item What initial/boundary data are available?
\item Can the solution be checked by substitution?
\end{enumerate}\end{promptbox}
\textbf{Workspace}\workspace
\begin{reflectionbox}\begin{enumerate}
\item Which sentence or datum in these problems supplied the main mathematical constraint?
\item Which assumptions were essential to the model, and how would the model change if one were removed?
\item When would an analytical solution become inconvenient enough to justify a numerical method?
\end{enumerate}\end{reflectionbox}
\chapter{Heat-seeking particle/path and pursuit dynamics}\label{app:A063}
\textbf{Application ID:} \pid{A063}\quad \textbf{Source locus:} 13.6–13.7\quad \textbf{Source pages:} 925, 930
\section*{Application Context}
This family models a state whose derivative is governed by the state itself, external inputs, or coupled physical laws. Initial conditions select a unique trajectory when the model permits one. The canonical application family is \emph{Heat-seeking particle/path and pursuit dynamics}.
\section*{Calculus Toolkit}
\begin{tabularx}{\textwidth}{@{}lX@{}}\toprule
Primary technique & Gradient-driven ODE path \\
Prerequisites & derivatives, separation or linear ODE methods, initial conditions \\
Numerical extension & Ch. 17 (initial-value ODEs); Ch. 18 for two-point boundary values \\
Analytical/numerical class & NUMERICAL-NATURAL \\
\bottomrule\end{tabularx}
\section*{Modeling Questions to Ask}\begin{itemize}
\item What is the state variable?
\item What law specifies its derivative?
\item What initial/boundary data are available?
\item Can the solution be checked by substitution?
\end{itemize}
\begin{errorbox}An ODE solution is not determined until initial or boundary data are applied. Verify the solution by substitution and check the initial condition.\end{errorbox}
\section*{Problems}
\needspace{8\baselineskip}
\subsection*{\pid{APP-063-F-01} — Foundation}
\begin{problembox}
A state variable in heat-seeking particle/path and pursuit dynamics satisfies $y'=0.11y$ with $y(0)=65$. Solve the initial-value problem and find $y(6)$.
\end{problembox}
\begin{promptbox}\begin{enumerate}
\item What is the state variable?
\item What law specifies its derivative?
\item What initial/boundary data are available?
\item Can the solution be checked by substitution?
\end{enumerate}\end{promptbox}
\textbf{Workspace}\workspace
\needspace{8\baselineskip}
\subsection*{\pid{APP-063-S-01} — Standard}
\begin{problembox}
A bounded-growth model for heat-seeking particle/path and pursuit dynamics satisfies $y'=0.30y(1-y/560)$ with $y(0)=55$. Use the logistic solution to estimate $y(8)$.
\end{problembox}
\begin{promptbox}\begin{enumerate}
\item What is the state variable?
\item What law specifies its derivative?
\item What initial/boundary data are available?
\item Can the solution be checked by substitution?
\end{enumerate}\end{promptbox}
\textbf{Workspace}\workspace
\needspace{8\baselineskip}
\subsection*{\pid{APP-063-A-01} — Advanced}
\begin{problembox}
A numerical extension of heat-seeking particle/path and pursuit dynamics uses $y'=-0.5y+10$, $y(0)=0$. Apply Euler's method with step $h=0.5$ to approximate $y(2)$, then compare with the exact solution $y=20(1-e^{-0.5t})$.
\end{problembox}
\begin{promptbox}\begin{enumerate}
\item What is the state variable?
\item What law specifies its derivative?
\item What initial/boundary data are available?
\item Can the solution be checked by substitution?
\end{enumerate}\end{promptbox}
\textbf{Workspace}\workspace
\begin{reflectionbox}\begin{enumerate}
\item Which sentence or datum in these problems supplied the main mathematical constraint?
\item Which assumptions were essential to the model, and how would the model change if one were removed?
\item When would an analytical solution become inconvenient enough to justify a numerical method?
\end{enumerate}\end{reflectionbox}
\part{Data, Statistics, and Probability}
\chapter{Sales, inflation, health-care expenditures, government expenditures and GDP data}\label{app:A064}
\textbf{Application ID:} \pid{A064}\quad \textbf{Source locus:} P; 5; 13\quad \textbf{Source pages:} 9, 131, 359, 597, 883, 903
\section*{Application Context}
This family begins with measured or simulated data and uses statistics, regression, probability, or Monte Carlo techniques to quantify patterns and uncertainty. The canonical application family is \emph{Sales, inflation, health-care expenditures, government expenditures and GDP data}.
\section*{Calculus Toolkit}
\begin{tabularx}{\textwidth}{@{}lX@{}}\toprule
Primary technique & Regression, trend estimation and uncertainty \\
Prerequisites & functions, descriptive statistics, model fitting \\
Numerical extension & Ch. 14 (statistics), Ch. 15 (modeling), Ch. 7 (random numbers) \\
Analytical/numerical class & NUMERICAL-NATURAL \\
\bottomrule\end{tabularx}
\section*{Modeling Questions to Ask}\begin{itemize}
\item What are the observations and response variable?
\item What model class is being assumed?
\item What statistic or fitted parameter is requested?
\item How should residuals or uncertainty be interpreted?
\end{itemize}
\begin{errorbox}Do not treat a fitted relationship as an exact law. Distinguish observations, model predictions, residuals, and uncertainty.\end{errorbox}
\section*{Problems}
\needspace{8\baselineskip}
\subsection*{\pid{APP-064-F-01} — Foundation}
\begin{problembox}
Five observations in sales, inflation, health-care expenditures, government expenditures and gdp data are [14, 16, 19, 17, 24]. Compute the sample mean and sample variance.
\end{problembox}
\begin{promptbox}\begin{enumerate}
\item What are the observations and response variable?
\item What model class is being assumed?
\item What statistic or fitted parameter is requested?
\item How should residuals or uncertainty be interpreted?
\end{enumerate}\end{promptbox}
\textbf{Workspace}\workspace
\needspace{8\baselineskip}
\subsection*{\pid{APP-064-S-01} — Standard}
\begin{problembox}
A small dataset for sales, inflation, health-care expenditures, government expenditures and gdp data has points $(0,6),(1,8),(2,10),(3,13)$. Find the least-squares line $y=mx+b$.
\end{problembox}
\begin{promptbox}\begin{enumerate}
\item What are the observations and response variable?
\item What model class is being assumed?
\item What statistic or fitted parameter is requested?
\item How should residuals or uncertainty be interpreted?
\end{enumerate}\end{promptbox}
\textbf{Workspace}\workspace
\needspace{8\baselineskip}
\subsection*{\pid{APP-064-A-01} — Advanced}
\begin{problembox}
For a simple uncertainty exercise related to sales, inflation, health-care expenditures, government expenditures and gdp data, ten uniform samples on $[0,1]$ are [0.11, 0.27, 0.33, 0.48, 0.52, 0.61, 0.73, 0.81, 0.89, 0.95]. Estimate $P(X^2<0.5)$ by Monte Carlo frequency and compare with the exact probability.
\end{problembox}
\begin{promptbox}\begin{enumerate}
\item What are the observations and response variable?
\item What model class is being assumed?
\item What statistic or fitted parameter is requested?
\item How should residuals or uncertainty be interpreted?
\end{enumerate}\end{promptbox}
\textbf{Workspace}\workspace
\begin{reflectionbox}\begin{enumerate}
\item Which sentence or datum in these problems supplied the main mathematical constraint?
\item Which assumptions were essential to the model, and how would the model change if one were removed?
\item When would an analytical solution become inconvenient enough to justify a numerical method?
\end{enumerate}\end{reflectionbox}
\chapter{Systolic blood pressure, intelligence-test and physiological data}\label{app:A065}
\textbf{Application ID:} \pid{A065}\quad \textbf{Source locus:} 13\quad \textbf{Source pages:} 902, 954
\section*{Application Context}
This family begins with measured or simulated data and uses statistics, regression, probability, or Monte Carlo techniques to quantify patterns and uncertainty. The canonical application family is \emph{Systolic blood pressure, intelligence-test and physiological data}.
\section*{Calculus Toolkit}
\begin{tabularx}{\textwidth}{@{}lX@{}}\toprule
Primary technique & Statistical description and fitted models \\
Prerequisites & functions, descriptive statistics, model fitting \\
Numerical extension & Ch. 14 (statistics), Ch. 15 (modeling), Ch. 7 (random numbers) \\
Analytical/numerical class & NUMERICAL-NATURAL \\
\bottomrule\end{tabularx}
\section*{Modeling Questions to Ask}\begin{itemize}
\item What are the observations and response variable?
\item What model class is being assumed?
\item What statistic or fitted parameter is requested?
\item How should residuals or uncertainty be interpreted?
\end{itemize}
\begin{errorbox}Do not treat a fitted relationship as an exact law. Distinguish observations, model predictions, residuals, and uncertainty.\end{errorbox}
\section*{Problems}
\needspace{8\baselineskip}
\subsection*{\pid{APP-065-F-01} — Foundation}
\begin{problembox}
Five observations in systolic blood pressure, intelligence-test and physiological data are [10, 12, 15, 13, 20]. Compute the sample mean and sample variance.
\end{problembox}
\begin{promptbox}\begin{enumerate}
\item What are the observations and response variable?
\item What model class is being assumed?
\item What statistic or fitted parameter is requested?
\item How should residuals or uncertainty be interpreted?
\end{enumerate}\end{promptbox}
\textbf{Workspace}\workspace
\needspace{8\baselineskip}
\subsection*{\pid{APP-065-S-01} — Standard}
\begin{problembox}
A small dataset for systolic blood pressure, intelligence-test and physiological data has points $(0,7),(1,9),(2,11),(3,14)$. Find the least-squares line $y=mx+b$.
\end{problembox}
\begin{promptbox}\begin{enumerate}
\item What are the observations and response variable?
\item What model class is being assumed?
\item What statistic or fitted parameter is requested?
\item How should residuals or uncertainty be interpreted?
\end{enumerate}\end{promptbox}
\textbf{Workspace}\workspace
\needspace{8\baselineskip}
\subsection*{\pid{APP-065-A-01} — Advanced}
\begin{problembox}
For a simple uncertainty exercise related to systolic blood pressure, intelligence-test and physiological data, ten uniform samples on $[0,1]$ are [0.11, 0.27, 0.33, 0.48, 0.52, 0.61, 0.73, 0.81, 0.89, 0.95]. Estimate $P(X^2<0.5)$ by Monte Carlo frequency and compare with the exact probability.
\end{problembox}
\begin{promptbox}\begin{enumerate}
\item What are the observations and response variable?
\item What model class is being assumed?
\item What statistic or fitted parameter is requested?
\item How should residuals or uncertainty be interpreted?
\end{enumerate}\end{promptbox}
\textbf{Workspace}\workspace
\begin{reflectionbox}\begin{enumerate}
\item Which sentence or datum in these problems supplied the main mathematical constraint?
\item Which assumptions were essential to the model, and how would the model change if one were removed?
\item When would an analytical solution become inconvenient enough to justify a numerical method?
\end{enumerate}\end{reflectionbox}
\chapter{Learning curve and memory model}\label{app:A066}
\textbf{Application ID:} \pid{A066}\quad \textbf{Source locus:} 6–7\quad \textbf{Source pages:} 421, 422, 436, 531
\section*{Application Context}
This family begins with measured or simulated data and uses statistics, regression, probability, or Monte Carlo techniques to quantify patterns and uncertainty. The canonical application family is \emph{Learning curve and memory model}.
\section*{Calculus Toolkit}
\begin{tabularx}{\textwidth}{@{}lX@{}}\toprule
Primary technique & Nonlinear model fitting \\
Prerequisites & functions, descriptive statistics, model fitting \\
Numerical extension & Ch. 14 (statistics), Ch. 15 (modeling), Ch. 7 (random numbers) \\
Analytical/numerical class & NUMERICAL-NATURAL \\
\bottomrule\end{tabularx}
\section*{Modeling Questions to Ask}\begin{itemize}
\item What are the observations and response variable?
\item What model class is being assumed?
\item What statistic or fitted parameter is requested?
\item How should residuals or uncertainty be interpreted?
\end{itemize}
\begin{errorbox}Do not treat a fitted relationship as an exact law. Distinguish observations, model predictions, residuals, and uncertainty.\end{errorbox}
\section*{Problems}
\needspace{8\baselineskip}
\subsection*{\pid{APP-066-F-01} — Foundation}
\begin{problembox}
Five observations in learning curve and memory model are [11, 13, 16, 14, 21]. Compute the sample mean and sample variance.
\end{problembox}
\begin{promptbox}\begin{enumerate}
\item What are the observations and response variable?
\item What model class is being assumed?
\item What statistic or fitted parameter is requested?
\item How should residuals or uncertainty be interpreted?
\end{enumerate}\end{promptbox}
\textbf{Workspace}\workspace
\needspace{8\baselineskip}
\subsection*{\pid{APP-066-S-01} — Standard}
\begin{problembox}
A small dataset for learning curve and memory model has points $(0,5),(1,7),(2,9),(3,12)$. Find the least-squares line $y=mx+b$.
\end{problembox}
\begin{promptbox}\begin{enumerate}
\item What are the observations and response variable?
\item What model class is being assumed?
\item What statistic or fitted parameter is requested?
\item How should residuals or uncertainty be interpreted?
\end{enumerate}\end{promptbox}
\textbf{Workspace}\workspace
\needspace{8\baselineskip}
\subsection*{\pid{APP-066-A-01} — Advanced}
\begin{problembox}
For a simple uncertainty exercise related to learning curve and memory model, ten uniform samples on $[0,1]$ are [0.11, 0.27, 0.33, 0.48, 0.52, 0.61, 0.73, 0.81, 0.89, 0.95]. Estimate $P(X^2<0.5)$ by Monte Carlo frequency and compare with the exact probability.
\end{problembox}
\begin{promptbox}\begin{enumerate}
\item What are the observations and response variable?
\item What model class is being assumed?
\item What statistic or fitted parameter is requested?
\item How should residuals or uncertainty be interpreted?
\end{enumerate}\end{promptbox}
\textbf{Workspace}\workspace
\begin{reflectionbox}\begin{enumerate}
\item Which sentence or datum in these problems supplied the main mathematical constraint?
\item Which assumptions were essential to the model, and how would the model change if one were removed?
\item When would an analytical solution become inconvenient enough to justify a numerical method?
\end{enumerate}\end{reflectionbox}
\chapter{Average typing speed and experimental-rate datasets}\label{app:A067}
\textbf{Application ID:} \pid{A067}\quad \textbf{Source locus:} 3\quad \textbf{Source pages:} 198, 207
\section*{Application Context}
This family begins with measured or simulated data and uses statistics, regression, probability, or Monte Carlo techniques to quantify patterns and uncertainty. The canonical application family is \emph{Average typing speed and experimental-rate datasets}.
\section*{Calculus Toolkit}
\begin{tabularx}{\textwidth}{@{}lX@{}}\toprule
Primary technique & Statistics and fitted rates \\
Prerequisites & functions, descriptive statistics, model fitting \\
Numerical extension & Ch. 14 (statistics), Ch. 15 (modeling), Ch. 7 (random numbers) \\
Analytical/numerical class & NUMERICAL-NATURAL \\
\bottomrule\end{tabularx}
\section*{Modeling Questions to Ask}\begin{itemize}
\item What are the observations and response variable?
\item What model class is being assumed?
\item What statistic or fitted parameter is requested?
\item How should residuals or uncertainty be interpreted?
\end{itemize}
\begin{errorbox}Do not treat a fitted relationship as an exact law. Distinguish observations, model predictions, residuals, and uncertainty.\end{errorbox}
\section*{Problems}
\needspace{8\baselineskip}
\subsection*{\pid{APP-067-F-01} — Foundation}
\begin{problembox}
Five observations in average typing speed and experimental-rate datasets are [12, 14, 17, 15, 22]. Compute the sample mean and sample variance.
\end{problembox}
\begin{promptbox}\begin{enumerate}
\item What are the observations and response variable?
\item What model class is being assumed?
\item What statistic or fitted parameter is requested?
\item How should residuals or uncertainty be interpreted?
\end{enumerate}\end{promptbox}
\textbf{Workspace}\workspace
\needspace{8\baselineskip}
\subsection*{\pid{APP-067-S-01} — Standard}
\begin{problembox}
A small dataset for average typing speed and experimental-rate datasets has points $(0,6),(1,8),(2,10),(3,13)$. Find the least-squares line $y=mx+b$.
\end{problembox}
\begin{promptbox}\begin{enumerate}
\item What are the observations and response variable?
\item What model class is being assumed?
\item What statistic or fitted parameter is requested?
\item How should residuals or uncertainty be interpreted?
\end{enumerate}\end{promptbox}
\textbf{Workspace}\workspace
\needspace{8\baselineskip}
\subsection*{\pid{APP-067-A-01} — Advanced}
\begin{problembox}
For a simple uncertainty exercise related to average typing speed and experimental-rate datasets, ten uniform samples on $[0,1]$ are [0.11, 0.27, 0.33, 0.48, 0.52, 0.61, 0.73, 0.81, 0.89, 0.95]. Estimate $P(X^2<0.5)$ by Monte Carlo frequency and compare with the exact probability.
\end{problembox}
\begin{promptbox}\begin{enumerate}
\item What are the observations and response variable?
\item What model class is being assumed?
\item What statistic or fitted parameter is requested?
\item How should residuals or uncertainty be interpreted?
\end{enumerate}\end{promptbox}
\textbf{Workspace}\workspace
\begin{reflectionbox}\begin{enumerate}
\item Which sentence or datum in these problems supplied the main mathematical constraint?
\item Which assumptions were essential to the model, and how would the model change if one were removed?
\item When would an analytical solution become inconvenient enough to justify a numerical method?
\end{enumerate}\end{reflectionbox}
\chapter{Agronomy, forestry, biomass and ecological field data}\label{app:A068}
\textbf{Application ID:} \pid{A068}\quad \textbf{Source locus:} 6; 13\quad \textbf{Source pages:} 422, 442, 882, 966
\section*{Application Context}
This family begins with measured or simulated data and uses statistics, regression, probability, or Monte Carlo techniques to quantify patterns and uncertainty. The canonical application family is \emph{Agronomy, forestry, biomass and ecological field data}.
\section*{Calculus Toolkit}
\begin{tabularx}{\textwidth}{@{}lX@{}}\toprule
Primary technique & Regression, interpolation, and statistics \\
Prerequisites & functions, descriptive statistics, model fitting \\
Numerical extension & Ch. 14 (statistics), Ch. 15 (modeling), Ch. 7 (random numbers) \\
Analytical/numerical class & NUMERICAL-NATURAL \\
\bottomrule\end{tabularx}
\section*{Modeling Questions to Ask}\begin{itemize}
\item What are the observations and response variable?
\item What model class is being assumed?
\item What statistic or fitted parameter is requested?
\item How should residuals or uncertainty be interpreted?
\end{itemize}
\begin{errorbox}Do not treat a fitted relationship as an exact law. Distinguish observations, model predictions, residuals, and uncertainty.\end{errorbox}
\section*{Problems}
\needspace{8\baselineskip}
\subsection*{\pid{APP-068-F-01} — Foundation}
\begin{problembox}
Five observations in agronomy, forestry, biomass and ecological field data are [13, 15, 18, 16, 23]. Compute the sample mean and sample variance.
\end{problembox}
\begin{promptbox}\begin{enumerate}
\item What are the observations and response variable?
\item What model class is being assumed?
\item What statistic or fitted parameter is requested?
\item How should residuals or uncertainty be interpreted?
\end{enumerate}\end{promptbox}
\textbf{Workspace}\workspace
\needspace{8\baselineskip}
\subsection*{\pid{APP-068-S-01} — Standard}
\begin{problembox}
A small dataset for agronomy, forestry, biomass and ecological field data has points $(0,7),(1,9),(2,11),(3,14)$. Find the least-squares line $y=mx+b$.
\end{problembox}
\begin{promptbox}\begin{enumerate}
\item What are the observations and response variable?
\item What model class is being assumed?
\item What statistic or fitted parameter is requested?
\item How should residuals or uncertainty be interpreted?
\end{enumerate}\end{promptbox}
\textbf{Workspace}\workspace
\needspace{8\baselineskip}
\subsection*{\pid{APP-068-A-01} — Advanced}
\begin{problembox}
For a simple uncertainty exercise related to agronomy, forestry, biomass and ecological field data, ten uniform samples on $[0,1]$ are [0.11, 0.27, 0.33, 0.48, 0.52, 0.61, 0.73, 0.81, 0.89, 0.95]. Estimate $P(X^2<0.5)$ by Monte Carlo frequency and compare with the exact probability.
\end{problembox}
\begin{promptbox}\begin{enumerate}
\item What are the observations and response variable?
\item What model class is being assumed?
\item What statistic or fitted parameter is requested?
\item How should residuals or uncertainty be interpreted?
\end{enumerate}\end{promptbox}
\textbf{Workspace}\workspace
\begin{reflectionbox}\begin{enumerate}
\item Which sentence or datum in these problems supplied the main mathematical constraint?
\item Which assumptions were essential to the model, and how would the model change if one were removed?
\item When would an analytical solution become inconvenient enough to justify a numerical method?
\end{enumerate}\end{reflectionbox}
\chapter{Probability simulation and uncertainty propagation}\label{app:A069}
\textbf{Application ID:} \pid{A069}\quad \textbf{Source locus:} 5–14\quad \textbf{Source pages:} 351, 581, 584, 606, 607, 667, 678, 989, 997
\section*{Application Context}
This family begins with measured or simulated data and uses statistics, regression, probability, or Monte Carlo techniques to quantify patterns and uncertainty. The canonical application family is \emph{Probability simulation and uncertainty propagation}.
\section*{Calculus Toolkit}
\begin{tabularx}{\textwidth}{@{}lX@{}}\toprule
Primary technique & Random variates and Monte Carlo \\
Prerequisites & functions, descriptive statistics, model fitting \\
Numerical extension & Ch. 14 (statistics), Ch. 15 (modeling), Ch. 7 (random numbers) \\
Analytical/numerical class & NUMERICAL-NATURAL \\
\bottomrule\end{tabularx}
\section*{Modeling Questions to Ask}\begin{itemize}
\item What are the observations and response variable?
\item What model class is being assumed?
\item What statistic or fitted parameter is requested?
\item How should residuals or uncertainty be interpreted?
\end{itemize}
\begin{errorbox}Do not treat a fitted relationship as an exact law. Distinguish observations, model predictions, residuals, and uncertainty.\end{errorbox}
\section*{Problems}
\needspace{8\baselineskip}
\subsection*{\pid{APP-069-F-01} — Foundation}
\begin{problembox}
Five observations in probability simulation and uncertainty propagation are [14, 16, 19, 17, 24]. Compute the sample mean and sample variance.
\end{problembox}
\begin{promptbox}\begin{enumerate}
\item What are the observations and response variable?
\item What model class is being assumed?
\item What statistic or fitted parameter is requested?
\item How should residuals or uncertainty be interpreted?
\end{enumerate}\end{promptbox}
\textbf{Workspace}\workspace
\needspace{8\baselineskip}
\subsection*{\pid{APP-069-S-01} — Standard}
\begin{problembox}
A small dataset for probability simulation and uncertainty propagation has points $(0,5),(1,7),(2,9),(3,12)$. Find the least-squares line $y=mx+b$.
\end{problembox}
\begin{promptbox}\begin{enumerate}
\item What are the observations and response variable?
\item What model class is being assumed?
\item What statistic or fitted parameter is requested?
\item How should residuals or uncertainty be interpreted?
\end{enumerate}\end{promptbox}
\textbf{Workspace}\workspace
\needspace{8\baselineskip}
\subsection*{\pid{APP-069-A-01} — Advanced}
\begin{problembox}
For a simple uncertainty exercise related to probability simulation and uncertainty propagation, ten uniform samples on $[0,1]$ are [0.11, 0.27, 0.33, 0.48, 0.52, 0.61, 0.73, 0.81, 0.89, 0.95]. Estimate $P(X^2<0.5)$ by Monte Carlo frequency and compare with the exact probability.
\end{problembox}
\begin{promptbox}\begin{enumerate}
\item What are the observations and response variable?
\item What model class is being assumed?
\item What statistic or fitted parameter is requested?
\item How should residuals or uncertainty be interpreted?
\end{enumerate}\end{promptbox}
\textbf{Workspace}\workspace
\begin{reflectionbox}\begin{enumerate}
\item Which sentence or datum in these problems supplied the main mathematical constraint?
\item Which assumptions were essential to the model, and how would the model change if one were removed?
\item When would an analytical solution become inconvenient enough to justify a numerical method?
\end{enumerate}\end{reflectionbox}
\part{Spectral and Mode Applications}
\chapter{Vibrating string, wave motion and periodic mechanical signals}\label{app:A070}
\textbf{Application ID:} \pid{A070}\quad \textbf{Source locus:} 2; 7\quad \textbf{Source pages:} 142, 161, 531
\section*{Application Context}
This family represents periodic or modal behavior using sinusoidal components, Fourier coefficients, eigenvectors, or related transforms. The canonical application family is \emph{Vibrating string, wave motion and periodic mechanical signals}.
\section*{Calculus Toolkit}
\begin{tabularx}{\textwidth}{@{}lX@{}}\toprule
Primary technique & Fourier/spectral decomposition \\
Prerequisites & trigonometric functions, complex/vector algebra as needed \\
Numerical extension & Ch. 12 (FFT), Ch. 13 (spectral applications) \\
Analytical/numerical class & NUMERICAL-NATURAL \\
\bottomrule\end{tabularx}
\section*{Modeling Questions to Ask}\begin{itemize}
\item What is the period or sampling interval?
\item Which frequencies can be represented?
\item What does a spectral coefficient mean physically?
\item How is the transform normalized?
\end{itemize}
\begin{errorbox}Frequency, angular frequency, period, and sampling index are different quantities. Keep their units and normalizations explicit.\end{errorbox}
\section*{Problems}
\needspace{8\baselineskip}
\subsection*{\pid{APP-070-F-01} — Foundation}
\begin{problembox}
A periodic signal associated with vibrating string, wave motion and periodic mechanical signals is $s(t)=4\cos(1.2566t)$. Find its period and ordinary frequency.
\end{problembox}
\begin{promptbox}\begin{enumerate}
\item What is the period or sampling interval?
\item Which frequencies can be represented?
\item What does a spectral coefficient mean physically?
\item How is the transform normalized?
\end{enumerate}\end{promptbox}
\textbf{Workspace}\workspace
\needspace{8\baselineskip}
\subsection*{\pid{APP-070-S-01} — Standard}
\begin{problembox}
A four-sample signal in vibrating string, wave motion and periodic mechanical signals is $x=[1,0,-1,0]$. Compute its unnormalized discrete Fourier transform $X_k=\sum_{n=0}^3 x_ne^{-2\pi i kn/4}$.
\end{problembox}
\begin{promptbox}\begin{enumerate}
\item What is the period or sampling interval?
\item Which frequencies can be represented?
\item What does a spectral coefficient mean physically?
\item How is the transform normalized?
\end{enumerate}\end{promptbox}
\textbf{Workspace}\workspace
\needspace{8\baselineskip}
\subsection*{\pid{APP-070-A-01} — Advanced}
\begin{problembox}
A two-mode model for vibrating string, wave motion and periodic mechanical signals has system matrix $A=\begin{pmatrix}2&1\\1&2\end{pmatrix}$. Find its eigenvalues and eigenvectors and interpret them as independent modal directions.
\end{problembox}
\begin{promptbox}\begin{enumerate}
\item What is the period or sampling interval?
\item Which frequencies can be represented?
\item What does a spectral coefficient mean physically?
\item How is the transform normalized?
\end{enumerate}\end{promptbox}
\textbf{Workspace}\workspace
\begin{reflectionbox}\begin{enumerate}
\item Which sentence or datum in these problems supplied the main mathematical constraint?
\item Which assumptions were essential to the model, and how would the model change if one were removed?
\item When would an analytical solution become inconvenient enough to justify a numerical method?
\end{enumerate}\end{reflectionbox}
\chapter{Periodic temperature and respiratory-cycle signals}\label{app:A071}
\textbf{Application ID:} \pid{A071}\quad \textbf{Source locus:} 2; 4; 5\quad \textbf{Source pages:} 142, 293, 413
\section*{Application Context}
This family represents periodic or modal behavior using sinusoidal components, Fourier coefficients, eigenvectors, or related transforms. The canonical application family is \emph{Periodic temperature and respiratory-cycle signals}.
\section*{Calculus Toolkit}
\begin{tabularx}{\textwidth}{@{}lX@{}}\toprule
Primary technique & Periodic function analysis \\
Prerequisites & trigonometric functions, complex/vector algebra as needed \\
Numerical extension & Ch. 12 (FFT), Ch. 13 (spectral applications) \\
Analytical/numerical class & NUMERICAL-NATURAL \\
\bottomrule\end{tabularx}
\section*{Modeling Questions to Ask}\begin{itemize}
\item What is the period or sampling interval?
\item Which frequencies can be represented?
\item What does a spectral coefficient mean physically?
\item How is the transform normalized?
\end{itemize}
\begin{errorbox}Frequency, angular frequency, period, and sampling index are different quantities. Keep their units and normalizations explicit.\end{errorbox}
\section*{Problems}
\needspace{8\baselineskip}
\subsection*{\pid{APP-071-F-01} — Foundation}
\begin{problembox}
A periodic signal associated with periodic temperature and respiratory-cycle signals is $s(t)=5\cos(1.0472t)$. Find its period and ordinary frequency.
\end{problembox}
\begin{promptbox}\begin{enumerate}
\item What is the period or sampling interval?
\item Which frequencies can be represented?
\item What does a spectral coefficient mean physically?
\item How is the transform normalized?
\end{enumerate}\end{promptbox}
\textbf{Workspace}\workspace
\needspace{8\baselineskip}
\subsection*{\pid{APP-071-S-01} — Standard}
\begin{problembox}
A four-sample signal in periodic temperature and respiratory-cycle signals is $x=[1,0,-1,0]$. Compute its unnormalized discrete Fourier transform $X_k=\sum_{n=0}^3 x_ne^{-2\pi i kn/4}$.
\end{problembox}
\begin{promptbox}\begin{enumerate}
\item What is the period or sampling interval?
\item Which frequencies can be represented?
\item What does a spectral coefficient mean physically?
\item How is the transform normalized?
\end{enumerate}\end{promptbox}
\textbf{Workspace}\workspace
\needspace{8\baselineskip}
\subsection*{\pid{APP-071-A-01} — Advanced}
\begin{problembox}
A two-mode model for periodic temperature and respiratory-cycle signals has system matrix $A=\begin{pmatrix}2&1\\1&2\end{pmatrix}$. Find its eigenvalues and eigenvectors and interpret them as independent modal directions.
\end{problembox}
\begin{promptbox}\begin{enumerate}
\item What is the period or sampling interval?
\item Which frequencies can be represented?
\item What does a spectral coefficient mean physically?
\item How is the transform normalized?
\end{enumerate}\end{promptbox}
\textbf{Workspace}\workspace
\begin{reflectionbox}\begin{enumerate}
\item Which sentence or datum in these problems supplied the main mathematical constraint?
\item Which assumptions were essential to the model, and how would the model change if one were removed?
\item When would an analytical solution become inconvenient enough to justify a numerical method?
\end{enumerate}\end{reflectionbox}
\chapter{Structural normal modes and vibration-oriented beam/cable models}\label{app:A072}
\textbf{Application ID:} \pid{A072}\quad \textbf{Source locus:} 5.9; 10\quad \textbf{Source pages:} 393, 697
\section*{Application Context}
This family represents periodic or modal behavior using sinusoidal components, Fourier coefficients, eigenvectors, or related transforms. The canonical application family is \emph{Structural normal modes and vibration-oriented beam/cable models}.
\section*{Calculus Toolkit}
\begin{tabularx}{\textwidth}{@{}lX@{}}\toprule
Primary technique & Eigenvalue/spectral analysis \\
Prerequisites & trigonometric functions, complex/vector algebra as needed \\
Numerical extension & Ch. 12 (FFT), Ch. 13 (spectral applications) \\
Analytical/numerical class & NUMERICAL-NATURAL \\
\bottomrule\end{tabularx}
\section*{Modeling Questions to Ask}\begin{itemize}
\item What is the period or sampling interval?
\item Which frequencies can be represented?
\item What does a spectral coefficient mean physically?
\item How is the transform normalized?
\end{itemize}
\begin{errorbox}Frequency, angular frequency, period, and sampling index are different quantities. Keep their units and normalizations explicit.\end{errorbox}
\section*{Problems}
\needspace{8\baselineskip}
\subsection*{\pid{APP-072-F-01} — Foundation}
\begin{problembox}
A periodic signal associated with structural normal modes and vibration-oriented beam/cable models is $s(t)=3\cos(1.5708t)$. Find its period and ordinary frequency.
\end{problembox}
\begin{promptbox}\begin{enumerate}
\item What is the period or sampling interval?
\item Which frequencies can be represented?
\item What does a spectral coefficient mean physically?
\item How is the transform normalized?
\end{enumerate}\end{promptbox}
\textbf{Workspace}\workspace
\needspace{8\baselineskip}
\subsection*{\pid{APP-072-S-01} — Standard}
\begin{problembox}
A four-sample signal in structural normal modes and vibration-oriented beam/cable models is $x=[1,0,-1,0]$. Compute its unnormalized discrete Fourier transform $X_k=\sum_{n=0}^3 x_ne^{-2\pi i kn/4}$.
\end{problembox}
\begin{promptbox}\begin{enumerate}
\item What is the period or sampling interval?
\item Which frequencies can be represented?
\item What does a spectral coefficient mean physically?
\item How is the transform normalized?
\end{enumerate}\end{promptbox}
\textbf{Workspace}\workspace
\needspace{8\baselineskip}
\subsection*{\pid{APP-072-A-01} — Advanced}
\begin{problembox}
A two-mode model for structural normal modes and vibration-oriented beam/cable models has system matrix $A=\begin{pmatrix}2&1\\1&2\end{pmatrix}$. Find its eigenvalues and eigenvectors and interpret them as independent modal directions.
\end{problembox}
\begin{promptbox}\begin{enumerate}
\item What is the period or sampling interval?
\item Which frequencies can be represented?
\item What does a spectral coefficient mean physically?
\item How is the transform normalized?
\end{enumerate}\end{promptbox}
\textbf{Workspace}\workspace
\begin{reflectionbox}\begin{enumerate}
\item Which sentence or datum in these problems supplied the main mathematical constraint?
\item Which assumptions were essential to the model, and how would the model change if one were removed?
\item When would an analytical solution become inconvenient enough to justify a numerical method?
\end{enumerate}\end{reflectionbox}
\chapter{Multiscale or localized signal/image analysis}\label{app:A073}
\textbf{Application ID:} \pid{A073}\quad \textbf{Source locus:} 9; 13\quad \textbf{Source pages:} general extension
\section*{Application Context}
This family represents periodic or modal behavior using sinusoidal components, Fourier coefficients, eigenvectors, or related transforms. The canonical application family is \emph{Multiscale or localized signal/image analysis}.
\section*{Calculus Toolkit}
\begin{tabularx}{\textwidth}{@{}lX@{}}\toprule
Primary technique & Wavelet transforms \\
Prerequisites & trigonometric functions, complex/vector algebra as needed \\
Numerical extension & Ch. 12 (FFT), Ch. 13 (spectral applications) \\
Analytical/numerical class & NUMERICAL-NATURAL \\
\bottomrule\end{tabularx}
\section*{Modeling Questions to Ask}\begin{itemize}
\item What is the period or sampling interval?
\item Which frequencies can be represented?
\item What does a spectral coefficient mean physically?
\item How is the transform normalized?
\end{itemize}
\begin{errorbox}Frequency, angular frequency, period, and sampling index are different quantities. Keep their units and normalizations explicit.\end{errorbox}
\section*{Problems}
\needspace{8\baselineskip}
\subsection*{\pid{APP-073-F-01} — Foundation}
\begin{problembox}
A periodic signal associated with multiscale or localized signal/image analysis is $s(t)=4\cos(1.2566t)$. Find its period and ordinary frequency.
\end{problembox}
\begin{promptbox}\begin{enumerate}
\item What is the period or sampling interval?
\item Which frequencies can be represented?
\item What does a spectral coefficient mean physically?
\item How is the transform normalized?
\end{enumerate}\end{promptbox}
\textbf{Workspace}\workspace
\needspace{8\baselineskip}
\subsection*{\pid{APP-073-S-01} — Standard}
\begin{problembox}
A four-sample signal in multiscale or localized signal/image analysis is $x=[1,0,-1,0]$. Compute its unnormalized discrete Fourier transform $X_k=\sum_{n=0}^3 x_ne^{-2\pi i kn/4}$.
\end{problembox}
\begin{promptbox}\begin{enumerate}
\item What is the period or sampling interval?
\item Which frequencies can be represented?
\item What does a spectral coefficient mean physically?
\item How is the transform normalized?
\end{enumerate}\end{promptbox}
\textbf{Workspace}\workspace
\needspace{8\baselineskip}
\subsection*{\pid{APP-073-A-01} — Advanced}
\begin{problembox}
A two-mode model for multiscale or localized signal/image analysis has system matrix $A=\begin{pmatrix}2&1\\1&2\end{pmatrix}$. Find its eigenvalues and eigenvectors and interpret them as independent modal directions.
\end{problembox}
\begin{promptbox}\begin{enumerate}
\item What is the period or sampling interval?
\item Which frequencies can be represented?
\item What does a spectral coefficient mean physically?
\item How is the transform normalized?
\end{enumerate}\end{promptbox}
\textbf{Workspace}\workspace
\begin{reflectionbox}\begin{enumerate}
\item Which sentence or datum in these problems supplied the main mathematical constraint?
\item Which assumptions were essential to the model, and how would the model change if one were removed?
\item When would an analytical solution become inconvenient enough to justify a numerical method?
\end{enumerate}\end{reflectionbox}
\part{Advanced Field, Boundary, and PDE Applications}
\chapter{Heat equation and temperature distribution}\label{app:A074}
\textbf{Application ID:} \pid{A074}\quad \textbf{Source locus:} 13; advanced DE\quad \textbf{Source pages:} 901, 902, 925, 930, 967
\section*{Application Context}
This family extends calculus to boundary-value, field, inverse, or PDE problems where discretization and linear algebra frequently become part of the solution. The canonical application family is \emph{Heat equation and temperature distribution}.
\section*{Calculus Toolkit}
\begin{tabularx}{\textwidth}{@{}lX@{}}\toprule
Primary technique & Parabolic PDEs and boundary/initial conditions \\
Prerequisites & multivariable calculus, differential equations, linear algebra \\
Numerical extension & Chs. 18–21 (boundary values, inverse problems, PDEs, computational geometry) \\
Analytical/numerical class & NUMERICAL-REQUIRED \\
\bottomrule\end{tabularx}
\section*{Modeling Questions to Ask}\begin{itemize}
\item What PDE/BVP/field equation is being modeled?
\item What boundary and initial conditions are required?
\item What quantity is conserved or balanced?
\item Which part is analytical and which part requires discretization?
\end{itemize}
\begin{errorbox}Separate the mathematical model, boundary/initial conditions, and discretization. A numerically solvable scheme is not automatically stable or convergent.\end{errorbox}
\section*{Problems}
\needspace{8\baselineskip}
\subsection*{\pid{APP-074-F-01} — Foundation}
\begin{problembox}
For an advanced model related to heat equation and temperature distribution, verify that $u(x,t)=e^{-\pi^2t}\sin(\pi x)$ satisfies the heat equation $u_t=u_{xx}$ and the boundary conditions $u(0,t)=u(1,t)=0$.
\end{problembox}
\begin{promptbox}\begin{enumerate}
\item What PDE/BVP/field equation is being modeled?
\item What boundary and initial conditions are required?
\item What quantity is conserved or balanced?
\item Which part is analytical and which part requires discretization?
\end{enumerate}\end{promptbox}
\textbf{Workspace}\workspace
\needspace{8\baselineskip}
\subsection*{\pid{APP-074-S-01} — Standard}
\begin{problembox}
A finite-difference model for heat equation and temperature distribution has one interior grid value $u$ surrounded by boundary values 14, 22, 32, and 40. For the five-point Laplace stencil, compute the interior value.
\end{problembox}
\begin{promptbox}\begin{enumerate}
\item What PDE/BVP/field equation is being modeled?
\item What boundary and initial conditions are required?
\item What quantity is conserved or balanced?
\item Which part is analytical and which part requires discretization?
\end{enumerate}\end{promptbox}
\textbf{Workspace}\workspace
\needspace{8\baselineskip}
\subsection*{\pid{APP-074-A-01} — Advanced}
\begin{problembox}
As a computational extension of heat equation and temperature distribution, solve $-u''=1$ on $0<x<1$ with $u(0)=u(1)=0$ using a centered finite difference on grid points $x=0,1/3,2/3,1$. Find the two interior unknowns and compare with the exact solution $u=x(1-x)/2$.
\end{problembox}
\begin{promptbox}\begin{enumerate}
\item What PDE/BVP/field equation is being modeled?
\item What boundary and initial conditions are required?
\item What quantity is conserved or balanced?
\item Which part is analytical and which part requires discretization?
\end{enumerate}\end{promptbox}
\textbf{Workspace}\workspace
\begin{reflectionbox}\begin{enumerate}
\item Which sentence or datum in these problems supplied the main mathematical constraint?
\item Which assumptions were essential to the model, and how would the model change if one were removed?
\item When would an analytical solution become inconvenient enough to justify a numerical method?
\end{enumerate}\end{reflectionbox}
\chapter{Wave equation and wave propagation}\label{app:A075}
\textbf{Application ID:} \pid{A075}\quad \textbf{Source locus:} 13; advanced DE\quad \textbf{Source pages:} 901, 968
\section*{Application Context}
This family extends calculus to boundary-value, field, inverse, or PDE problems where discretization and linear algebra frequently become part of the solution. The canonical application family is \emph{Wave equation and wave propagation}.
\section*{Calculus Toolkit}
\begin{tabularx}{\textwidth}{@{}lX@{}}\toprule
Primary technique & Hyperbolic PDEs \\
Prerequisites & multivariable calculus, differential equations, linear algebra \\
Numerical extension & Chs. 18–21 (boundary values, inverse problems, PDEs, computational geometry) \\
Analytical/numerical class & NUMERICAL-REQUIRED \\
\bottomrule\end{tabularx}
\section*{Modeling Questions to Ask}\begin{itemize}
\item What PDE/BVP/field equation is being modeled?
\item What boundary and initial conditions are required?
\item What quantity is conserved or balanced?
\item Which part is analytical and which part requires discretization?
\end{itemize}
\begin{errorbox}Separate the mathematical model, boundary/initial conditions, and discretization. A numerically solvable scheme is not automatically stable or convergent.\end{errorbox}
\section*{Problems}
\needspace{8\baselineskip}
\subsection*{\pid{APP-075-F-01} — Foundation}
\begin{problembox}
For an advanced model related to wave equation and wave propagation, verify that $u(x,t)=e^{-\pi^2t}\sin(\pi x)$ satisfies the heat equation $u_t=u_{xx}$ and the boundary conditions $u(0,t)=u(1,t)=0$.
\end{problembox}
\begin{promptbox}\begin{enumerate}
\item What PDE/BVP/field equation is being modeled?
\item What boundary and initial conditions are required?
\item What quantity is conserved or balanced?
\item Which part is analytical and which part requires discretization?
\end{enumerate}\end{promptbox}
\textbf{Workspace}\workspace
\needspace{8\baselineskip}
\subsection*{\pid{APP-075-S-01} — Standard}
\begin{problembox}
A finite-difference model for wave equation and wave propagation has one interior grid value $u$ surrounded by boundary values 10, 23, 30, and 41. For the five-point Laplace stencil, compute the interior value.
\end{problembox}
\begin{promptbox}\begin{enumerate}
\item What PDE/BVP/field equation is being modeled?
\item What boundary and initial conditions are required?
\item What quantity is conserved or balanced?
\item Which part is analytical and which part requires discretization?
\end{enumerate}\end{promptbox}
\textbf{Workspace}\workspace
\needspace{8\baselineskip}
\subsection*{\pid{APP-075-A-01} — Advanced}
\begin{problembox}
As a computational extension of wave equation and wave propagation, solve $-u''=1$ on $0<x<1$ with $u(0)=u(1)=0$ using a centered finite difference on grid points $x=0,1/3,2/3,1$. Find the two interior unknowns and compare with the exact solution $u=x(1-x)/2$.
\end{problembox}
\begin{promptbox}\begin{enumerate}
\item What PDE/BVP/field equation is being modeled?
\item What boundary and initial conditions are required?
\item What quantity is conserved or balanced?
\item Which part is analytical and which part requires discretization?
\end{enumerate}\end{promptbox}
\textbf{Workspace}\workspace
\begin{reflectionbox}\begin{enumerate}
\item Which sentence or datum in these problems supplied the main mathematical constraint?
\item Which assumptions were essential to the model, and how would the model change if one were removed?
\item When would an analytical solution become inconvenient enough to justify a numerical method?
\end{enumerate}\end{reflectionbox}
\chapter{Motion of a liquid, fluid fields and continuum flow}\label{app:A076}
\textbf{Application ID:} \pid{A076}\quad \textbf{Source locus:} 15\quad \textbf{Source pages:} 1122, 1123, 1126
\section*{Application Context}
This family extends calculus to boundary-value, field, inverse, or PDE problems where discretization and linear algebra frequently become part of the solution. The canonical application family is \emph{Motion of a liquid, fluid fields and continuum flow}.
\section*{Calculus Toolkit}
\begin{tabularx}{\textwidth}{@{}lX@{}}\toprule
Primary technique & Vector calculus leading to PDEs \\
Prerequisites & multivariable calculus, differential equations, linear algebra \\
Numerical extension & Chs. 18–21 (boundary values, inverse problems, PDEs, computational geometry) \\
Analytical/numerical class & NUMERICAL-REQUIRED \\
\bottomrule\end{tabularx}
\section*{Modeling Questions to Ask}\begin{itemize}
\item What PDE/BVP/field equation is being modeled?
\item What boundary and initial conditions are required?
\item What quantity is conserved or balanced?
\item Which part is analytical and which part requires discretization?
\end{itemize}
\begin{errorbox}Separate the mathematical model, boundary/initial conditions, and discretization. A numerically solvable scheme is not automatically stable or convergent.\end{errorbox}
\section*{Problems}
\needspace{8\baselineskip}
\subsection*{\pid{APP-076-F-01} — Foundation}
\begin{problembox}
For an advanced model related to motion of a liquid, fluid fields and continuum flow, verify that $u(x,t)=e^{-\pi^2t}\sin(\pi x)$ satisfies the heat equation $u_t=u_{xx}$ and the boundary conditions $u(0,t)=u(1,t)=0$.
\end{problembox}
\begin{promptbox}\begin{enumerate}
\item What PDE/BVP/field equation is being modeled?
\item What boundary and initial conditions are required?
\item What quantity is conserved or balanced?
\item Which part is analytical and which part requires discretization?
\end{enumerate}\end{promptbox}
\textbf{Workspace}\workspace
\needspace{8\baselineskip}
\subsection*{\pid{APP-076-S-01} — Standard}
\begin{problembox}
A finite-difference model for motion of a liquid, fluid fields and continuum flow has one interior grid value $u$ surrounded by boundary values 11, 20, 31, and 40. For the five-point Laplace stencil, compute the interior value.
\end{problembox}
\begin{promptbox}\begin{enumerate}
\item What PDE/BVP/field equation is being modeled?
\item What boundary and initial conditions are required?
\item What quantity is conserved or balanced?
\item Which part is analytical and which part requires discretization?
\end{enumerate}\end{promptbox}
\textbf{Workspace}\workspace
\needspace{8\baselineskip}
\subsection*{\pid{APP-076-A-01} — Advanced}
\begin{problembox}
As a computational extension of motion of a liquid, fluid fields and continuum flow, solve $-u''=1$ on $0<x<1$ with $u(0)=u(1)=0$ using a centered finite difference on grid points $x=0,1/3,2/3,1$. Find the two interior unknowns and compare with the exact solution $u=x(1-x)/2$.
\end{problembox}
\begin{promptbox}\begin{enumerate}
\item What PDE/BVP/field equation is being modeled?
\item What boundary and initial conditions are required?
\item What quantity is conserved or balanced?
\item Which part is analytical and which part requires discretization?
\end{enumerate}\end{promptbox}
\textbf{Workspace}\workspace
\begin{reflectionbox}\begin{enumerate}
\item Which sentence or datum in these problems supplied the main mathematical constraint?
\item Which assumptions were essential to the model, and how would the model change if one were removed?
\item When would an analytical solution become inconvenient enough to justify a numerical method?
\end{enumerate}\end{reflectionbox}
\chapter{Electric/gravitational fields, Gauss's law, charge and heat flux}\label{app:A077}
\textbf{Application ID:} \pid{A077}\quad \textbf{Source locus:} 15\quad \textbf{Source pages:} 1045, 1107, 1109, 1127
\section*{Application Context}
This family extends calculus to boundary-value, field, inverse, or PDE problems where discretization and linear algebra frequently become part of the solution. The canonical application family is \emph{Electric/gravitational fields, Gauss's law, charge and heat flux}.
\section*{Calculus Toolkit}
\begin{tabularx}{\textwidth}{@{}lX@{}}\toprule
Primary technique & Divergence, flux and field equations \\
Prerequisites & multivariable calculus, differential equations, linear algebra \\
Numerical extension & Chs. 18–21 (boundary values, inverse problems, PDEs, computational geometry) \\
Analytical/numerical class & NUMERICAL-REQUIRED \\
\bottomrule\end{tabularx}
\section*{Modeling Questions to Ask}\begin{itemize}
\item What PDE/BVP/field equation is being modeled?
\item What boundary and initial conditions are required?
\item What quantity is conserved or balanced?
\item Which part is analytical and which part requires discretization?
\end{itemize}
\begin{errorbox}Separate the mathematical model, boundary/initial conditions, and discretization. A numerically solvable scheme is not automatically stable or convergent.\end{errorbox}
\section*{Problems}
\needspace{8\baselineskip}
\subsection*{\pid{APP-077-F-01} — Foundation}
\begin{problembox}
For an advanced model related to electric/gravitational fields, gauss's law, charge and heat flux, verify that $u(x,t)=e^{-\pi^2t}\sin(\pi x)$ satisfies the heat equation $u_t=u_{xx}$ and the boundary conditions $u(0,t)=u(1,t)=0$.
\end{problembox}
\begin{promptbox}\begin{enumerate}
\item What PDE/BVP/field equation is being modeled?
\item What boundary and initial conditions are required?
\item What quantity is conserved or balanced?
\item Which part is analytical and which part requires discretization?
\end{enumerate}\end{promptbox}
\textbf{Workspace}\workspace
\needspace{8\baselineskip}
\subsection*{\pid{APP-077-S-01} — Standard}
\begin{problembox}
A finite-difference model for electric/gravitational fields, gauss's law, charge and heat flux has one interior grid value $u$ surrounded by boundary values 12, 21, 32, and 41. For the five-point Laplace stencil, compute the interior value.
\end{problembox}
\begin{promptbox}\begin{enumerate}
\item What PDE/BVP/field equation is being modeled?
\item What boundary and initial conditions are required?
\item What quantity is conserved or balanced?
\item Which part is analytical and which part requires discretization?
\end{enumerate}\end{promptbox}
\textbf{Workspace}\workspace
\needspace{8\baselineskip}
\subsection*{\pid{APP-077-A-01} — Advanced}
\begin{problembox}
As a computational extension of electric/gravitational fields, gauss's law, charge and heat flux, solve $-u''=1$ on $0<x<1$ with $u(0)=u(1)=0$ using a centered finite difference on grid points $x=0,1/3,2/3,1$. Find the two interior unknowns and compare with the exact solution $u=x(1-x)/2$.
\end{problembox}
\begin{promptbox}\begin{enumerate}
\item What PDE/BVP/field equation is being modeled?
\item What boundary and initial conditions are required?
\item What quantity is conserved or balanced?
\item Which part is analytical and which part requires discretization?
\end{enumerate}\end{promptbox}
\textbf{Workspace}\workspace
\begin{reflectionbox}\begin{enumerate}
\item Which sentence or datum in these problems supplied the main mathematical constraint?
\item Which assumptions were essential to the model, and how would the model change if one were removed?
\item When would an analytical solution become inconvenient enough to justify a numerical method?
\end{enumerate}\end{reflectionbox}
\chapter{Beam deflection, suspension bridges and boundary-constrained structural curves}\label{app:A078}
\textbf{Application ID:} \pid{A078}\quad \textbf{Source locus:} 7; 10\quad \textbf{Source pages:} 484, 697
\section*{Application Context}
This family extends calculus to boundary-value, field, inverse, or PDE problems where discretization and linear algebra frequently become part of the solution. The canonical application family is \emph{Beam deflection, suspension bridges and boundary-constrained structural curves}.
\section*{Calculus Toolkit}
\begin{tabularx}{\textwidth}{@{}lX@{}}\toprule
Primary technique & Boundary-value ODEs \\
Prerequisites & multivariable calculus, differential equations, linear algebra \\
Numerical extension & Chs. 18–21 (boundary values, inverse problems, PDEs, computational geometry) \\
Analytical/numerical class & NUMERICAL-NATURAL \\
\bottomrule\end{tabularx}
\section*{Modeling Questions to Ask}\begin{itemize}
\item What PDE/BVP/field equation is being modeled?
\item What boundary and initial conditions are required?
\item What quantity is conserved or balanced?
\item Which part is analytical and which part requires discretization?
\end{itemize}
\begin{errorbox}Separate the mathematical model, boundary/initial conditions, and discretization. A numerically solvable scheme is not automatically stable or convergent.\end{errorbox}
\section*{Problems}
\needspace{8\baselineskip}
\subsection*{\pid{APP-078-F-01} — Foundation}
\begin{problembox}
For an advanced model related to beam deflection, suspension bridges and boundary-constrained structural curves, verify that $u(x,t)=e^{-\pi^2t}\sin(\pi x)$ satisfies the heat equation $u_t=u_{xx}$ and the boundary conditions $u(0,t)=u(1,t)=0$.
\end{problembox}
\begin{promptbox}\begin{enumerate}
\item What PDE/BVP/field equation is being modeled?
\item What boundary and initial conditions are required?
\item What quantity is conserved or balanced?
\item Which part is analytical and which part requires discretization?
\end{enumerate}\end{promptbox}
\textbf{Workspace}\workspace
\needspace{8\baselineskip}
\subsection*{\pid{APP-078-S-01} — Standard}
\begin{problembox}
A finite-difference model for beam deflection, suspension bridges and boundary-constrained structural curves has one interior grid value $u$ surrounded by boundary values 13, 22, 30, and 40. For the five-point Laplace stencil, compute the interior value.
\end{problembox}
\begin{promptbox}\begin{enumerate}
\item What PDE/BVP/field equation is being modeled?
\item What boundary and initial conditions are required?
\item What quantity is conserved or balanced?
\item Which part is analytical and which part requires discretization?
\end{enumerate}\end{promptbox}
\textbf{Workspace}\workspace
\needspace{8\baselineskip}
\subsection*{\pid{APP-078-A-01} — Advanced}
\begin{problembox}
As a computational extension of beam deflection, suspension bridges and boundary-constrained structural curves, solve $-u''=1$ on $0<x<1$ with $u(0)=u(1)=0$ using a centered finite difference on grid points $x=0,1/3,2/3,1$. Find the two interior unknowns and compare with the exact solution $u=x(1-x)/2$.
\end{problembox}
\begin{promptbox}\begin{enumerate}
\item What PDE/BVP/field equation is being modeled?
\item What boundary and initial conditions are required?
\item What quantity is conserved or balanced?
\item Which part is analytical and which part requires discretization?
\end{enumerate}\end{promptbox}
\textbf{Workspace}\workspace
\begin{reflectionbox}\begin{enumerate}
\item Which sentence or datum in these problems supplied the main mathematical constraint?
\item Which assumptions were essential to the model, and how would the model change if one were removed?
\item When would an analytical solution become inconvenient enough to justify a numerical method?
\end{enumerate}\end{reflectionbox}
\chapter{Electromagnetic theory and field simulation}\label{app:A079}
\textbf{Application ID:} \pid{A079}\quad \textbf{Source locus:} 9 / advanced continuation\quad \textbf{Source pages:} 581
\section*{Application Context}
This family extends calculus to boundary-value, field, inverse, or PDE problems where discretization and linear algebra frequently become part of the solution. The canonical application family is \emph{Electromagnetic theory and field simulation}.
\section*{Calculus Toolkit}
\begin{tabularx}{\textwidth}{@{}lX@{}}\toprule
Primary technique & Field equations and spectral/PDE methods \\
Prerequisites & multivariable calculus, differential equations, linear algebra \\
Numerical extension & Chs. 18–21 (boundary values, inverse problems, PDEs, computational geometry) \\
Analytical/numerical class & NUMERICAL-REQUIRED \\
\bottomrule\end{tabularx}
\section*{Modeling Questions to Ask}\begin{itemize}
\item What PDE/BVP/field equation is being modeled?
\item What boundary and initial conditions are required?
\item What quantity is conserved or balanced?
\item Which part is analytical and which part requires discretization?
\end{itemize}
\begin{errorbox}Separate the mathematical model, boundary/initial conditions, and discretization. A numerically solvable scheme is not automatically stable or convergent.\end{errorbox}
\section*{Problems}
\needspace{8\baselineskip}
\subsection*{\pid{APP-079-F-01} — Foundation}
\begin{problembox}
For an advanced model related to electromagnetic theory and field simulation, verify that $u(x,t)=e^{-\pi^2t}\sin(\pi x)$ satisfies the heat equation $u_t=u_{xx}$ and the boundary conditions $u(0,t)=u(1,t)=0$.
\end{problembox}
\begin{promptbox}\begin{enumerate}
\item What PDE/BVP/field equation is being modeled?
\item What boundary and initial conditions are required?
\item What quantity is conserved or balanced?
\item Which part is analytical and which part requires discretization?
\end{enumerate}\end{promptbox}
\textbf{Workspace}\workspace
\needspace{8\baselineskip}
\subsection*{\pid{APP-079-S-01} — Standard}
\begin{problembox}
A finite-difference model for electromagnetic theory and field simulation has one interior grid value $u$ surrounded by boundary values 14, 23, 31, and 41. For the five-point Laplace stencil, compute the interior value.
\end{problembox}
\begin{promptbox}\begin{enumerate}
\item What PDE/BVP/field equation is being modeled?
\item What boundary and initial conditions are required?
\item What quantity is conserved or balanced?
\item Which part is analytical and which part requires discretization?
\end{enumerate}\end{promptbox}
\textbf{Workspace}\workspace
\needspace{8\baselineskip}
\subsection*{\pid{APP-079-A-01} — Advanced}
\begin{problembox}
As a computational extension of electromagnetic theory and field simulation, solve $-u''=1$ on $0<x<1$ with $u(0)=u(1)=0$ using a centered finite difference on grid points $x=0,1/3,2/3,1$. Find the two interior unknowns and compare with the exact solution $u=x(1-x)/2$.
\end{problembox}
\begin{promptbox}\begin{enumerate}
\item What PDE/BVP/field equation is being modeled?
\item What boundary and initial conditions are required?
\item What quantity is conserved or balanced?
\item Which part is analytical and which part requires discretization?
\end{enumerate}\end{promptbox}
\textbf{Workspace}\workspace
\begin{reflectionbox}\begin{enumerate}
\item Which sentence or datum in these problems supplied the main mathematical constraint?
\item Which assumptions were essential to the model, and how would the model change if one were removed?
\item When would an analytical solution become inconvenient enough to justify a numerical method?
\end{enumerate}\end{reflectionbox}
\chapter{Hyperbolic detection, geometric navigation, satellite antenna and reflector geometry}\label{app:A080}
\textbf{Application ID:} \pid{A080}\quad \textbf{Source locus:} 10–11\quad \textbf{Source pages:} 688, 695, 699, 746
\section*{Application Context}
This family extends calculus to boundary-value, field, inverse, or PDE problems where discretization and linear algebra frequently become part of the solution. The canonical application family is \emph{Hyperbolic detection, geometric navigation, satellite antenna and reflector geometry}.
\section*{Calculus Toolkit}
\begin{tabularx}{\textwidth}{@{}lX@{}}\toprule
Primary technique & Conics, intersections and computational geometry \\
Prerequisites & multivariable calculus, differential equations, linear algebra \\
Numerical extension & Chs. 18–21 (boundary values, inverse problems, PDEs, computational geometry) \\
Analytical/numerical class & NUMERICAL-NATURAL \\
\bottomrule\end{tabularx}
\section*{Modeling Questions to Ask}\begin{itemize}
\item What PDE/BVP/field equation is being modeled?
\item What boundary and initial conditions are required?
\item What quantity is conserved or balanced?
\item Which part is analytical and which part requires discretization?
\end{itemize}
\begin{errorbox}Separate the mathematical model, boundary/initial conditions, and discretization. A numerically solvable scheme is not automatically stable or convergent.\end{errorbox}
\section*{Problems}
\needspace{8\baselineskip}
\subsection*{\pid{APP-080-F-01} — Foundation}
\begin{problembox}
For an advanced model related to hyperbolic detection, geometric navigation, satellite antenna and reflector geometry, verify that $u(x,t)=e^{-\pi^2t}\sin(\pi x)$ satisfies the heat equation $u_t=u_{xx}$ and the boundary conditions $u(0,t)=u(1,t)=0$.
\end{problembox}
\begin{promptbox}\begin{enumerate}
\item What PDE/BVP/field equation is being modeled?
\item What boundary and initial conditions are required?
\item What quantity is conserved or balanced?
\item Which part is analytical and which part requires discretization?
\end{enumerate}\end{promptbox}
\textbf{Workspace}\workspace
\needspace{8\baselineskip}
\subsection*{\pid{APP-080-S-01} — Standard}
\begin{problembox}
A finite-difference model for hyperbolic detection, geometric navigation, satellite antenna and reflector geometry has one interior grid value $u$ surrounded by boundary values 10, 20, 32, and 40. For the five-point Laplace stencil, compute the interior value.
\end{problembox}
\begin{promptbox}\begin{enumerate}
\item What PDE/BVP/field equation is being modeled?
\item What boundary and initial conditions are required?
\item What quantity is conserved or balanced?
\item Which part is analytical and which part requires discretization?
\end{enumerate}\end{promptbox}
\textbf{Workspace}\workspace
\needspace{8\baselineskip}
\subsection*{\pid{APP-080-A-01} — Advanced}
\begin{problembox}
As a computational extension of hyperbolic detection, geometric navigation, satellite antenna and reflector geometry, solve $-u''=1$ on $0<x<1$ with $u(0)=u(1)=0$ using a centered finite difference on grid points $x=0,1/3,2/3,1$. Find the two interior unknowns and compare with the exact solution $u=x(1-x)/2$.
\end{problembox}
\begin{promptbox}\begin{enumerate}
\item What PDE/BVP/field equation is being modeled?
\item What boundary and initial conditions are required?
\item What quantity is conserved or balanced?
\item Which part is analytical and which part requires discretization?
\end{enumerate}\end{promptbox}
\textbf{Workspace}\workspace
\begin{reflectionbox}\begin{enumerate}
\item Which sentence or datum in these problems supplied the main mathematical constraint?
\item Which assumptions were essential to the model, and how would the model change if one were removed?
\item When would an analytical solution become inconvenient enough to justify a numerical method?
\end{enumerate}\end{reflectionbox}
\chapter{Architecture, bridge/building design and general engineering geometry}\label{app:A081}
\textbf{Application ID:} \pid{A081}\quad \textbf{Source locus:} 7; 10; 14–15\quad \textbf{Source pages:} 453, 571, 698, 1012, 1039, 1068
\section*{Application Context}
This family extends calculus to boundary-value, field, inverse, or PDE problems where discretization and linear algebra frequently become part of the solution. The canonical application family is \emph{Architecture, bridge/building design and general engineering geometry}.
\section*{Calculus Toolkit}
\begin{tabularx}{\textwidth}{@{}lX@{}}\toprule
Primary technique & Geometry, optimization and multiple integration \\
Prerequisites & multivariable calculus, differential equations, linear algebra \\
Numerical extension & Chs. 18–21 (boundary values, inverse problems, PDEs, computational geometry) \\
Analytical/numerical class & NUMERICAL-NATURAL \\
\bottomrule\end{tabularx}
\section*{Modeling Questions to Ask}\begin{itemize}
\item What PDE/BVP/field equation is being modeled?
\item What boundary and initial conditions are required?
\item What quantity is conserved or balanced?
\item Which part is analytical and which part requires discretization?
\end{itemize}
\begin{errorbox}Separate the mathematical model, boundary/initial conditions, and discretization. A numerically solvable scheme is not automatically stable or convergent.\end{errorbox}
\section*{Problems}
\needspace{8\baselineskip}
\subsection*{\pid{APP-081-F-01} — Foundation}
\begin{problembox}
For an advanced model related to architecture, bridge/building design and general engineering geometry, verify that $u(x,t)=e^{-\pi^2t}\sin(\pi x)$ satisfies the heat equation $u_t=u_{xx}$ and the boundary conditions $u(0,t)=u(1,t)=0$.
\end{problembox}
\begin{promptbox}\begin{enumerate}
\item What PDE/BVP/field equation is being modeled?
\item What boundary and initial conditions are required?
\item What quantity is conserved or balanced?
\item Which part is analytical and which part requires discretization?
\end{enumerate}\end{promptbox}
\textbf{Workspace}\workspace
\needspace{8\baselineskip}
\subsection*{\pid{APP-081-S-01} — Standard}
\begin{problembox}
A finite-difference model for architecture, bridge/building design and general engineering geometry has one interior grid value $u$ surrounded by boundary values 11, 21, 30, and 41. For the five-point Laplace stencil, compute the interior value.
\end{problembox}
\begin{promptbox}\begin{enumerate}
\item What PDE/BVP/field equation is being modeled?
\item What boundary and initial conditions are required?
\item What quantity is conserved or balanced?
\item Which part is analytical and which part requires discretization?
\end{enumerate}\end{promptbox}
\textbf{Workspace}\workspace
\needspace{8\baselineskip}
\subsection*{\pid{APP-081-A-01} — Advanced}
\begin{problembox}
As a computational extension of architecture, bridge/building design and general engineering geometry, solve $-u''=1$ on $0<x<1$ with $u(0)=u(1)=0$ using a centered finite difference on grid points $x=0,1/3,2/3,1$. Find the two interior unknowns and compare with the exact solution $u=x(1-x)/2$.
\end{problembox}
\begin{promptbox}\begin{enumerate}
\item What PDE/BVP/field equation is being modeled?
\item What boundary and initial conditions are required?
\item What quantity is conserved or balanced?
\item Which part is analytical and which part requires discretization?
\end{enumerate}\end{promptbox}
\textbf{Workspace}\workspace
\begin{reflectionbox}\begin{enumerate}
\item Which sentence or datum in these problems supplied the main mathematical constraint?
\item Which assumptions were essential to the model, and how would the model change if one were removed?
\item When would an analytical solution become inconvenient enough to justify a numerical method?
\end{enumerate}\end{reflectionbox}
\chapter{Sphereflake, surface geometry and spherical integration}\label{app:A082}
\textbf{Application ID:} \pid{A082}\quad \textbf{Source locus:} 9; 14\quad \textbf{Source pages:} 607, 1030
\section*{Application Context}
This family extends calculus to boundary-value, field, inverse, or PDE problems where discretization and linear algebra frequently become part of the solution. The canonical application family is \emph{Sphereflake, surface geometry and spherical integration}.
\section*{Calculus Toolkit}
\begin{tabularx}{\textwidth}{@{}lX@{}}\toprule
Primary technique & Recursive geometry and spherical surfaces \\
Prerequisites & multivariable calculus, differential equations, linear algebra \\
Numerical extension & Chs. 18–21 (boundary values, inverse problems, PDEs, computational geometry) \\
Analytical/numerical class & NUMERICAL-NATURAL \\
\bottomrule\end{tabularx}
\section*{Modeling Questions to Ask}\begin{itemize}
\item What PDE/BVP/field equation is being modeled?
\item What boundary and initial conditions are required?
\item What quantity is conserved or balanced?
\item Which part is analytical and which part requires discretization?
\end{itemize}
\begin{errorbox}Separate the mathematical model, boundary/initial conditions, and discretization. A numerically solvable scheme is not automatically stable or convergent.\end{errorbox}
\section*{Problems}
\needspace{8\baselineskip}
\subsection*{\pid{APP-082-F-01} — Foundation}
\begin{problembox}
For an advanced model related to sphereflake, surface geometry and spherical integration, verify that $u(x,t)=e^{-\pi^2t}\sin(\pi x)$ satisfies the heat equation $u_t=u_{xx}$ and the boundary conditions $u(0,t)=u(1,t)=0$.
\end{problembox}
\begin{promptbox}\begin{enumerate}
\item What PDE/BVP/field equation is being modeled?
\item What boundary and initial conditions are required?
\item What quantity is conserved or balanced?
\item Which part is analytical and which part requires discretization?
\end{enumerate}\end{promptbox}
\textbf{Workspace}\workspace
\needspace{8\baselineskip}
\subsection*{\pid{APP-082-S-01} — Standard}
\begin{problembox}
A finite-difference model for sphereflake, surface geometry and spherical integration has one interior grid value $u$ surrounded by boundary values 12, 22, 31, and 40. For the five-point Laplace stencil, compute the interior value.
\end{problembox}
\begin{promptbox}\begin{enumerate}
\item What PDE/BVP/field equation is being modeled?
\item What boundary and initial conditions are required?
\item What quantity is conserved or balanced?
\item Which part is analytical and which part requires discretization?
\end{enumerate}\end{promptbox}
\textbf{Workspace}\workspace
\needspace{8\baselineskip}
\subsection*{\pid{APP-082-A-01} — Advanced}
\begin{problembox}
As a computational extension of sphereflake, surface geometry and spherical integration, solve $-u''=1$ on $0<x<1$ with $u(0)=u(1)=0$ using a centered finite difference on grid points $x=0,1/3,2/3,1$. Find the two interior unknowns and compare with the exact solution $u=x(1-x)/2$.
\end{problembox}
\begin{promptbox}\begin{enumerate}
\item What PDE/BVP/field equation is being modeled?
\item What boundary and initial conditions are required?
\item What quantity is conserved or balanced?
\item Which part is analytical and which part requires discretization?
\end{enumerate}\end{promptbox}
\textbf{Workspace}\workspace
\begin{reflectionbox}\begin{enumerate}
\item Which sentence or datum in these problems supplied the main mathematical constraint?
\item Which assumptions were essential to the model, and how would the model change if one were removed?
\item When would an analytical solution become inconvenient enough to justify a numerical method?
\end{enumerate}\end{reflectionbox}
\chapter{Inverse/parameter-recovery problems from measured fields or responses}\label{app:A083}
\textbf{Application ID:} \pid{A083}\quad \textbf{Source locus:} 13–15\quad \textbf{Source pages:} general computational extension
\section*{Application Context}
This family extends calculus to boundary-value, field, inverse, or PDE problems where discretization and linear algebra frequently become part of the solution. The canonical application family is \emph{Inverse/parameter-recovery problems from measured fields or responses}.
\section*{Calculus Toolkit}
\begin{tabularx}{\textwidth}{@{}lX@{}}\toprule
Primary technique & Inverse problems and regularization \\
Prerequisites & multivariable calculus, differential equations, linear algebra \\
Numerical extension & Chs. 18–21 (boundary values, inverse problems, PDEs, computational geometry) \\
Analytical/numerical class & NUMERICAL-REQUIRED \\
\bottomrule\end{tabularx}
\section*{Modeling Questions to Ask}\begin{itemize}
\item What PDE/BVP/field equation is being modeled?
\item What boundary and initial conditions are required?
\item What quantity is conserved or balanced?
\item Which part is analytical and which part requires discretization?
\end{itemize}
\begin{errorbox}Separate the mathematical model, boundary/initial conditions, and discretization. A numerically solvable scheme is not automatically stable or convergent.\end{errorbox}
\section*{Problems}
\needspace{8\baselineskip}
\subsection*{\pid{APP-083-F-01} — Foundation}
\begin{problembox}
For an advanced model related to inverse/parameter-recovery problems from measured fields or responses, verify that $u(x,t)=e^{-\pi^2t}\sin(\pi x)$ satisfies the heat equation $u_t=u_{xx}$ and the boundary conditions $u(0,t)=u(1,t)=0$.
\end{problembox}
\begin{promptbox}\begin{enumerate}
\item What PDE/BVP/field equation is being modeled?
\item What boundary and initial conditions are required?
\item What quantity is conserved or balanced?
\item Which part is analytical and which part requires discretization?
\end{enumerate}\end{promptbox}
\textbf{Workspace}\workspace
\needspace{8\baselineskip}
\subsection*{\pid{APP-083-S-01} — Standard}
\begin{problembox}
A finite-difference model for inverse/parameter-recovery problems from measured fields or responses has one interior grid value $u$ surrounded by boundary values 13, 23, 32, and 41. For the five-point Laplace stencil, compute the interior value.
\end{problembox}
\begin{promptbox}\begin{enumerate}
\item What PDE/BVP/field equation is being modeled?
\item What boundary and initial conditions are required?
\item What quantity is conserved or balanced?
\item Which part is analytical and which part requires discretization?
\end{enumerate}\end{promptbox}
\textbf{Workspace}\workspace
\needspace{8\baselineskip}
\subsection*{\pid{APP-083-A-01} — Advanced}
\begin{problembox}
As a computational extension of inverse/parameter-recovery problems from measured fields or responses, solve $-u''=1$ on $0<x<1$ with $u(0)=u(1)=0$ using a centered finite difference on grid points $x=0,1/3,2/3,1$. Find the two interior unknowns and compare with the exact solution $u=x(1-x)/2$.
\end{problembox}
\begin{promptbox}\begin{enumerate}
\item What PDE/BVP/field equation is being modeled?
\item What boundary and initial conditions are required?
\item What quantity is conserved or balanced?
\item Which part is analytical and which part requires discretization?
\end{enumerate}\end{promptbox}
\textbf{Workspace}\workspace
\begin{reflectionbox}\begin{enumerate}
\item Which sentence or datum in these problems supplied the main mathematical constraint?
\item Which assumptions were essential to the model, and how would the model change if one were removed?
\item When would an analytical solution become inconvenient enough to justify a numerical method?
\end{enumerate}\end{reflectionbox}
\part{Mixed Application Review}
\chapter{Mixed Review — Technique Not Given}
The following problems intentionally omit the technique label. Decide how to model each problem before calculating.
\section*{Mixed Review 1: \pid{MR-01}}
\begin{problembox}For a short planning horizon, a quantity associated with population, cell-phone subscribers, gdp, apartment rental and other baseline trend models is modeled by $Q(t)=53e^{0.045t}$, where $t$ is measured in years. Find the instantaneous growth rate at $t=5$ and interpret it.\end{problembox}
\workspace
\section*{Mixed Review 2: \pid{MR-02}}
\begin{problembox}A rigid segment of length 10 m is used in a angular speed and angular rate of change model. One endpoint moves horizontally away from a reference point at 1.00 m/s while the other endpoint moves vertically. When the horizontal distance is 6 m, how fast is the vertical distance changing?\end{problembox}
\workspace
\section*{Mixed Review 3: \pid{MR-03}}
\begin{problembox}A calibration equation for break-even point and cost–revenue intersections is $e^{-x}=x/2.0$. Use bisection on $[0,2]$ until the interval width is below $10^{-4}$, and report the solution.\end{problembox}
\workspace
\section*{Mixed Review 4: \pid{MR-04}}
\begin{problembox}A cost model for beam strength, illumination, engine efficiency and mechanical design is $C(x)=3100/x+84+0.6x$ for $x>0$. Find the operating value $x$ that minimizes cost and the minimum cost.\end{problembox}
\workspace
\section*{Mixed Review 5: \pid{MR-05}}
\begin{problembox}In a arc length, cable length, gateway arch, pursuit and stream length application, liquid fills a vertical prismatic tank of height 6 m and horizontal cross-sectional area 4 m$^2$. Use a weight density of 9.8 kN/m$^3$. Find the work required to pump all liquid just to the top edge.\end{problembox}
\workspace
\section*{Mixed Review 6: \pid{MR-06}}
\begin{problembox}In a probability integrals and expectation-like calculations application, liquid fills a vertical prismatic tank of height 5 m and horizontal cross-sectional area 2 m$^2$. Use a weight density of 9.8 kN/m$^3$. Find the work required to pump all liquid just to the top edge.\end{problembox}
\workspace
\section*{Mixed Review 7: \pid{MR-07}}
\begin{problembox}A repeated contribution in special-function physics: pendulum refinements, elliptic geometry and spherical harmonics has magnitudes $315, 204.8, 133.1,\ldots$, with common ratio $r=0.65$. Determine the infinite accumulated total and state the convergence condition used.\end{problembox}
\workspace
\section*{Mixed Review 8: \pid{MR-08}}
\begin{problembox}A scalar field used in mass, moments and density distributions in two/three dimensions is $F(x,y,z)=x^2+2y^2+3z$. At $P=(3,2,1)$, find the gradient and the directional derivative toward the direction $(1,2,1)$.\end{problembox}
\workspace
\section*{Mixed Review 9: \pid{MR-09}}
\begin{problembox}A bounded-growth model for torricelli flow, emptying tanks and hydraulic transients satisfies $y'=0.30y(1-y/540)$ with $y(0)=45$. Use the logistic solution to estimate $y(8)$.\end{problembox}
\workspace
\section*{Mixed Review 10: \pid{MR-10}}
\begin{problembox}A small dataset for sales, inflation, health-care expenditures, government expenditures and gdp data has points $(0,6),(1,8),(2,10),(3,13)$. Find the least-squares line $y=mx+b$.\end{problembox}
\workspace
\section*{Mixed Review 11: \pid{MR-11}}
\begin{problembox}A four-sample signal in periodic temperature and respiratory-cycle signals is $x=[1,0,-1,0]$. Compute its unnormalized discrete Fourier transform $X_k=\sum_{n=0}^3 x_ne^{-2\pi i kn/4}$.\end{problembox}
\workspace
\section*{Mixed Review 12: \pid{MR-12}}
\begin{problembox}A finite-difference model for beam deflection, suspension bridges and boundary-constrained structural curves has one interior grid value $u$ surrounded by boundary values 13, 22, 30, and 40. For the five-point Laplace stencil, compute the interior value.\end{problembox}
\workspace
\part{Capstone Modeling Problems}
\chapter{Interdisciplinary Capstones}
\section*{\pid{CAP-01} — Solar-storage design}
\begin{problembox}A cylindrical thermal-storage tank must hold 120 m$^3$. The lateral material costs twice as much per square meter as the top and bottom material. Model total relative material cost, determine the radius and height minimizing it, and then estimate how a 1\% radius fabrication error changes the optimized volume using differentials.\end{problembox}
\workspace\workspace
\section*{\pid{CAP-02} — Tracking and interception}
\begin{problembox}A sensor reports a target position $\mathbf r(t)=(20+4t,\,10+2t,\,100-5t^2)$ meters. Determine velocity and acceleration, find the time when altitude reaches 40 m, and compute the horizontal distance traveled by then.\end{problembox}
\workspace\workspace
\section*{\pid{CAP-03} — Heat-and-data model}
\begin{problembox}Temperature in a small test chamber follows $T'= -0.18(T-22)$ with $T(0)=90$. Find the exact model, determine when $T=35$, and use two Euler steps of size 1 minute to approximate the first two minutes for comparison.\end{problembox}
\workspace\workspace
\section*{\pid{CAP-04} — Field-flux calculation}
\begin{problembox}A vector field is $\mathbf F=(x,y,z)$ and $S$ is the sphere of radius 2. Use the divergence theorem to compute outward flux and interpret the result as total source strength inside the sphere.\end{problembox}
\workspace\workspace
\part{Reference Material}
\chapter{Application Crosswalk}
\small\begin{longtable}{@{}p{0.065\textwidth}p{0.30\textwidth}p{0.16\textwidth}p{0.22\textwidth}p{0.14\textwidth}@{}}\toprule
ID & Application & Source locus & Primary technique & Problem IDs \\\midrule\endfirsthead\toprule ID & Application & Source locus & Primary technique & Problem IDs \\\midrule\endhead
A001 & Population, cell-phone subscribers, GDP, apartment rental and other baseline trend models & P.1–P.3 & Functions, graphs, linear models and rates & APP-001-F-01, APP-001-S-01, APP-001-A-01  \\
A002 & Carbon-dioxide concentration and environmental trend models & P.1 & Function modeling from data & APP-002-F-01, APP-002-S-01, APP-002-A-01  \\
A003 & Biodiesel production, sales, salary and simple production models & P.1–P.3; 9 & Function modeling; sequences & APP-003-F-01, APP-003-S-01, APP-003-A-01  \\
A004 & Ferris-wheel height, angular position and periodic geometry & P.4; 12 & Trigonometric functions and parametric motion & APP-004-F-01, APP-004-S-01, APP-004-A-01  \\
A005 & Automobile aerodynamics, thermostat and ripple models & P.3; 2.6 & Functional dependence and geometry & APP-005-F-01, APP-005-S-01, APP-005-A-01  \\
A006 & Average velocity, speed and acceleration & 1; 2.1–2.3; 12.3 & Difference quotients, derivatives and higher derivatives & APP-006-F-01, APP-006-S-01, APP-006-A-01  \\
A007 & Particle, rectilinear and vertical motion & 2; 4; 12 & Position, velocity and acceleration & APP-007-F-01, APP-007-S-01, APP-007-A-01  \\
A008 & Angular speed and angular rate of change & P.4; 2.6; 5.8 & Chain rule and related rates & APP-008-F-01, APP-008-S-01, APP-008-A-01  \\
A009 & Moving ladder, moving shadow, air-traffic and angle-of-elevation problems & 2.6; 10; 12 & Implicit differentiation and related rates & APP-009-F-01, APP-009-S-01, APP-009-A-01  \\
A010 & Flow rate, oil leak and changing-volume systems & 4.4–4.5; 5; 15 & Derivative/integral rate relations & APP-010-F-01, APP-010-S-01, APP-010-A-01  \\
A011 & Marginal cost, revenue, productivity and utility & 13.3; 13.8–13.10 & Partial derivatives and marginal analysis & APP-011-F-01, APP-011-S-01, APP-011-A-01  \\
A012 & Sales growth, declining sales, depreciation and investment growth & 2; 5; 6; 9 & Rates of change and exponential models & APP-012-F-01, APP-012-S-01, APP-012-A-01  \\
A013 & Concentration, trachea contraction and physiological rate-of-change models & 3.9; 6 & Differentials, sensitivity and rate models & APP-013-F-01, APP-013-S-01, APP-013-A-01  \\
A014 & Error and differential approximation & 3.9; 13.4 & Differentials and propagated error & APP-014-F-01, APP-014-S-01, APP-014-A-01  \\
A015 & Break-even point and cost–revenue intersections & P.2; 7.1 & Zeros/intersections of economic functions & APP-015-F-01, APP-015-S-01, APP-015-A-01  \\
A016 & Boiling temperature, atmospheric-pressure/altitude inversions and calibration problems & 5; 6; 13 & Solve nonlinear scalar equations & APP-016-F-01, APP-016-S-01, APP-016-A-01  \\
A017 & Optical/refraction and geometric intersection problems & 3; 10; 13 & Nonlinear equations from geometry & APP-017-F-01, APP-017-S-01, APP-017-A-01  \\
A018 & Polynomial design/intersection and machine-part equations & 3; 7 & Polynomial zeros & APP-018-F-01, APP-018-S-01, APP-018-A-01  \\
A019 & Maximum area, minimum area, perimeter, length and surface area & 3.7; 13.8–13.10 & Scalar and constrained optimization & APP-019-F-01, APP-019-S-01, APP-019-A-01  \\
A020 & Maximum volume and packaging/design optimization & 3.7; 13.8–13.10 & Extrema in one and several variables & APP-020-F-01, APP-020-S-01, APP-020-A-01  \\
A021 & Minimum distance, time and travel-path optimization & 3.7; 10; 13 & Geometric optimization & APP-021-F-01, APP-021-S-01, APP-021-A-01  \\
A022 & Beam strength, illumination, engine efficiency and mechanical design & 3.7; 13 & Engineering objective functions & APP-022-F-01, APP-022-S-01, APP-022-A-01  \\
A023 & Highway design, honeycomb and structural geometry optimization & 3; 12 & Extrema and curvature constraints & APP-023-F-01, APP-023-S-01, APP-023-A-01  \\
A024 & Inventory management/replenishment and manufacturing optimization & 1; 2; 7 & Cost/objective optimization & APP-024-F-01, APP-024-S-01, APP-024-A-01  \\
A025 & Maximum profit/revenue and minimum cost & 3.7; 13.9–13.10 & Business optimization and constraints & APP-025-F-01, APP-025-S-01, APP-025-A-01  \\
A026 & Building a pipeline & 13.9 & Multivariable geometric optimization & APP-026-F-01, APP-026-S-01, APP-026-A-01  \\
A027 & Area and area between curves & 4; 7.1 & Definite integration and accumulation & APP-027-F-01, APP-027-S-01, APP-027-A-01  \\
A028 & Volumes by disks/shells and general solids & 7.2–7.3; 14 & One- and multiple-integral volume & APP-028-F-01, APP-028-S-01, APP-028-A-01  \\
A029 & Arc length, cable length, Gateway Arch, pursuit and stream length & 7.4 & Arc-length integrals & APP-029-F-01, APP-029-S-01, APP-029-A-01  \\
A030 & Surface area and surfaces of revolution & 7.4; 14.5 & Surface-area integrals & APP-030-F-01, APP-030-S-01, APP-030-A-01  \\
A031 & Work: lifting, pumping, towing, springs, gases, rockets and force fields & 7.5; 15 & Work integrals and line integrals & APP-031-F-01, APP-031-S-01, APP-031-A-01  \\
A032 & Fluid pressure and fluid force & 7.7 & Pressure-force integration & APP-032-F-01, APP-032-S-01, APP-032-A-01  \\
A033 & Center of mass, centroid, moments and inertia & 7.6; 14.4 & Moment integrals & APP-033-F-01, APP-033-S-01, APP-033-A-01  \\
A034 & Consumer and producer surplus & 7.7 / applied integration & Area between economic curves & APP-034-F-01, APP-034-S-01, APP-034-A-01  \\
A035 & Average temperature, production and profit & 14 & Average-value/multiple integration & APP-035-F-01, APP-035-S-01, APP-035-A-01  \\
A036 & Probability integrals and expectation-like calculations & 4–9; 14 & Integration of probability models & APP-036-F-01, APP-036-S-01, APP-036-A-01  \\
A037 & Topographic maps, terrain and scattered spatial data & 2; 13 & Interpolation, gradients and surfaces & APP-037-F-01, APP-037-S-01, APP-037-A-01  \\
A038 & Mercator map and coordinate-transform applications & 5.9 & Transcendental functions and integration & APP-038-F-01, APP-038-S-01, APP-038-A-01  \\
A039 & Tractrix and implicit/transcendental curve models & 5; 8; 10 & Implicit curves, special functions and quadrature & APP-039-F-01, APP-039-S-01, APP-039-A-01  \\
A040 & Sound intensity, decibel-like and exponential/logarithmic models & 1; 5; 6 & Logarithmic/exponential functions & APP-040-F-01, APP-040-S-01, APP-040-A-01  \\
A041 & Compound interest, annuities, home mortgage, present value and capitalized cost & 5; 6; 9 & Exponential models, series and accumulation & APP-041-F-01, APP-041-S-01, APP-041-A-01  \\
A042 & Cantor set, disappearing table, bouncing-ball distance and geometric series & 9 & Sequences, series and convergence & APP-042-F-01, APP-042-S-01, APP-042-A-01  \\
A043 & Special-function physics: pendulum refinements, elliptic geometry and spherical harmonics & 10; 12; advanced & Special functions and nonlinear models & APP-043-F-01, APP-043-S-01, APP-043-A-01  \\
A044 & Resultant force/velocity and vector mechanics & 11.1–11.2 & Vector addition and components & APP-044-F-01, APP-044-S-01, APP-044-A-01  \\
A045 & Torque, work from vector forces and rigid-mechanics geometry & 11.3–11.5 & Dot/cross products and vector geometry & APP-045-F-01, APP-045-S-01, APP-045-A-01  \\
A046 & Cable tension, load supports, navigation and surveying & 10–11 & Linear/vector systems & APP-046-F-01, APP-046-S-01, APP-046-A-01  \\
A047 & Ideal-gas, meteorology and apparent-temperature multivariable models & 13.2–13.6 & Partial derivatives and gradients & APP-047-F-01, APP-047-S-01, APP-047-A-01  \\
A048 & Cobb–Douglas production, marginal analysis and economic surfaces & 13 & Partial derivatives, gradients and constrained extrema & APP-048-F-01, APP-048-S-01, APP-048-A-01  \\
A049 & Hardy–Weinberg law, Shannon diversity and constrained population composition & 13.9–13.10 & Constrained multivariable optimization & APP-049-F-01, APP-049-S-01, APP-049-A-01  \\
A050 & Mass, moments and density distributions in two/three dimensions & 14; 15 & Multiple integrals and vector fields & APP-050-F-01, APP-050-S-01, APP-050-A-01  \\
A051 & Field strength, electric potential and gravitational potential & 7; 13; 15 & Scalar/vector fields and gradients & APP-051-F-01, APP-051-S-01, APP-051-A-01  \\
A052 & Radioactive decay and carbon dating & 5.5; 6.2–6.3 & Exponential decay ODE & APP-052-F-01, APP-052-S-01, APP-052-A-01  \\
A053 & Bacterial culture, tree/dog growth and population growth & 5–6 & Growth/logistic ODEs & APP-053-F-01, APP-053-S-01, APP-053-A-01  \\
A054 & Endangered species, epidemic and ecological population models & 6; 8 & Coupled/separable population ODEs & APP-054-F-01, APP-054-S-01, APP-054-A-01  \\
A055 & Chemical mixture, tracer-drug and intravenous-feeding models & 6.4 & First-order linear compartment ODEs & APP-055-F-01, APP-055-S-01, APP-055-A-01  \\
A056 & Newton cooling, heat transfer and thermal transients & 6.2 & First-order thermal ODEs & APP-056-F-01, APP-056-S-01, APP-056-A-01  \\
A057 & Torricelli flow, emptying tanks and hydraulic transients & 6; 7 & Separable ODEs and fluid models & APP-057-F-01, APP-057-S-01, APP-057-A-01  \\
A058 & Electric circuits, resistance/inductance dynamics and control & 5–6; 13 & First-order and coupled circuit ODEs & APP-058-F-01, APP-058-S-01, APP-058-A-01  \\
A059 & Spring motion, pendulum and harmonic motion & 2; 6; 12 & Second-order ODEs; oscillation & APP-059-F-01, APP-059-S-01, APP-059-A-01  \\
A060 & Falling bodies and velocity in a resisting medium & 1; 6 & Initial-value ODEs with drag & APP-060-F-01, APP-060-S-01, APP-060-A-01  \\
A061 & Projectile trajectories, baseball, football, and shot paths; tossing bales & 3; 10; 12–13 & Parametric/vector ODE trajectory models & APP-061-F-01, APP-061-S-01, APP-061-A-01  \\
A062 & Orbital motion, satellites, planetary motion and Kepler-law simulation & 10–12 & Second-order gravitational ODEs & APP-062-F-01, APP-062-S-01, APP-062-A-01  \\
A063 & Heat-seeking particle/path and pursuit dynamics & 13.6–13.7 & Gradient-driven ODE path & APP-063-F-01, APP-063-S-01, APP-063-A-01  \\
A064 & Sales, inflation, health-care expenditures, government expenditures and GDP data & P; 5; 13 & Regression, trend estimation and uncertainty & APP-064-F-01, APP-064-S-01, APP-064-A-01  \\
A065 & Systolic blood pressure, intelligence-test and physiological data & 13 & Statistical description and fitted models & APP-065-F-01, APP-065-S-01, APP-065-A-01  \\
A066 & Learning curve and memory model & 6–7 & Nonlinear model fitting & APP-066-F-01, APP-066-S-01, APP-066-A-01  \\
A067 & Average typing speed and experimental-rate datasets & 3 & Statistics and fitted rates & APP-067-F-01, APP-067-S-01, APP-067-A-01  \\
A068 & Agronomy, forestry, biomass and ecological field data & 6; 13 & Regression, interpolation, and statistics & APP-068-F-01, APP-068-S-01, APP-068-A-01  \\
A069 & Probability simulation and uncertainty propagation & 5–14 & Random variates and Monte Carlo & APP-069-F-01, APP-069-S-01, APP-069-A-01  \\
A070 & Vibrating string, wave motion and periodic mechanical signals & 2; 7 & Fourier/spectral decomposition & APP-070-F-01, APP-070-S-01, APP-070-A-01  \\
A071 & Periodic temperature and respiratory-cycle signals & 2; 4; 5 & Periodic function analysis & APP-071-F-01, APP-071-S-01, APP-071-A-01  \\
A072 & Structural normal modes and vibration-oriented beam/cable models & 5.9; 10 & Eigenvalue/spectral analysis & APP-072-F-01, APP-072-S-01, APP-072-A-01  \\
A073 & Multiscale or localized signal/image analysis & 9; 13 & Wavelet transforms & APP-073-F-01, APP-073-S-01, APP-073-A-01  \\
A074 & Heat equation and temperature distribution & 13; advanced DE & Parabolic PDEs and boundary/initial conditions & APP-074-F-01, APP-074-S-01, APP-074-A-01  \\
A075 & Wave equation and wave propagation & 13; advanced DE & Hyperbolic PDEs & APP-075-F-01, APP-075-S-01, APP-075-A-01  \\
A076 & Motion of a liquid, fluid fields and continuum flow & 15 & Vector calculus leading to PDEs & APP-076-F-01, APP-076-S-01, APP-076-A-01  \\
A077 & Electric/gravitational fields, Gauss's law, charge and heat flux & 15 & Divergence, flux and field equations & APP-077-F-01, APP-077-S-01, APP-077-A-01  \\
A078 & Beam deflection, suspension bridges and boundary-constrained structural curves & 7; 10 & Boundary-value ODEs & APP-078-F-01, APP-078-S-01, APP-078-A-01  \\
A079 & Electromagnetic theory and field simulation & 9 / advanced continuation & Field equations and spectral/PDE methods & APP-079-F-01, APP-079-S-01, APP-079-A-01  \\
A080 & Hyperbolic detection, geometric navigation, satellite antenna and reflector geometry & 10–11 & Conics, intersections and computational geometry & APP-080-F-01, APP-080-S-01, APP-080-A-01  \\
A081 & Architecture, bridge/building design and general engineering geometry & 7; 10; 14–15 & Geometry, optimization and multiple integration & APP-081-F-01, APP-081-S-01, APP-081-A-01  \\
A082 & Sphereflake, surface geometry and spherical integration & 9; 14 & Recursive geometry and spherical surfaces & APP-082-F-01, APP-082-S-01, APP-082-A-01  \\
A083 & Inverse/parameter-recovery problems from measured fields or responses & 13–15 & Inverse problems and regularization & APP-083-F-01, APP-083-S-01, APP-083-A-01  \\
\bottomrule\end{longtable}\normalsize
\chapter{Sources and Scope Notes}
\begin{itemize}
\item The application inventory and calculus-topic crosswalk are derived from the supplied calculus source and its application index; problems in this workbook are newly authored.
\item The modeling pedagogy is an original synthesis of Eugene Don and Benay Don, \emph{How to Solve Word Problems in Calculus: A Solved Problem Approach} (McGraw-Hill, 2001), especially its emphasis on constructing a mathematical model before solving it and on learning general procedures rather than memorizing isolated templates.
\item Numerical extensions use the method taxonomy of Press, Teukolsky, Vetterling, and Flannery, \emph{Numerical Recipes: The Art of Scientific Computing}, Third Edition (Cambridge University Press, 2007).
\item Proof-oriented explanations use the writing and proof-strategy principles of Chartrand, Polimeni, and Zhang, \emph{Mathematical Proofs: A Transition to Advanced Mathematics}, Fourth Edition, only where justification improves understanding.
\end{itemize}
\end{document}
